\chapter{Anatomy}

\medterm{Abdominal Aorta}-- From diaphragmatic hiatus at T12 to common iliac bifurcation at L4. Major branches: celiac trunk (T12), SMA (L1), IMA (L3), renal arteries (L1-L2), gonadal arteries, lumbar arteries. Anterior to the vertebral column. Abdominal aortic aneurysm is common in elderly males (13th leading cause of death). Rupture causes hypotension, back pain, and pulsatile mass with over 80 percent mortality.

\medterm{Abdominal Cavity}-- The largest body cavity, bounded superiorly by the diaphragm, inferiorly by the pelvic inlet (where it is continuous with the pelvic cavity, collectively termed the abdominopelvic cavity), anterolaterally by the abdominal muscles, and posteriorly by the lumbar vertebral column, psoas, and quadratus lumborum muscles. It is lined by the peritoneum, a continuous serous membrane dividing the space into the peritoneal cavity (containing greater and lesser sacs connected via the epiploic foramen of Winslow) and the retroperitoneal space. It houses major digestive organs (stomach, small and large intestines, liver, gallbladder, pancreas), the spleen, and retroperitoneal structures (kidneys, adrenal glands, abdominal aorta, inferior vena cava). Retroperitoneal organs follow the SAD PUCKER mnemonic: Suprarenal glands, Aorta/IVC, Duodenum (2nd--4th parts), Pancreas (head/body), Ureters, Colon (ascending/descending), Kidneys, Esophagus, Rectum. Pathologies include peritonitis, hemoperitoneum, ascites, intra-abdominal abscesses, and trauma-induced organ laceration.

\medterm{Abducens Nerve (Cranial Nerve VI)}-- The sixth cranial nerve (CN VI), a purely somatic motor nerve that innervates the lateral rectus muscle of the eye, responsible for abduction of the globe. Its nucleus is situated in the dorsal caudal pons in the floor of the fourth ventricle, encircled by fibers of the facial nerve (facial colliculus). It emerges from the brainstem at the pontomedullary junction, ascends the clivus, traverses Dorello's canal beneath the petroclinoid ligament, courses through the cavernous sinus (adjacent to the internal carotid artery), and enters the orbit through the superior orbital fissure within the common tendinous ring (annulus of Zinn). Because it has the longest intracranial subarachnoid course of any cranial nerve, it is exceptionally vulnerable to elevated intracranial pressure (a 'false localizing sign'). Abducens nerve palsy causes horizontal diplopia (worse on lateral gaze toward the affected side) and an inability to abduct the affected eye, resulting in convergent strabismus (esotropia).

\medterm{Abduction}-- A movement of a limb or body part in the coronal plane away from the median plane (midline) of the body. For the fingers, abduction is spreading them away from the longitudinal axis of the third (middle) digit; for the toes, it is movement away from the second digit. For the thumb, abduction is movement anteriorly away from the palm in the sagittal plane. Produced by muscles such as the middle deltoid and supraspinatus (arm), gluteus medius and minimus (thigh), and dorsal interossei (hand and foot).

\medterm{Accessory Nerve (Cranial Nerve XI)}-- The eleventh cranial nerve (CN XI), a purely motor nerve that innervates the sternocleidomastoid (SCM) and trapezius muscles. Although historically described as having cranial and spinal roots, modern neuroanatomy recognizes it as arising from lower motor neurons in the anterior horn of the upper cervical spinal cord (C1--C5/C6). These spinal rootlets unite, ascend through the foramen magnum into the posterior cranial fossa, and exit the skull through the jugular foramen alongside CN IX and CN X. It traverses the posterior triangle of the neck superficially between the posterior border of the SCM and the anterior border of the trapezius, embedded within the investing layer of deep cervical fascia. Its superficial course in the posterior neck makes it vulnerable to injury during lymph node biopsies, radical neck dissection, or penetrating trauma. Injury produces shoulder drop, winged scapula (lateral winging due to trapezius weakness, unlike medial winging from long thoracic nerve injury), and weakness in contralateral head turning and shoulder shrugging.

\medterm{Acetabulum}-- Cup-shaped lateral hip socket from fused ilium, ischium, and pubis. Articulates with the femoral head at the ball-and-socket hip joint. The acetabular labrum deepens the socket. Oriented with approximately 15 degrees anteversion and 45 degrees inclination. Fractures are serious high-energy trauma injuries.

\medterm{Achilles Tendon (Calcaneal Tendon)}-- The thickest and strongest tendon in the human body, connecting the gastrocnemius and soleus (triceps surae) muscles to the calcaneus (heel bone). It plays a crucial role in plantarflexion of the foot, enabling walking, running, and jumping. The tendon is vulnerable to tendonitis and acute rupture, which presents with a sudden "pop" and a positive Thompson test. Also known as the calcaneal tendon.

\medterm{Acromioclavicular Joint}-- A plane synovial joint between the acromion of the scapula and the lateral end of the clavicle. It permits gliding movements and is stabilized primarily by the acromioclavicular and coracoclavicular ligaments.

\medterm{Acromion}-- Bony process forming the shoulder's highest point, extending laterally from the scapular spine. Articulates with the clavicle at the acromioclavicular joint and arches over the glenohumeral joint. Attachment for deltoid and trapezius. The subacromial space contains the supraspinatus tendon and subacromial bursa. Its shape influences impingement syndrome susceptibility.

\medterm{Adam's Apple}-- The common name for the laryngeal prominence, the visible projection at the front of the neck formed by the angle where the two laminae of the thyroid cartilage meet. It protects the larynx and vocal cords, and it usually becomes more prominent during male puberty as the larynx enlarges under androgen influence, although it is present in people of all sexes. Its size does not determine voice quality or health.

\medterm{Adduction}-- A movement of a limb or body part in the coronal plane toward the median plane (midline) of the body. For the digits of the hand, adduction draws the fingers toward the longitudinal axis of the third (middle) digit; for the toes, it brings them toward the second digit. For the thumb, adduction brings it toward the radial border of the palm. Produced by muscles including pectoralis major and latissimus dorsi (arm), adductor longus, brevis, and magnus (thigh), and palmar interossei (hand).

\medterm{Adductor Muscles}-- A group of muscles situated in the medial compartment of the thigh that adduct and stabilize the hip: adductor longus, adductor brevis, adductor magnus, gracilis, and pectineus. They originate from the pubic and ischiopubic rami and insert along the linea aspera of the femur (gracilis inserts onto the anteromedial tibia at the pes anserinus). Innervation is provided primarily by the obturator nerve (L2--L4); pectineus is dual-innervated (predominantly by the femoral nerve, L2--L3, with variable obturator assistance), and adductor magnus is a hybrid muscle having an adductor part (obturator nerve) and a hamstring part (tibial division of the sciatic nerve, L4). Adductor longus is the most frequently injured muscle in athletic 'groin strains.' The adductor (Hunter) canal is an aponeurotic tunnel beneath the sartorius conveying the femoral artery, femoral vein, nerve to vastus medialis, and saphenous nerve into the popliteal fossa.

\medterm{Adenoids}-- Lymphoid tissue in the posterior wall and roof of the nasopharynx, part of Waldeyer ring of pharyngeal lymphoid tissue. Enlarged adenoids (adenoidal hypertrophy) in children cause nasal obstruction, mouth breathing, snoring, obstructive sleep apnea, serous otitis media (eustachian tube obstruction), and facial changes ('adenoid facies'). Inflammation (adenoiditis) is treated medically or with adenoidectomy. They enlarge in childhood and atrophy after puberty. Also called the pharyngeal tonsil.

\medterm{Adipose Tissue}-- A specialized connective tissue composed mainly of adipocytes that store energy as triglycerides. It provides thermal insulation, cushioning for organs, and endocrine regulation through hormones and signaling molecules such as leptin and adiponectin. White adipose tissue is the principal energy store in adults, whereas brown adipose tissue contains abundant mitochondria and generates heat through non-shivering thermogenesis, especially in newborns. Excess visceral adipose tissue is associated with insulin resistance, dyslipidemia, hypertension, and increased cardiometabolic risk.

\medterm{Adrenal Gland}-- Paired endocrine glands atop each kidney (suprarenal), composed of cortex and medulla. Cortex (three zones): zona glomerulosa (aldosterone, mineralocorticoids), zona fasciculata (cortisol, glucocorticoids), zona reticularis (androgens). Medulla: chromaffin cells secrete epinephrine and norepinephrine (catecholamines). Regulated by ACTH (cortex) and sympathetic nervous system (medulla). Conditions: Cushing syndrome, Addison disease, pheochromocytoma, Conn syndrome, congenital adrenal hyperplasia.

\medterm{Alveolus}-- The terminal air sac (200-300 microns) where gas exchange occurs. Approximately 300-500 million per lung, providing ~70 m2 surface area. Wall: type I pneumocytes (95\% surface, thin for diffusion), type II pneumocytes (5\%, secrete pulmonary surfactant reducing surface tension), alveolar macrophages (dust cells, phagocytosis). Capillary network in close apposition (blood-air barrier: type I pneumocyte, fused basement membrane, capillary endothelium). Surfactant deficiency causes RDS in preterm infants. Conditions: pneumonia, ARDS, pulmonary edema, emphysema (alveolar destruction), pulmonary fibrosis.

\medterm{Amygdala}-- Almond-shaped nucleus cluster in the anterior medial temporal lobe, part of the limbic system, primarily processing emotions (especially fear/anxiety). Receives thalamic/cortical sensory input; projects to hypothalamus, brainstem, and cortex for autonomic/behavioral fear responses. Essential for emotional memory (fear conditioning) and social signal processing. Lesions (Urbach-Wiethe disease) impair fear recognition. Hyperactive in anxiety disorders and PTSD.

\medterm{Anal Canal}-- The terminal 3-4 cm of the gastrointestinal tract, extending from the anorectal junction (puborectalis/pectinate line) to the anal verge. Above the pectinate line: lined by columnar epithelium, innervated by autonomic nerves, draining to portal system, supplied by superior rectal artery. Below the pectinate line: stratified squamous epithelium, somatic innervation (inferior rectal nerve—pain-sensitive), systemic drainage. Internal anal sphincter (smooth, involuntary, sympathetic/parasympathetic) and external anal sphincter (striated, voluntary, pudendal). Hemorrhoids, fissures, fistulae, and anorectal abscesses occur here.

\medterm{Anatomical Axis}-- An imaginary reference line around which a body part, joint, or bone rotates, or along which a structure extends. The three cardinal axes of the human body are: the transverse (horizontal/coronal) axis, around which flexion and extension occur; the anteroposterior (sagittal) axis, around which abduction and adduction occur; and the longitudinal (vertical) axis, around which medial and lateral rotation occur. In long bones, the mechanical axis connects the centers of the proximal and distal joints, which may diverge from the anatomical axis of the bone shaft.

\medterm{Anatomical Planes}-- Hypothetical imaginary flat two-dimensional geometric surfaces passing through the human body in the standard anatomical position, used as reference frameworks for describing anatomical locations, sections, movements, and radiological imaging (CT, MRI, ultrasound). The three cardinal orthogonal planes are: (1) Sagittal planes (including the median/midsagittal plane), dividing the body into right and left halves; (2) Coronal (frontal) planes, dividing the body into anterior (front) and posterior (back) portions; and (3) Transverse (axial/horizontal) planes, dividing the body into superior (upper) and inferior (lower) portions. Oblique planes pass at any non-orthogonal angle.

\medterm{Anatomical Position}-- The universal standardized reference posture of the human body used in all anatomical, surgical, and clinical descriptions to ensure unambiguous spatial terminology. It is defined as: an individual standing fully erect, facing forward, head level with eyes looking straight ahead toward the horizon, lower limbs together with feet parallel and directed forward, and upper limbs hanging relaxed at the sides with the palms facing forward (supinated) and the thumbs pointing away from the body. All relational terms (anterior, posterior, medial, lateral, superior, inferior) are defined strictly relative to this posture, regardless of the patient's actual physical position (e.g., supine, prone, or lateral decubitus).

\medterm{Anatomical Snuffbox}-- Triangular depression on the dorsal wrist, bounded by extensor pollicis longus medially and extensor pollicis brevis/abductor pollicis longus laterally. Floor contains the scaphoid, trapezium, and radial artery. Tenderness after a fall on an outstretched hand suggests scaphoid fracture. The radial artery pulse is palpable here.

\medterm{Ankle Joint}-- Uniaxial hinge synovial joint between distal tibia, lateral malleolus, and superior talus, allowing dorsiflexion/plantarflexion. Medial (deltoid) ligament resists eversion; lateral complex resists inversion. The ankle mortise provides bony stability. Ankle sprains are the most common sports injury (anterior talofibular ligament most commonly affected). Fractures classified by the Weber system.

\medterm{Antecubital Fossa}-- Triangular depression on the anterior elbow, bounded laterally by brachioradialis, medially by pronator teres, and superiorly by a line connecting the epicondyles. Floor formed by brachialis and supinator. Contains biceps tendon, brachial artery, and median nerve. Primary site for venipuncture and blood pressure measurement.

\medterm{Anterior}-- A directional term indicating a position situated at, toward, or closer to the front (ventral) surface of the body or an organ. In human anatomy, because of upright bipedal posture, anterior is generally synonymous with ventral. Opposed to posterior (dorsal). Examples include the sternum lying anterior to the heart, and the patella situated on the anterior surface of the knee.

\medterm{Anterior Cruciate Ligament (ACL)}-- A primary stabilizing intra-articular, extrasynovial ligament of the knee joint. It originates from the posteromedial aspect of the lateral femoral condyle in the intercondylar notch and passes anteriorly, medially, and distally to insert onto the anterior intercondylar area of the tibial plateau, blending with the anterior horn of the lateral meniscus. Composed of two functional bundles: the anteromedial bundle (tight in flexion) and the posterolateral bundle (tight in extension). Its primary function is to resist anterior translation of the tibia relative to the femur (providing ~85 percent of total resistance) and secondary restraint against internal tibial rotation and varus/valgus stress. ACL rupture is among the most common major sports injuries, typically occurring from non-contact deceleration, cutting, or landing with a valgus knee and internally rotated tibia. Clinical diagnosis is made using the Lachman test (most sensitive), anterior drawer test, and pivot-shift test; confirmed by MRI.

\medterm{Anterior Tibial Artery}-- Arises from the popliteal artery, passes through the anterior intermuscular septum to the anterior compartment, continuing to the ankle as the dorsalis pedis artery. Supplies anterior leg and dorsal foot muscles. Palpable anterior to the lateral malleolus as the dorsalis pedis pulse. Important for vascular assessment; vulnerable to injury in tibial fractures and anterior compartment syndrome.

\medterm{Anus}-- The terminal opening of the gastrointestinal tract through which feces are expelled. It is surrounded by the internal smooth-muscle sphincter and external skeletal-muscle sphincter.

\medterm{Aorta}-- Largest artery, from the left ventricle distributing oxygenated blood systemically. Approximately 30 cm long, 2.5 cm diameter. Four segments: ascending aorta, aortic arch (brachiocephalic trunk, left common carotid, left subclavian), descending thoracic aorta, descending abdominal aorta (bifurcating into common iliacs at L4). Three layers: intima, media (elastic tissue), adventitia. Aneurysm and dissection are life-threatening.

\medterm{Aortic Arch}-- Curves from second right sternocostal joint superiorly and left, then posteriorly and inferiorly to T4. From right to left: brachiocephalic trunk, left common carotid, left subclavian. Related to left recurrent laryngeal nerve (hooking under it), trachea, esophagus, left phrenic nerve. Arch aneurysms may cause hoarseness and SVC syndrome.

\medterm{Aortic Valve}-- Semilunar valve located between the left ventricle and the ascending aorta, consisting of three pocket-like cusps: right, left, and posterior (non-coronary). Superior to each cusp is an aortic sinus (sinus of Valsalva); the right and left coronary arteries arise from the right and left aortic sinuses, respectively, whereas the posterior sinus gives rise to no coronary vessel. The valve opens passively during ventricular systole and closes during diastole, preventing systemic regurgitation. Aortic stenosis is the most prevalent valvular disease in elderly patients, typically resulting from calcific degeneration and presenting with angina, syncope, heart failure, and a harsh systolic crescendo-decrescendo murmur radiating to the carotids. Bicuspid aortic valve is the most common congenital cardiac valvular abnormality, predisposing to premature calcific stenosis, regurgitation, and ascending aortic aneurysm.

\medterm{Appendix}-- A blind-ended tubular diverticulum of the cecum, typically 6-10 cm long, rich in lymphoid tissue. Its position is variable (retrocecal, pelvic, or retroileal); the base usually projects near McBurney's point, one-third of the distance from the anterior superior iliac spine to the umbilicus. Innervated by T10 (visceral pain initially periumbilical, later often localizing to the right lower quadrant). Acute appendicitis may cause perforation, abscess, and peritonitis; appendectomy is definitive treatment.

\medterm{Arachnoid Mater}-- The middle of the three meninges, a delicate avascular membrane lying between the dura mater and pia mater. The subarachnoid space between arachnoid and pia contains cerebrospinal fluid (CSF) and major vessels. Arachnoid granulations (villi) project into the dural venous sinuses and mediate CSF resorption. The arachnoid is separated from the dura by the subdural space (site of subdural hematoma). Inflammation (meningitis) or hemorrhage (subarachnoid hemorrhage) in this space has serious consequences.

\medterm{Arrector Pili Muscle}-- A small band of involuntary smooth muscle attached to a hair follicle and the superficial dermis. Sympathetic stimulation causes it to contract, pulling the hair upright and producing temporary piloerection, commonly called goosebumps. Contraction also compresses nearby sebaceous glands and may help release sebum onto the skin surface.

\medterm{Arteriole}-- A small artery that delivers blood to a capillary bed. Arteriolar smooth muscle regulates vascular resistance and controls tissue perfusion and systemic blood pressure.

\medterm{Artery}-- A blood vessel that carries oxygenated blood away from the heart to the tissues, except for the pulmonary arteries which carry deoxygenated blood to the lungs. Arteries have thick, elastic, muscular walls that withstand high pressure and distribute blood throughout the body. The aorta is the largest artery. Arterial structure: tunica intima (endothelium), tunica media (smooth muscle and elastic fibers), tunica adventitia (connective tissue). Arteries branch into smaller arterioles and eventually capillaries. Diseases: atherosclerosis, aneurysms, arteritis, emboli.

\medterm{Ascending Aorta}-- From left ventricle at the aortic valve to the sternal angle (approximately 5 cm), enclosed in pericardium. Gives rise to right/left coronary arteries at the aortic sinuses. Most common site of aneurysm and dissection, especially in Marfan syndrome, bicuspid aortic valve, and hypertension. Aneurysms may compress adjacent structures and cause aortic regurgitation.

\medterm{Ascending Colon}-- Extends from cecum to hepatic flexure below the liver, approximately 15-20 cm, retroperitoneal (fixed to posterior abdominal wall). Anteriorly: abdominal wall. Posteriorly: right kidney. Absorbs water/electrolytes and compacts fecal matter. Tumors may present as a palpable right-sided mass. The right colic flexure marks its transition to the transverse colon.

\medterm{Atlas (C1)}-- First cervical vertebra articulating with occipital condyles at the atlanto-occipital joint for nodding movement. Lacks a body and spinous process; has anterior/posterior arches and lateral masses. Transverse foramina transmit vertebral arteries. The anterior arch has a facet for the axis dens. Commonly fractured in axial loading injuries such as diving accidents.

\medterm{Atrium}-- One of the two upper chambers of the heart that receive blood returning to the heart. The right atrium receives deoxygenated blood from the superior and inferior venae cavae and the coronary sinus; the left atrium receives oxygenated blood from the four pulmonary veins. The interatrial septum separates the two atria; the fossa ovalis marks the site of the fetal foramen ovale. Atrial contraction (atrial kick) contributes 15--25\% of ventricular filling. The atrial appendages (auricles) increase atrial capacity. Conditions: atrial fibrillation, atrial septal defect, atrial myxoma.

\medterm{Axilla}-- Pyramidal space between the upper lateral thorax and medial arm. Bounded by anterior fold (pectoralis major/minor), posterior fold (latissimus dorsi/teres major), medial wall (ribs/serratus anterior), and lateral wall (humerus). Contains axillary artery/vein, brachial plexus cords, and lymph nodes. Axillary nodes are significant in breast cancer staging.

\medterm{Axillary Artery}-- The continuation of the subclavian artery, running from the outer border of the first rib to the lower border of teres major, where it becomes the brachial artery. Divided into three parts by the pectoralis minor. Branches (part 1): superior thoracic; (part 2): thoracoacromial, lateral thoracic; (part 3): subscapular, anterior and posterior circumflex humeral. Supplies the shoulder, axilla, and upper limb. Relations: axillary vein (anterior/medial), brachial plexus cords. Site of axillary artery catheterization; injury causes axillary hematoma and upper limb ischemia; circumflex humeral vessels are important to the humeral head blood supply.

\medterm{Axillary Nerve}-- From posterior cord (C5-C6), passing through the quadrangular space with the posterior circumflex humeral artery. Innervates deltoid and teres minor; provides sensation to the regimental badge area (lateral deltoid skin). Vulnerable in anterior shoulder dislocations and surgical neck humeral fractures. Injury causes loss of lateral deltoid sensation and weakness of abduction beyond 15 degrees.

\medterm{Axis (C2)}-- Second cervical vertebra with the dens (odontoid process) projecting superiorly, serving as a pivot for atlas rotation, allowing approximately 50 percent of cervical rotation. Has a large spinous process and flat superior articular facets. Dens fractures (Types I-III) may cause atlanto-axial instability and spinal cord compression, often requiring surgical fixation.

\medterm{Axon}-- The single long process of a neuron that conducts action potentials away from the cell body toward the axon terminal for synaptic transmission. Axons vary in length (microscopic to >1 m for sciatic motor fibers) and are wrapped by Schwann cells (PNS) or oligodendrocytes (CNS) to form myelin, which enables saltatory conduction and speeds transmission. Axons end in synaptic boutons releasing neurotransmitters. Axonal injury in the CNS does not regenerate effectively; Wallerian degeneration occurs distal to injury. Axonal transport (anterograde and retrograde) carries organelles and proteins.

\medterm{Azygos Vein}-- An unpaired vein ascending on the right side of the vertebral column that drains the posterior thoracic wall and empties into the superior vena cava. It provides an important collateral pathway between the superior and inferior venae cavae.

\medterm{Basal Ganglia}-- Deep gray matter nuclei in the cerebral hemispheres for motor control, cognition, and emotion. Key structures: caudate nucleus, putamen, globus pallidus (internal/external segments), subthalamic nucleus, substantia nigra. Direct pathway facilitates movement; indirect inhibits it. Receive cortical/thalamic input, project back to cortex via thalamus. Diseases: Parkinson's (substantia nigra dopamine depletion), Huntington's (caudate atrophy), hemiballismus (subthalamic lesion).

\medterm{Basement Membrane}-- A specialized, dense, thin sheet of extracellular matrix that underlies all epithelial layers and surrounds muscle fibers, adipocytes, and Schwann cells, anchoring them to the underlying connective tissue. It consists of two ultrastructural layers: the basal lamina (secreted by epithelial cells, composed of type IV collagen, laminin, enactin/nidogen, and perlecan/heparan sulfate proteoglycans) and the reticular lamina (secreted by connective tissue fibroblasts, containing type III collagen/reticulin and type VII anchoring fibrils). Functions include mechanical anchoring via hemidesmosomes, cellular polarity maintenance, semipermeable size- and charge-selective filtration (critical in the renal glomerular basement membrane), tissue compartmentalization, and regulation of cell growth, differentiation, and repair. Pathologies include Goodpasture syndrome (autoantibodies against type IV collagen alpha-3 chain), Alport syndrome (hereditary type IV collagen mutations causing nephritis and hearing loss), bullous pemphigoid (autoantibodies against hemidesmosomal BP180/BP230 at the basement membrane), and tumor invasion (epithelial carcinomas must breach the basement membrane to metastasize).

\medterm{Biceps Brachii Muscle}-- A prominent two-headed muscle in the anterior compartment of the arm. The long head originates from the supraglenoid tubercle of the scapula and the glenoid labrum, coursing intracapsularly across the humeral head and exiting via the intertubercular (bicipital) groove; the short head originates from the tip of the coracoid process alongside the coracobrachialis. Both heads fuse into a common distal tendon that inserts onto the radial tuberosity, with a medial aponeurotic expansion (bicipital aponeurosis) inserting into the deep fascia of the forearm. Innervated by the musculocutaneous nerve (C5--C7, primarily C5--C6). It functions as the body's most powerful supinator of the flexed forearm, a major elbow flexor, and a secondary shoulder flexor and stabilizer. Rupture of the proximal long head tendon causes distal retraction ('Popeye deformity'), while distal tendon rupture severely weakens supination and requires urgent surgical repair. Bicipital tenosynovitis is a common source of anterior shoulder pain.

\medterm{Bile Duct}-- A channel that carries bile from the liver and gallbladder to the duodenum. The right and left hepatic ducts unite as the common hepatic duct; union with the cystic duct forms the common bile duct, which usually joins the pancreatic duct before entering the duodenum at the major duodenal papilla.

\medterm{Bone}-- A rigid, mineralised connective tissue that forms the skeleton, providing structural support, protection of vital organs, hematopoiesis (bone marrow), and mineral storage (calcium, phosphate). Bone composition: organic matrix (90\% type I collagen) and inorganic mineral (hydroxyapatite, calcium phosphate). Types: cortical (compact, 80\% of bone mass) and cancellous (trabecular, 20\%). Cells: osteoblasts (bone formation), osteocytes (mechanosensation), osteoclasts (bone resorption). Ossification: intramembranous (flat bones) and endochondral (long bones). Remodelling: continuous cycle of resorption and formation.

\medterm{Bone Marrow}-- Vascular tissue within the medullary cavities and cancellous bone. Red marrow produces blood cells, whereas yellow marrow primarily stores fat and can revert to red marrow when demand increases.

\medterm{Brachial Artery}-- The continuation of the axillary artery, running from the lower border of teres major to the cubital fossa, where it divides into radial and ulnar arteries. Lies medial to the biceps; palpated for blood pressure and brachial pulse. Branches: profunda brachii (deep brachial), superior/inferior ulnar collaterals. Site for blood pressure measurement, arterial access (catheterization), and the median nerve travels alongside. Trauma here causes forearm ischemia; collateral circulation around the elbow (anastomosis) limits ischemic damage.

\medterm{Brachialis Muscle}-- A powerful, broad muscle lying deep to the biceps brachii in the lower two-thirds of the anterior arm. It originates from the distal half of the anterior surface of the humerus and adjacent intermuscular septa, descending across the anterior capsule of the elbow joint to insert onto the ulnar tuberosity and the coronoid process of the ulna. It is innervated primarily by the musculocutaneous nerve (C5--C7); its lateral portion also receives sensory and proprioceptive twigs from the radial nerve (C7). Because it inserts onto the ulna rather than the radius, its mechanical advantage is completely independent of forearm rotation, making it the primary 'workhorse' flexor of the elbow in all positions (pronation, supination, or midway). Trauma to the elbow or repetitive stress can lead to brachialis tendinopathy or post-traumatic myositis ossificans.

\medterm{Brachial Plexus}-- Nerve network (C5-T1) formed by roots, trunks, divisions, and cords, giving rise to upper limb nerves. Roots unite into superior/middle/inferior trunks, dividing into anterior/posterior divisions, recombining into lateral/posterior/medial cords. Major branches: musculocutaneous, axillary, radial, median, ulnar. Birth injuries cause Erb's palsy (C5-C6) or Klumpke's palsy (C8-T1). Vulnerable to shoulder trauma and thoracic outlet syndrome.

\medterm{Brachial Plexus Cords}-- Formed by divisions recombining behind the pectoralis minor. Lateral cord (anterior divisions of upper/middle trunks), posterior cord (all three posterior divisions), medial cord (anterior division of lower trunk). Named for relationship to the axillary artery. Lateral cord: lateral pectoral nerve, contributes to median nerve. Posterior cord: subscapular, thoracodorsal, axillary, radial nerves. Medial cord: medial pectoral, cutaneous, and ulnar nerves.

\medterm{Brachial Plexus Roots}-- Ventral rami of C5-T1 emerging from intervertebral foramina, passing between anterior and middle scalene muscles (interscalene triangle). C5-C6 form the upper trunk, C7 continues as middle trunk, C8-T1 form the lower trunk. Root injuries cause specific dermatome/myotome deficits. C5 avulsion causes loss of deltoid/biceps function. Roots are vulnerable in birth trauma, motorcycle accidents, and disc herniation radiculopathy.

\medterm{Brachial Plexus Trunks}-- Formed in the posterior neck triangle. Upper (C5-C6), middle (C7), and lower (C8-T1) trunks each divide into anterior/posterior divisions behind the clavicle. Upper trunk is most commonly injured in birth injuries (Erb-Duchenne palsy), causing waiter's tip position. Lower trunk injury (Klumpke's palsy from upward arm traction) causes claw hand.

\medterm{Brachiocephalic Trunk}-- The first and largest branch of the aortic arch. It ascends to the right sternoclavicular joint, where it divides into the right common carotid and right subclavian arteries.

\medterm{Brain}-- The central organ of the nervous system, enclosed in the cranial cavity and composed of the cerebrum, cerebellum, and brainstem. The brain receives, processes, and integrates sensory information; initiates and coordinates motor responses; and serves as the seat of cognition, emotion, memory, and consciousness. The cerebrum (largest part) has two hemispheres connected by the corpus callosum, with four lobes (frontal, parietal, temporal, occipital) responsible for higher functions. The cerebellum coordinates movement, balance, and motor learning. The brainstem (midbrain, pons, medulla) controls vital autonomic functions (breathing, heart rate, blood pressure) and contains cranial nerve nuclei. Protected by the skull, meninges (dura, arachnoid, pia), and cerebrospinal fluid. Blood supply via internal carotid and vertebral arteries (circle of Willis). Conditions: stroke, tumor, trauma, neurodegenerative diseases, epilepsy, infection.

\medterm{Brainstem}-- Connects cerebrum to spinal cord, with midbrain (superior), pons (middle), and medulla oblongata (inferior). Contains cranial nerve nuclei (III-XII), ascending/descending tracts, and vital autonomic centers for cardiac, respiratory, and vasomotor function. Reticular formation regulates consciousness and sleep-wake cycles. Contains all brain-spinal cord pathways. Lesions are often devastating and fatal, causing respiratory arrest, cardiac arrhythmias, and bilateral cranial nerve deficits.

\medterm{Bronchi}-- The major airways conducting air from the trachea into the lungs. The trachea bifurcates at the carina (T4-T5) into right and left main bronchi. The right main bronchus is wider, shorter, and more vertical—aspirated foreign bodies preferentially lodge here. Main bronchi divide into lobar (secondary), segmental (tertiary), and progressively smaller bronchi and bronchioles. Bronchi contain cartilage and submucosal glands; bronchioles lack cartilage. Bronchial walls contain smooth muscle (bronchoconstriction in asthma), goblet cells, and ciliated epithelium that clears mucus. Blood supply: bronchial arteries.

\medterm{Bronchiole}-- Small airways (less than 2 mm diameter) without cartilage, branching from segmental bronchi to terminal bronchioles (conducting zone) to respiratory bronchioles (transitional, alveoli bud from walls) to alveolar ducts to alveolar sacs to alveoli. Lined by ciliated simple cuboidal epithelium (Clara cells in terminal bronchioles, secrete surfactant components). No goblet cells, no submucosal glands. Smooth muscle in walls regulates airflow (bronchoconstriction/dilation). Conditions: bronchiolitis, asthma, bronchiolitis obliterans, small airway disease.

\medterm{Bronchus}-- One of the two main branches of the trachea that conduct air into the lungs. The right main bronchus is wider, shorter, and more vertical than the left, making it more susceptible to aspiration. Each main bronchus divides into lobar bronchi (secondary), segmental bronchi (tertiary), and progressively smaller branches (bronchioles). Bronchial walls contain ciliated pseudostratified columnar epithelium, smooth muscle, cartilage plates, and mucous glands. Bronchoconstriction occurs in asthma. Bronchiectasis is irreversible dilatation. Bronchoscopy allows visualization and biopsy.

\medterm{Cadaver}-- A dead human body used for anatomical dissection, surgical training, and medical research. Cadavers are preserved through embalming (arterial injection of formaldehyde-based fixatives) to prevent decomposition and maintain tissue integrity. Anatomical dissection of cadavers remains a cornerstone of medical education, providing three-dimensional understanding of human anatomy, anatomical variation, and surgical approaches. Whole-body donation programs supply cadavers to medical schools; ethical considerations include informed consent, respectful treatment, and compliance with anatomical gift laws.

\medterm{Calcaneus}-- Largest tarsal bone forming the heel. Articulates with talus proximally and cuboid anteriorly. Lever for the calf muscles via the Achilles tendon on its posterior tuberosity. Predominantly trabecular bone, the most commonly fractured tarsal bone from falls from height. Bohler's angle assesses fractures on lateral radiographs.

\medterm{Cantlie Line}-- Functional liver division extending from the gallbladder fossa to the IVC, corresponding to the middle hepatic vein plane. Unlike the falciform ligament (anatomical division), it represents the true functional blood supply/biliary drainage division. Right hepatic artery and portal vein supply the right half; left supply the left. Basis for anatomic hepatic segmentectomy and liver transplantation.

\medterm{Capillary}-- The smallest blood vessel, consisting primarily of a single endothelial cell layer and basement membrane. Capillaries connect arterioles with venules and are the principal site of exchange between blood and tissues.

\medterm{Capillary Bed}-- A network of interconnected capillaries that links arterioles to venules within a tissue. It provides a large surface area for exchange of oxygen, carbon dioxide, nutrients, hormones, and metabolic waste between blood and interstitial fluid. Precapillary sphincters and arteriolar tone regulate how much blood enters the bed according to local metabolic demand.

\medterm{Capitate}-- Largest carpal bone, centrally located in the distal row as the carpal keystone. Rounded head articulates with scaphoid and lunate; distal surface with the third metacarpal. Last carpal bone to ossify. Wrist joint axis, bearing significant axial load. Avascular necrosis may occur with fractures similar to scaphoid injuries.

\medterm{Cardiac Muscle}-- Involuntary, striated muscle forming the myocardium. Cardiac muscle cells are electrically connected by intercalated discs, allowing coordinated contraction of the heart.

\medterm{Carotid Artery}-- The paired arterial system supplying the head and neck. Each common carotid artery divides into internal and external carotid arteries; the internal carotid supplies the brain and orbit, while the external carotid supplies most external head and neck structures.

\medterm{Carpal Tunnel}-- Narrow passageway on the palmar wrist, bounded by carpal bones dorsally and the flexor retinaculum as the roof. Transmits the median nerve and flexor tendons. Carpal tunnel syndrome causes numbness in the lateral 3.5 digits and thenar weakness. It is the most common entrapment neuropathy, associated with repetitive motion, diabetes, hypothyroidism, and pregnancy.

\medterm{Cartilage}-- A tough, flexible avascular connective tissue found in joints, the ear, nose, rib cage, and intervertebral discs. Composed of chondrocytes embedded in an extracellular matrix of collagen fibers, proteoglycans, and water, providing resilience and compressive strength. Types: hyaline cartilage (most common, covers articular surfaces of joints, costal cartilages), elastic cartilage (ear, epiglottis), and fibrocartilage (intervertebral discs, menisci, pubic symphysis). Cartilage lacks blood vessels, nerves, and lymphatics; it receives nutrition by diffusion from synovial fluid, limiting healing capacity.

\medterm{Caudal}-- A directional term relating to or situated toward the tail (coccyx) or inferior aspect of the trunk in humans. In embryology and comparative anatomy, caudal denotes the tail-end of the embryonic axis, opposed to rostral/cranial. In regional anatomy, it is frequently used synonymously with inferior (e.g., the caudal end of the spinal cord terminates at L1--L2 as the conus medullaris).

\medterm{Cecum}-- First and widest large intestinal part in the right iliac fossa below the ileocecal valve, a blind-ended pouch approximately 6 cm long. Appendix arises from its posteromedial wall approximately 2 cm below the valve, containing abundant lymphoid tissue. Cecal volvulus causes large bowel obstruction. Absorbs water/electrolytes and compacts fecal matter.

\medterm{Celiac Artery}-- The first major ventral branch of the abdominal aorta at T12, arising from the aorta just below the diaphragm and supplying the foregut (stomach, proximal duodenum, liver, gallbladder, spleen, pancreas). It trifurcates into the left gastric, splenic, and common hepatic arteries, with rich collateral circulation between these branches. Occlusion/stenosis can cause chronic mesenteric ischemia. The celiac trunk is encased in the celiac plexus; it is important in vascular surgery (celiac axis clamping) and anatomy of the pancreatic bed.

\medterm{Celiac Trunk}-- First major abdominal aortic branch at T12-L1, the foregut artery. Divides into left gastric (lesser curvature), splenic (spleen, pancreas, greater curvature), and common hepatic (proper hepatic/gastroduodenal) arteries. Surrounded by the celiac plexus. Compression by the median arcuate ligament causes celiac trunk compression syndrome with postprandial pain.

\medterm{Cell}-- The fundamental structural and functional unit of all living organisms. Human cells are eukaryotic, characterized by a nucleus containing DNA and membrane-bound organelles. Major components: cell membrane (phospholipid bilayer with embedded proteins, regulates transport and signaling), cytoplasm (cytosol + organelles), nucleus (chromosomes, nucleolus), mitochondria (ATP production through oxidative phosphorylation), endoplasmic reticulum (protein\/lipid synthesis), Golgi apparatus (modification and sorting of proteins), lysosomes (intracellular digestion), peroxisomes (fatty acid oxidation), and cytoskeleton (microfilaments, intermediate filaments, microtubules--maintains shape and enables movement). Cells differentiate into ~200 specialized types with distinct functions (neurons, myocytes, hepatocytes, etc.). Cell division occurs via mitosis (somatic cells) and meiosis (gametes). Cell death occurs by apoptosis (programmed) or necrosis (pathologic).

\medterm{Central}-- A directional term indicating a location situated at or toward the center, interior, or core of the body, a region, or an organ. Opposed to peripheral. Used in functional anatomy to designate major control hubs (e.g., the Central Nervous System, central venous catheterization in the internal jugular or subclavian vein, and the central canal of the spinal cord).

\medterm{Cerebellum}-- The second largest part of the brain (after the cerebrum) and the largest structure located within the posterior cranial fossa, situated posterior to the pons and medulla oblongata and separated from the occipital lobes by the tentorium cerebelli. Although accounting for only approximately 10 percent of total brain volume, it contains over 50 percent of all central neurons. It consists of two lateral hemispheres joined by a midline vermis and is divided into three functional zones: the vestibulocerebellum (archicerebellum/flocculonodular lobe, governing balance and vestibulo-ocular reflexes), the spinocerebellum (paleocerebellum, coordinating posture and axial/limb tone), and the cerebrocerebellum (neocerebellum, planning and fine-tuning voluntary skilled movements and motor learning). Afferent input arrives via mossy and climbing fibers through the cerebellar peduncles; efferent output leaves via deep cerebellar nuclei (dentate, emboliform, globose, fastigial) projecting to the red nucleus, thalamus, and motor cortex. Cerebellar lesions produce ipsilateral deficits characterized by ataxia, dysmetria (past-pointing), dysdiadochokinesia, intention tremor, nystagmus, and scanning staccato speech.

\medterm{Cerebral Cortex}-- The folded outer layer of gray matter covering the cerebral hemispheres. It contains interconnected networks responsible for conscious perception, voluntary movement, language, memory, and higher cognition.

\medterm{Cerebral Lobes}-- Four lobes divided by sulci. Frontal (anterior to central sulcus): executive function, motor control (precentral gyrus), Broca's area. Parietal (between central and parieto-occipital sulci): somatosensory (postcentral gyrus), spatial awareness. Temporal (below lateral sulcus): auditory processing, memory (hippocampus), Wernicke's area. Occipital (posterior): visual processing. Insula (deep within lateral sulcus): taste, emotion, interoception.

\medterm{Cerebral Peduncle}-- Paired structures of the ventral midbrain carrying descending motor tracts (corticospinal, corticobulbar, corticopontine) connecting the cerebrum to the pons and below. Lesions produce contralateral hemiparesis; Weber syndrome (peduncle + CN III) causes ipsilateral CN III palsy with contralateral hemiparesis. The peduncles are separated by the interpeduncular fossa and flanked by the substantia nigra (dopaminergic, degenerated in Parkinson disease).

\medterm{Cerebrospinal Fluid}-- A clear, colorless, ultrafiltered bodily fluid circulating within the cerebral ventricles, central canal of the spinal cord, and the subarachnoid space surrounding the brain and spinal cord. It is produced continuously at a rate of approximately 500 mL/day (0.35 mL/min; total volume ~150 mL) predominantly by the choroid plexus of the lateral, third, and fourth ventricles. CSF flows from the lateral ventricles through the interventricular foramina of Monro into the third ventricle, passes down the cerebral aqueduct of Sylvius into the fourth ventricle, and exits into the subarachnoid space via the midline foramen of Magendie and paired lateral foramina of Luschka. It circulates around the brain and spinal cord and is resorbed into dural venous sinuses (principally the superior sagittal sinus) via arachnoid granulations (Pacchionian bodies) driven by hydrostatic pressure gradients. It provides physical buoyancy (reducing apparent brain weight from 1400 g to ~50 g), mechanical protection against trauma, and biochemical regulation/waste removal (glymphatic system). Normal opening pressure on lumbar puncture is 10--20 cm H2O. Disruption of flow or resorption leads to hydrocephalus (communicating or non-communicating); CSF analysis is essential in diagnosing meningitis, subarachnoid hemorrhage, and multiple sclerosis.

\medterm{Cerebrum}-- The largest part of the brain, comprising two cerebral hemispheres connected by the corpus callosum. Each hemisphere has four lobes: frontal (executive function, motor, Broca's area), parietal (somatosensory, spatial awareness), temporal (auditory, memory, Wernicke's area), occipital (visual processing). Cerebral cortex (gray matter) over white matter (myelinated tracts). Basal ganglia deep within. Lateral ventricles produce CSF. Blood supply: anterior, middle, posterior cerebral arteries (circle of Willis). Conditions: stroke, tumor, dementia, epilepsy, trauma.

\medterm{Cervical Vertebrae (Cervical Spine, C1-C7)}-- The seven uppermost vertebrae of the vertebral column (C1--C7) forming the cervical spine, supporting the skull and permitting extensive movement of the head and neck. Typical cervical vertebrae (C3--C6) feature small, broad, oval bodies with uncinate processes (forming the uncovertebral joints of Luschka), large triangular vertebral foramina accommodating the cervical enlargement of the spinal cord, bifid spinous processes, and transverse processes containing transverse foramina (foramina transversaria) that transmit the vertebral artery, vertebral veins, and sympathetic plexuses (C1--C6 transmit the vertebral artery; C7 transmits only accessory vertebral veins). Three vertebrae are atypical: C1 (atlas, ring-shaped without body or spinous process, articulates with occipital condyles for nodding movements), C2 (axis, possessing the dens/odontoid process serving as a pivot for atlas rotation), and C7 (vertebra prominens, featuring a long, non-bifid, easily palpable spinous process and a transition toward thoracic morphology). Clinical conditions include cervical radiculopathy (nerve root compression via disc herniation or osteophytes), cervical spondylotic myelopathy (spinal cord compression), whiplash injuries, and high-energy fractures (Jefferson burst fracture of C1, Hangman's fracture of C2).

\medterm{Cervix}-- The inferior, cylindrical portion of the uterus that projects into the vagina. Its canal connects the uterine cavity with the vagina through the internal and external os.

\medterm{Chest Wall}-- The musculoskeletal boundary of the thoracic cavity, formed by the skin, fascia, muscles, ribs, costal cartilages, sternum, thoracic vertebrae, and associated neurovascular structures. It protects the thoracic organs and participates in breathing through coordinated movement of the ribs and respiratory muscles.

\medterm{Chordae Tendineae}-- Thin tendinous cords connecting the papillary muscles to the free edges of the atrioventricular valve leaflets (mitral and tricuspid). They tether the leaflets and, with papillary muscle contraction during ventricular systole, prevent valve prolapse and regurgitation. Rupture (from papillary muscle infarction) causes acute severe mitral regurgitation and cardiogenic shock. Chordal dysfunction contributes to mitral valve prolapse and functional MR.

\medterm{Choroid}-- The vascular layer of the eye between the sclera and retina, part of the uveal tract (with ciliary body and iris). Contains dense capillary network (choriocapillaris) supplying outer retina (photoreceptors). Pigmented (melanin) to absorb stray light. Bruchs membrane separates choroid from retinal pigment epithelium. Conditions: choroidal nevus, melanoma, neovascularization (wet AMD), choroiditis, central serous chorioretinopathy.

\medterm{Ciliary Body}-- Circumferential structure between iris and choroid, with ciliary muscle (accommodation control via zonular fibers) and ciliary processes (aqueous humor production). Part of the uveal tract (iris, ciliary body, choroid). Parasympathetic CN III fibers cause muscle contraction, relaxing zonular fibers and allowing lens convexity for near vision. Ciliary processes secrete aqueous humor into the posterior chamber, flowing through the pupil to the anterior chamber and draining at the trabecular meshwork and Schlemm's canal.

\medterm{Circle of Willis}-- A ring-like arterial anastomosis at the base of the brain that provides collateral blood flow between the anterior (internal carotid) and posterior (vertebrobasilar) circulations. It is formed by the anterior cerebral, anterior communicating, internal carotid, posterior cerebral, and posterior communicating arteries. Aneurysms commonly occur at its junctions (berry aneurysms), which can rupture and cause subarachnoid hemorrhage.

\medterm{Clavicle}-- Only horizontally oriented long bone, serving as a shoulder-trunk strut. Medial end articulates with the manubrium (sternoclavicular joint); lateral end with the acromion (acromioclavicular joint). Most common fracture at the middle-lateral third junction from falls on an outstretched hand. Protects the underlying brachial plexus and subclavian vessels.

\medterm{Coccyx}-- Terminal vertebral column part with 3-5 fused vertebrae. Articulates with the sacrum, allowing limited flexion important during childbirth. Provides attachment for gluteus maximus, coccygeus, and levator ani (pelvic floor). A common site of pain (coccydynia) following falls, direct trauma, and prolonged sitting.

\medterm{Cochlea}-- Snail-shaped hearing organ in the inner ear, coiled approximately 2.5 turns around the modiolus. Three fluid-filled chambers: scala vestibuli (perilymph), scala media (endolymph, organ of Corti), scala tympani (perilymph). The organ of Corti has inner/outer hair cells transducing mechanical vibrations into electrical signals via stereocilia deflection. Basilar membrane vibrates differentially along its length (tonotopic: high frequencies at base, low at apex). Hair cell damage from noise, ototoxicity, and ageing causes irreversible sensorineural hearing loss, most commonly affecting high frequencies.

\medterm{Collecting Duct}-- A tubular structure in the kidney that receives urine from distal convoluted tubules and carries it through the medulla to the renal papilla. Its permeability to water is regulated by antidiuretic hormone.

\medterm{Colon}-- The largest part of the large intestine, extending from the cecum to the rectum. Anatomically divided into ascending colon (right side), transverse colon (across the abdomen), descending colon (left side), and sigmoid colon (S-shaped, leading to the rectum). Functions: absorption of water and electrolytes, fermentation of undigested carbohydrate by gut microbiota, formation and storage of feces. The colon has a distinct appearance with taeniae coli (three longitudinal muscle bands), haustra (sacculations), and appendices epiploicae (fat appendages). Common diseases: diverticulosis, colitis, colorectal cancer, irritable bowel syndrome.

\medterm{Common Bile Duct}-- The duct formed by union of the common hepatic and cystic ducts. It conveys bile to the duodenum, usually joining the pancreatic duct at the hepatopancreatic ampulla.

\medterm{Common Iliac Arteries}-- Terminal abdominal aortic branches at L4, each approximately 4 cm long, dividing into internal and external iliac arteries. They supply the pelvis and lower extremities. The right common iliac artery crosses anterior to the right common iliac vein. Aneurysms may occur with abdominal aortic aneurysm and compress the ureters, causing hydronephrosis. The corresponding common iliac veins unite to form the inferior vena cava at approximately L5.

\medterm{Common Peroneal Nerve}-- Sciatic division (L4-S2) winding around the fibular neck, the body's most vulnerable major nerve to external compression. Divides into superficial peroneal (peroneus longus/brevis, dorsal foot sensation) and deep peroneal (tibialis anterior, extensors, first web space sensation). Injury causes foot drop (loss of dorsiflexion/eversion) with dorsal foot sensory loss.

\medterm{Connective Tissue}-- One of the four basic animal tissue types (along with epithelial, muscle, and nervous tissue), originating from embryonic mesoderm and mesenchyme, that provides structural support, binding, metabolic storage, and defense throughout the body. It consists of cells widely separated by an abundant extracellular matrix (ECM) composed of ground substance (glycosaminoglycans like hyaluronic acid, proteoglycans, glycoproteins) and protein fibers (type I/II/III collagen fibers providing tensile strength, elastic fibers providing elasticity, and reticular fibers providing stroma). Resident cells include fibroblasts (primary ECM producers), adipocytes, mast cells, and tissue macrophages; transient cells include leukocytes. Classified into connective tissue proper (loose/areolar, dense regular in tendons/ligaments, dense irregular in dermis), specialized connective tissue (adipose, cartilage, bone, blood, hemopoietic tissue), and embryonic connective tissue (mesenchyme, Wharton's jelly). Pathologies include Marfan syndrome (fibrillin-1 mutation), Ehlers-Danlos syndromes (collagen synthesis defects), and systemic autoimmune connective tissue diseases (lupus, scleroderma).

\medterm{Contralateral}-- A relational anatomical term denoting a structure, lesion, or symptom located on or affecting the opposite side of the body from another reference structure. Opposed to ipsilateral. For example, because the corticospinal tracts decussate in the medullary pyramids, a stroke in the left motor cortex results in contralateral (right-sided) hemiparesis.

\medterm{Coracoid Process}-- Hook-like scapular projection extending anteriorly and laterally. Attachment for 3 muscles (pectoralis minor, coracobrachialis, short head of biceps) and 3 ligaments. Important surgical landmark. Fractures uncommon but may occur with direct trauma or shoulder dislocations.

\medterm{Cornea}-- Transparent, avascular, dome-shaped anterior eye surface providing approximately two-thirds of total refractive power. Five layers: epithelium (stratified squamous, rapidly regenerating), Bowman's membrane, stroma (thickest, regular collagen), Descemet's membrane, endothelium (single cell layer maintaining transparency). Innervated by CN V1 (ophthalmic trigeminal), extremely sensitive. Conditions: abrasion, ulcers, keratoconus (progressive thinning and cone-shaped protrusion). Contact lens wear is a major risk factor for microbial keratitis, which can rapidly lead to corneal ulceration, perforation, and sight loss if not treated promptly.

\medterm{Coronal Plane}-- A vertical anatomical plane (also known as the frontal plane) oriented parallel to the coronal suture of the skull, passing vertically from side to side to divide the body, organ, or specimen into anterior (front) and posterior (back) portions. Movements taking place in the coronal plane include abduction and adduction. In radiological imaging (CT and MRI), coronal sections display structures viewed from the front, essential for evaluating the paranasal sinuses, orbits, kidneys, and musculoskeletal joints.

\medterm{Coronary Arteries}-- The arterial system supplying oxygenated blood directly to the myocardium, arising from the ascending aorta immediately superior to the aortic valve cusps. The right coronary artery (RCA) arises from the right (anterior) aortic sinus and courses in the coronary sulcus, giving off the conus artery, sinus node artery (in 60 percent of individuals), right marginal artery, and, in right-dominant circulations (approximately 85 percent), the posterior descending artery (PDA/posterior interventricular artery) and AV nodal artery (in 90 percent), supplying the right atrium, right ventricle, SA/AV nodes, and inferior wall of the left ventricle. The left coronary artery (LCA/left main) arises from the left (posterior) aortic sinus and quickly bifurcates into the left anterior descending artery (LAD/anterior interventricular, running in the anterior interventricular groove to supply the anterior wall of the left ventricle, anterior two-thirds of the interventricular septum, and apex) and the left circumflex artery (LCx, coursing in the left coronary sulcus to supply the lateral and posterior left ventricular walls). Unlike most systemic beds, myocardial perfusion occurs almost entirely during ventricular diastole. Atherosclerotic occlusion causes coronary artery disease, angina pectoris, and acute myocardial infarction (LAD occlusion is known as the 'widow-maker').

\medterm{Coronary Sinus}-- A large vein on the back of the heart that collects deoxygenated blood from the heart muscle and drains it into the right atrium.

\medterm{Corpus Callosum}-- Brain's largest white matter tract with approximately 200 million axons connecting left and right hemispheres. Four parts: rostrum, genu, body, splenium. Facilitates interhemispheric communication for sensory, motor, and cognitive integration. Genu connects prefrontal cortices; body connects motor/sensory cortices; splenium connects temporal/occipital cortices. Callosotomy is a last-resort epilepsy treatment producing split-brain syndromes with disconnection phenomena.

\medterm{Cranial Cavity}-- The hollow space inside the skull that contains and protects the brain.

\medterm{Cranial Nerve}-- Twelve paired nerves (CN I-XII) emerging directly from the brain/brainstem, primarily serving head and neck. I Olfactory (smell), II Optic (vision), III Oculomotor (eye movement, pupil), IV Trochlear (superior oblique), V Trigeminal (face sensation, mastication), VI Abducens (lateral rectus), VII Facial (expression, taste anterior 2/3), VIII Vestibulocochlear (hearing, balance), IX Glossopharyngeal (taste posterior 1/3, swallowing), X Vagus (parasympathetic to thorax/abdomen), XI Accessory (SCM, trapezius), XII Hypoglossal (tongue movement). Mnemonics: "Oh, Oh, Oh, To Touch And Feel Very Good Velvet, Ah, Heaven."

\medterm{Cuboid}-- Cube-shaped lateral foot tarsal bone articulating with calcaneus proximally and 4th/5th metatarsals distally. Attachment for the long plantar ligament and peroneus longus tendon. Keystone of the lateral longitudinal arch. Cuboid syndrome (subluxation) is a common lateral foot pain cause in athletes from excessive pronation.

\medterm{Deep}-- A directional term indicating a location situated further inward, away from the surface of the body, or deeper within an organ. Opposed to superficial. Examples include the internal carotid artery running deep to the sternocleidomastoid muscle, and bone marrow lying deep within cortical bone.

\medterm{Deltoid Muscle}-- Large triangular muscle forming the shoulder contour, from lateral clavicle, acromion, and scapular spine to the deltoid tuberosity. Innervated by the axillary nerve (C5-C6). Anterior fibers flex/medially rotate, middle fibers abduct (primary abductor beyond 15 degrees), posterior fibers extend/laterally rotate. Essential for abduction; common injection site. Axillary nerve injury causes loss of abduction.

\medterm{Dendrite}-- The branched, tapering neuronal processes that receive synaptic input from other neurons and conduct graded potentials toward the cell body. Dendrites contain receptors, ion channels, and postsynaptic densities. Dendritic spines (small protrusions) are major sites of excitatory (glutamatergic) synapse formation; their density and morphology change with learning and are altered in many neuropsychiatric and neurodegenerative conditions (e.g., Down syndrome, schizophrenia). Electrical signals are graded and summate at the axon hillock, where an action potential is generated if threshold is reached.

\medterm{Dens}-- A prominent tooth-like process that projects superiorly from the anterior body of the axis (C2 vertebra), serving as the pivotal axis around which the atlas (C1) rotates. See \hyperlink{Odontoid Process}{Odontoid Process}.

\medterm{Depression}-- An inferior movement of a body part, moving the structure downward away from the head. Examples include lowering the shoulders from a shrugged position (depression of the scapulae by the pectoralis minor, lower trapezius, and gravity) and opening the mouth (depression of the mandible by the lateral pterygoid, digastric, and infrahyoid muscles). Opposed to elevation.

\medterm{Dermis}-- Thick connective tissue layer beneath the epidermis providing structural support, elasticity, and nourishment. Two layers: papillary (superficial, loose connective tissue with dermal papillae interdigitating with epidermal ridges) and reticular (deep, dense irregular connective tissue). Contains collagen (type I, tensile strength), elastin, fibroblasts, blood vessels, lymphatics, nerves, hair follicles, sebaceous glands, and sweat glands. The dermal-epidermal junction is critical for adhesion; disruption causes blistering disorders such as pemphigus vulgaris and toxic epidermal necrolysis.

\medterm{Descending Colon}-- From the splenic flexure to the left iliac fossa, approximately 25-30 cm, and mostly retroperitoneal and fixed to the posterior abdominal wall. The phrenicocolic ligament supports the splenic flexure. Anteriorly it is related to small intestine; posteriorly to the left kidney. It absorbs water and electrolytes and compacts fecal matter. Left-sided pathology, including diverticulitis and cancer, is common here and in the sigmoid colon.

\medterm{Descending Thoracic Aorta}-- From aortic arch at T4 to diaphragm at T12. Gives off bronchial, esophageal, mediastinal, and posterior intercostal arteries. In the posterior mediastinum, initially left of the vertebral column. Related to esophagus, thoracic duct, and azygos vein. Aneurysms are mostly atherosclerotic, treatable with thoracic endovascular aortic repair.

\medterm{Diaphragm}-- Dome-shaped musculotendinous partition separating thoracic/abdominal cavities, primary inspiration muscle. Central tendon with muscular periphery from xiphoid, lower six costal cartilages, and lumbar vertebrae via crura and arcuate ligaments. Innervated by phrenic nerves (C3-C5). Contraction flattens the dome increasing thoracic volume. Three openings: caval (T8), esophageal (T10), aortic (T12). Paralysis causes hemiparesis or paradoxical breathing.

\medterm{Diaphysis}-- The shaft of a long bone, comprising thick cortical (compact) bone surrounding the medullary cavity, which contains bone marrow in children and is largely yellow (fatty) marrow in adults. Enclosed by periosteum; the nutrient artery enters here. Fractures of the diaphysis heal by callus formation (enchondral and intramembranous). Blood supply via nutrient artery (avulsed in displaced shaft fractures, delaying healing). Contrast with epiphysis (ends) and metaphysis (growth zone).

\medterm{Distal}-- A directional term used primarily in the extremities to describe a position situated further away from the point of attachment of a limb to the trunk, or further along the flow/course of a vessel, nerve, or tubular organ. Opposed to proximal. Examples include the fingers being distal to the elbow, the foot being distal to the knee, and the distal colon leading to the rectum.

\medterm{Dorsal}-- A directional term relating to or situated toward the back surface of the trunk, or the superior surface of the foot (dorsum of foot) and posterior surface of the hand (dorsum of hand). In human bipedal anatomy, dorsal is largely synonymous with posterior. In spinal cord anatomy, the dorsal roots carry primary sensory (afferent) fibers whose cell bodies reside in the dorsal root ganglia. Opposed to ventral.

\medterm{Dorsal Root Ganglion}-- A swelling on the dorsal root of a spinal nerve containing the cell bodies of primary sensory neurons. Its pseudounipolar neurons convey peripheral sensory information to the spinal cord.

\medterm{Duodenum}-- Shortest first small intestinal segment (25 cm), C-shaped around the pancreatic head. Four parts. Descending part has the major duodenal papilla for bile/pancreatic duct entry via the hepatopancreatic ampulla. Brunner's glands produce alkaline mucus. Most common peptic ulcer site. Primary iron and calcium absorption site.

\medterm{Dural Venous Sinuses}-- Venous channels between dural layers, draining brain blood and CSF from the subarachnoid space into the internal jugular veins. Major sinuses: superior sagittal, transverse, sigmoid, cavernous, straight. Valveless with bidirectional flow possible. Thrombosis (cerebral venous sinus thrombosis) causes increased intracranial pressure, headache, seizures, and focal deficits.

\medterm{Dura Mater}-- The outermost, thickest meningeal layer, a tough fibrous membrane. In the cranium it is two layers: periosteal (outer, adherent to skull) and meningeal (inner), which separate to form dural venous sinuses (e.g., superior sagittal, cavernous). Dural folds: falx cerebri, tentorium cerebelli, falx cerebelli. Cranial dura is pain-sensitive near vessels and at the base. Epidural hematoma (arterial, middle meningeal, lens-shaped) lies outside the dura; subdural hematoma (venous, crescent-shaped) lies between dura and arachnoid.

\medterm{Ear}-- Organ of hearing and balance, divided into external, middle, and inner ear. External: auricle (pinna), external auditory canal, tympanic membrane. Middle: air-filled cavity in temporal bone, auditory ossicles (malleus, incus, stapes), Eustachian tube (to nasopharynx). Inner (bony labyrinth): cochlea (hearing), vestibule (utricle, saccule -- linear acceleration), semicircular canals (rotational acceleration). CN VIII (vestibulocochlear). Conditions: otitis media/externa, hearing loss, vertigo, Meniere's disease, acoustic neuroma, cholesteatoma.

\medterm{Eccrine Gland}-- A coiled sweat gland distributed throughout the skin that secretes a watery fluid directly onto the skin surface through a duct. Eccrine sweating supports thermoregulation and is stimulated by sympathetic cholinergic fibres; the glands are especially numerous on the palms, soles, and forehead.

\medterm{Elastic Cartilage}-- Flexible cartilage containing abundant elastic fibers in its matrix. It supports structures that require shape retention with flexibility, including the external ear, epiglottis, and auditory tube.

\medterm{Elbow Joint}-- A synovial hinge complex comprising the humeroulnar and humeroradial joints. It primarily permits flexion and extension, while the proximal radioulnar component permits forearm rotation.

\medterm{Elevation}-- A superior movement of a body part in the frontal or sagittal plane, moving the structure upward toward the head. Examples include shrugging the shoulders (elevation of the scapulae and shoulder girdle by the upper trapezius and levator scapulae) and closing the jaw (elevation of the mandible by the masseter, temporalis, and medial pterygoid muscles). Opposed to depression.

\medterm{Endocardium}-- The smooth, innermost endothelial and subendothelial connective-tissue layer lining the internal cavities of the heart, continuous with the tunica intima of the great blood vessels. Histologically, it consists of a single layer of simple squamous endothelial cells overlying a thin subendothelial layer of loose connective tissue, elastic fibers, and a deep subendocardial layer containing collagen, nerves, and the specialized conduction fibers of the Purkinje system. The cardiac valves are duplications or folds of endocardium reinforced by a central core of dense fibrous connective tissue. Clinical significance: Infective endocarditis is a microbial infection of the endocardial surface, primarily targeting cardiac valves damaged by rheumatic heart disease or congenital anomalies, forming vegetations that cause valvular insufficiency, septic emboli, Osler nodes, and Janeway lesions.

\medterm{Endosteum}-- The thin connective tissue membrane lining the inner surface of bone, including the medullary cavity, Haversian canals, and trabeculae. Contains osteoprogenitor cells and osteoblasts, contributing to bone formation, remodeling, and repair (fracture healing). Participates with periosteum in the regulation of bone growth and mineral homeostasis.

\medterm{Epidermis}-- Outermost skin layer, stratified squamous keratinized epithelium approximately 0.1 mm thick, with five layers (superficial to deep): stratum corneum (dead keratinized cells), lucidum (clear layer, thick skin only), granulosum (keratinization begins), spinosum (spiny layer, Langerhans cells), basale (stem cells, melanocytes, Merkel cells). Keratinocytes (80 percent of cells) produce keratin for waterproofing and protection. Avascular, nourished by dermal diffusion. Melanocytes produce melanin for UV protection and skin pigmentation.

\medterm{Epididymis}-- Coiled duct approximately 6 meters long on the posterolateral aspect of each testis, divided into head (caput, receiving efferent tubules from the testis), body (corpus), and tail (cauda, continuing as the vas deferens). Functions include sperm maturation (gaining motility over approximately 12 days), storage (in the tail), and transport. The epididymis is a common site of infection (epididymitis), typically from sexually transmitted organisms in young men or urinary tract organisms in older men.

\medterm{Epigastric Region}-- Central superior region between the hypochondriac regions and above the umbilical region. Overlies the stomach, duodenum, pancreas, and portions of the liver. The aorta and inferior vena cava pass through it. Epigastric pain associates with gastric ulcers, pancreatitis, GERD, and aortic aneurysms. Critical for detecting pulsatile masses and organomegaly.

\medterm{Epiglottis}-- A leaf-shaped elastic cartilage at the base of the tongue that covers the laryngeal inlet during swallowing, directing food and liquid into the esophagus and preventing aspiration into the airway. Elevates during swallowing (with the larynx) under glossopharyngeal/vagal innervation. Acute epiglottitis (traditionally Haemophilus influenzae type b, now often viral or other bacteria) causes severe sore throat, drooling, stridor, and airway obstruction—a medical emergency. Inflamed epiglottis is visible on lateral neck X-ray ('thumb sign').

\medterm{Epiphyseal Plate}-- A layer of hyaline cartilage, also called the growth plate or physis, located between the epiphysis and metaphysis of a growing long bone. It permits longitudinal bone growth through endochondral ossification, progressing from resting cartilage through proliferation and hypertrophy to calcification and replacement by bone. Injury can cause premature closure, limb-length discrepancy, or angular deformity; after puberty it closes and becomes the epiphyseal line.

\medterm{Epiphysis}-- The expanded end of a long bone, usually covered by articular cartilage where it forms a synovial joint. It contains spongy (trabecular) bone with a thin outer shell of compact bone and contributes to the shape and function of the joint. In growing children, the epiphysis is separated from the metaphysis by the epiphyseal plate; after growth ends, the plate ossifies and becomes the epiphyseal line.

\medterm{Epithelial Tissue}-- One of the four primary tissue types, composed of closely aggregated cells with minimal extracellular matrix, adhering tightly via junctional complexes (tight junctions, adherens junctions, desmosomes) and anchored to an underlying basement membrane. It covers all external body surfaces, lines internal body cavities, tubes, and vessels, and forms the secretory units of all exocrine and endocrine glands. It is completely avascular, receiving oxygen and nutrients via diffusion from adjacent vascularized lamina propria. Classified by cell layers into simple (single layer: squamous in alveoli/endothelium, cuboidal in renal tubules, columnar in stomach/intestine) and stratified (multiple layers: squamous keratinized in epidermis, non-keratinized in oral/esophageal mucosa, transitional/urothelium in urinary tract), plus pseudostratified ciliated columnar epithelium in the respiratory tract. Functions include physical protection, selective permeability, absorption, transcellular transport, secretion, and sensory reception. More than 80 percent of human cancers originate from epithelial cells (carcinomas).

\medterm{Epithelium}-- One of the four primary tissue types, composed of cells that line the surfaces of the body (skin, mucous membranes, serous membranes) and form glands. Epithelial tissues are classified by cell shape (squamous--flat, cuboidal--cube-shaped, columnar--tall, transitional--variable) and layering (simple--single layer, stratified--multiple layers, pseudostratified--appears layered but all cells contact the basement membrane). Functions: protection (stratified squamous epithelium of skin), absorption (simple columnar of intestinal lining), secretion (glandular epithelium), filtration (simple squamous of capillary endothelium), and sensory reception (olfactory epithelium). Epithelium rests on a basement membrane, is avascular (receives nutrition by diffusion), and has high regenerative capacity. Epithelial-mesenchymal transition is a key process in development and cancer metastasis.

\medterm{Esophagus}-- A muscular tubular conduit approximately 25 cm long that transports food and liquids from the pharynx to the stomach. It originates at the lower border of the cricoid cartilage (C6 vertebral level), traverses the superior and posterior mediastinum, passes through the diaphragmatic hiatus at the T10 level, and terminates at the cardiac orifice of the stomach (T11 level). It exhibits three physiological constrictions where foreign bodies or carcinomas frequently lodge: cervical (cricopharyngeus at C6), bronchoaortic (crossing of aortic arch and left main bronchus at T4--T5), and diaphragmatic (hiatus at T10). Its wall is unique in having an inner circular and outer longitudinal muscularis that transitions from skeletal muscle in the upper third, to mixed in the middle third, and smooth muscle in the lower third. Lined by non-keratinized stratified squamous epithelium, it transitions abruptly to columnar gastric mucosa at the squamocolumnar junction (Z-line). Blood supply arises from the inferior thyroid, bronchial, esophageal (from descending aorta), and left gastric arteries. Venous drainage from the lower esophagus forms a key portosystemic anastomosis between the left gastric vein (portal) and azygos veins (systemic), dilation of which produces life-threatening esophageal varices in portal hypertension. Pathologies include gastroesophageal reflux disease (GERD), Barrett esophagus (premalignant intestinal metaplasia), achalasia, and esophageal adenocarcinoma or squamous cell carcinoma.

\medterm{Ethmoid Bone}-- Delicate bone between the orbits forming parts of the nasal septum and orbital/nasal roofs. Has the cribriform plate (olfactory nerve filaments), perpendicular plate (superior nasal septum), and labyrinths (nasal conchae, ethmoid air cells). The crista galli provides falx cerebri attachment. Cribriform fractures cause CSF rhinorrhea and anosmia.

\medterm{Eustachian Tube}-- Auditory/pharyngotympanic tube connecting the middle ear to the nasopharynx, approximately 36 mm long in adults. Cartilaginous medial two-thirds and bony lateral third. Functions: equalizing middle ear/atmospheric pressure, draining middle ear secretions, and protecting from nasopharyngeal pathogens. Opens during swallowing, yawning, sneezing via tensor/levator veli palatini muscles. Children have shorter, wider, more horizontal tubes predisposing to otitis media. Dysfunction causes middle ear effusion (glue ear), the most common cause of conductive hearing loss in children.

\medterm{Eversion}-- A specialized movement of the foot occurring at the subtalar and transverse tarsal (talocalcaneonavicular and calcaneocuboid) joints, in which the sole of the foot is turned laterally away from the median plane, tilting the lateral border of the foot upward. It is primarily produced by the fibularis (peroneus) longus and fibularis brevis muscles (innervated by the superficial fibular nerve). Excessive eversion under stress can tear the deltoid ligament or cause a bimalleolar/trimalleolar Potts fracture.

\medterm{Extension}-- A straightening movement, typically occurring in the sagittal plane, that increases the angle between two articulating bones or body parts, often restoring a joint to anatomical position from a flexed state. Movement beyond the normal anatomical limit is termed hyperextension. Produced by extensor muscle groups (e.g., triceps brachii at the elbow; quadriceps femoris at the knee; gluteus maximus and hamstrings at the hip; erector spinae at the vertebral column).

\medterm{Extensor Digitorum}-- Posterior forearm muscle from lateral epicondyle via the common extensor tendon, dividing into 4 tendons inserting on extensor expansions of digits 2-5. Innervated by posterior interosseous nerve (C7-C8). Extends MCP joints, assists PIP/DIP extension. Lateral epicondylitis (tennis elbow) is an overuse injury of its common origin. Extensor hood coordinates finger movements.

\medterm{External Auditory Canal}-- Tube from the concha to the tympanic membrane, approximately 2.5 cm long in adults. Lateral third is cartilaginous; medial two-thirds are bony (temporal bone). Skin-lined with ceruminous glands (producing cerumen/earwax) and hair follicles. Cerumen has antimicrobial and insect-repellent properties. Innervated by auriculotemporal nerve (V3) laterally and vagus nerve (Arnold's nerve) medially, explaining why canal touching may trigger coughing. Otitis externa (swimmer's ear) is common.

\medterm{External Iliac Artery}-- Larger common iliac terminal branch, continuing as the femoral artery below the inguinal ligament. Gives off inferior epigastric and deep circumflex iliac arteries. Related to external iliac vein medially and ureter anteriorly. Commonly used for cardiac catheterization and endovascular access. Atherosclerosis causes thigh/calf claudication.

\medterm{External Jugular Vein}-- Formed by posterior auricular and retromandibular vein divisions, crossing the sternocleidomastoid superficially to drain into the subclavian vein. Visible through the skin; used for peripheral IV access when other sites are unavailable. Distended in right heart failure and SVC obstruction. A useful physical examination landmark.

\medterm{External Oblique Muscle}-- Most superficial flat abdominal muscle from outer lower eight ribs to linea alba, pubic tubercle, and iliac crest. Innervated by T7-T12. Contralateral trunk rotation and flexion; compresses abdomen. Forms the inguinal ligament from its free inferior border. Superficial inguinal ring is an opening in its aponeurosis. Strains common in twisting athletes.

\medterm{Extracellular Matrix}-- A complex, dynamic, non-cellular three-dimensional macromolecular network secreted by local cells that provides essential physical scaffolding, mechanical resilience, and biochemical signaling to surrounding tissues. It comprises two main classes of macromolecules: (1) glycosaminoglycans (GAGs, such as hyaluronic acid, chondroitin sulfate, keratan sulfate) covalently linked to proteins forming proteoglycans, which form a highly hydrated, gel-like ground substance that resists compressive forces and facilitates nutrient diffusion; and (2) fibrous structural proteins, primarily collagens (type I in bone/tendon/dermis, type II in hyaline cartilage, type III in reticular meshworks, type IV in basement membranes) providing tensile strength, and elastin providing elastic recoil. Adhesive glycoproteins (fibronectin, laminin) link cells to the ECM via transmembrane integrin receptors, transducing mechanical signals to the cytoskeleton (mechanotransduction) and regulating cell survival, migration, proliferation, and tissue morphogenesis. Abnormal ECM remodeling occurs in hepatic cirrhosis, pulmonary fibrosis, systemic sclerosis, and cancer metastasis.

\medterm{Extraocular Muscles}-- Six muscles controlling eye movement: four rectus (superior, inferior, medial, lateral) and two oblique (superior, inferior). Medial rectus (CN III) adducts; lateral rectus (CN VI) abducts; superior rectus (CN III) elevates/intorts; inferior rectus (CN III) depresses/extorts; superior oblique (CN IV) depresses/intorts; inferior oblique (CN III) elevates/extorts. Hering's law: yoke muscles receive equal innervation. Sherrington's law: agonist contraction accompanies antagonist relaxation. Paralysis of individual muscles causes characteristic strabismus (squint) and diplopia.

\medterm{Eye}-- The organ of vision, a spherical structure (approximately 24 mm diameter) in the bony orbit. Three layers: fibrous (sclera, cornea), vascular/uveal (choroid, ciliary body, iris), neural (retina). Anterior segment: cornea, anterior chamber (aqueous humor), iris/pupil, lens. Posterior segment: vitreous humor, retina (photoreceptors: rods/cones), optic nerve. Extraocular muscles (6) control movement. Cranial nerves II, III, IV, VI. Conditions: refractive errors, cataract, glaucoma, macular degeneration, retinal detachment, diabetic retinopathy.

\medterm{Facet Joints}-- Zygapophysial joints connecting adjacent vertebral articular processes. Orientation varies: cervical (45 degrees, allowing rotation), thoracic (coronal, allowing rotation but limiting flexion/extension), lumbar (sagittal, allowing flexion/extension but limiting rotation). Innervated by medial branches of dorsal rami. Common source of low back pain treated with facet injections and medial branch blocks.

\medterm{Facial Nerve (Cranial Nerve VII)}-- The seventh cranial nerve (CN VII), a complex mixed nerve containing branchiomotor, parasympathetic secretomotor, special sensory (taste), and somatic sensory components. It emerges from the pontomedullary junction at the cerebellopontine angle, enters the internal acoustic meatus alongside the vestibulocochlear nerve (CN VIII), travels through the facial canal in the temporal bone (giving off the greater petrosal nerve for lacrimation, nerve to stapedius, and chorda tympani for taste to anterior 2/3 of tongue and submandibular/sublingual salivation), and exits the skull base via the stylomastoid foramen. It traverses the parotid gland (without innervating it) and divides into five terminal motor branches: Temporal, Zygomatic, Buccal, Marginal mandibular, and Cervical (mnemonic: 'To Zanzibar By Motor Car'), supplying all muscles of facial expression, the posterior belly of the digastric, stylohyoid, and platysma. Lower motor neuron lesions (Bell's palsy) cause complete ipsilateral facial paralysis (including forehead drop and inability to close the eye); upper motor neuron lesions (stroke) spare the forehead due to bilateral cortical innervation of the frontalis muscle.

\medterm{Fallopian Tube}-- The paired muscular tubes (10-12 cm) connecting the ovaries to the uterus, conducting oocytes and sperm and the site of fertilization (usually ampulla). Segments: infundibulum (fimbriae—sweep the oocyte), ampulla, isthmus, intramural. Lined by ciliated columnar epithelium and secretory cells; peristalsis and ciliary action transport the oocyte. Blockage causes infertility; inflammation (salpingitis, from ascending pelvic infection) predisposes to ectopic pregnancy. Tubal ligation is a sterilization procedure.

\medterm{False Ribs}-- Ribs 8-10 not articulating directly with the sternum; their cartilages join rib 7's cartilage forming the costal margin. Greater mobility and flexibility of the lower thoracic wall. More susceptible to fractures at their angles due to greater length and less rigid attachment. The costal margin is an important clinical landmark for abdominal examination.

\medterm{Femoral Artery}-- Continues from external iliac through the femoral triangle and adductor canal to become the popliteal artery. Branches: superficial/deep circumflex femoral, deep femoral (profunda femoris), muscular. Palpable in the femoral triangle; most common site for arterial catheterization. Related to femoral vein medially and nerve laterally. Atherosclerosis causes peripheral arterial disease.

\medterm{Femoral Nerve}-- Largest lumbar plexus nerve (L2-L4), passing through psoas major and entering the thigh under the inguinal ligament lateral to the femoral artery. Innervates quadriceps, sartorius, pectineus; sensation to anterior thigh and medial leg (saphenous nerve). Used for knee surgery nerve block anesthesia. Neuropathy (retroperitoneal hemorrhage, hip surgery) causes quadriceps weakness and loss of knee jerk.

\medterm{Femoral Triangle}-- Upper anterior thigh space bounded superiorly by the inguinal ligament, laterally by sartorius, and medially by adductor longus. Contents from lateral to medial: femoral nerve, artery, vein, and deep inguinal lymph nodes. The femoral artery is used for arterial catheterization. The femoral vein is accessed for central venous lines. Femoral hernias occur through the femoral canal medial to the vein.

\medterm{Femoral Vein}-- The major deep vein of the thigh, continuation of the popliteal vein (through the adductor hiatus), running with the femoral artery in the femoral sheath (medial to artery) through the femoral triangle, becoming the external iliac vein at the inguinal ligament. Receives the great saphenous vein. Site of central venous access (femoral line), femoral venipuncture, and DVT. The femoral sheath also contains the femoral canal (medial), site of femoral hernia.

\medterm{Femur}-- Longest, heaviest, strongest body bone. Proximal: head, neck, trochanters. Shaft: linea aspera. Distal: condyles, epicondyles. The neck connects head to shaft at approximately 125 degrees. Neck fractures are common in elderly osteoporotic patients with avascular necrosis risk. Bears entire upper body weight during standing.

\medterm{Fibrocartilage}-- Strong cartilage containing dense bundles of type I collagen, adapted to resist compression and tension. It forms intervertebral discs, the pubic symphysis, and menisci.

\medterm{Fibula}-- Slender lateral leg bone, not weight-bearing, serving for muscle attachment and forming the lateral malleolus. Head provides biceps femoris and lateral collateral ligament attachment; the common peroneal nerve wraps around its neck, vulnerable to injury. Lateral malleolus extends more distally and posteriorly than the medial malleolus. Often fractured with tibial fractures.

\medterm{Flexion}-- A movement, generally taking place in the sagittal plane, that decreases the angle between two articulating bones or body parts, usually bending a joint anteriorly (except at the knee, where flexion bends the joint posteriorly). In the foot, upward movement toward the shin is termed dorsiflexion, whereas downward movement is plantarflexion. Produced by flexor muscle groups (e.g., biceps brachii and brachialis at the elbow; iliopsoas and rectus femoris at the hip; hamstrings at the knee).

\medterm{Flexor Digitorum Profundus}-- Deep anterior forearm muscle with ulnar and humeroulnar heads, inserting on distal phalanges bases of digits 2-5. Medial half (digits 4-5) innervated by ulnar nerve; lateral half (digits 2-3) by anterior interosseous nerve (C8-T1). Flexes DIP, PIP, MCP joints, and wrist. Anterior interosseous palsy causes inability to pinch with the index finger.

\medterm{Flexor Digitorum Superficialis}-- Large anterior forearm muscle with humeroulnar, radial, and ulnar heads, inserting on middle phalanges of digits 2-5. Innervated by median nerve (C8-T1). Flexes PIP joints, MCP joints, and wrist. Important for median nerve function assessment. Tendons pass through carpal tunnel, susceptible to stenosing tenosynovitis (trigger finger).

\medterm{Floating Ribs}-- Ribs 11-12 with short costal cartilages terminating in abdominal wall musculature, not articulating with the sternum or other ribs. Most mobile and least protected, susceptible to fracture from direct trauma. Fractures may injure underlying kidneys, requiring hematuria and flank tenderness assessment in trauma.

\medterm{Fovea Centralis}-- Central macular pit (approximately 1.5 mm) with the highest cone density (approximately 200,000 per square millimeter) and no rods, providing the sharpest visual acuity and color vision. Avascular with overlying retinal layers displaced laterally for direct light access to cones. The foveola (central 0.35 mm) has the thinnest retina and highest cone density. Pathology here (macular holes, central serous chorioretinopathy, macular degeneration) severely impairs central vision.

\medterm{Frontal Bone}-- Forms the forehead and superior orbits, with the squama frontalis, orbital plates, and nasal part. Contains frontal sinuses draining into the nasal cavity via the frontonasal duct. Articulates with parietal bones at the coronal suture. The metopic suture fuses by age six. Fractures may cause pneumocephalus or CSF rhinorrhea through the sinuses.

\medterm{Frontal Lobe}-- The largest lobe of the cerebral cortex, lying anterior to the central sulcus, comprising the primary motor cortex (precentral gyrus, Brodmann 4), premotor and supplementary motor areas, prefrontal cortex, and Broca's area (dominant hemisphere, expressive language). Functions: motor planning and execution, executive function, judgment, reasoning, working memory, social behavior, and personality. Lesions cause contralateral hemiparesis (UMN), expressive aphasia (Broca), personality changes, apathy, and executive dysfunction; prefrontal injury (e.g., Phineas Gage) alters behavior. Orbital surface (olfaction).

\medterm{Gallbladder}-- A pear-shaped, hollow musculo-membranous reservoir (7--10 cm long, capacity 30--50 mL) located in a shallow fossa on the visceral surface of the right hepatic lobe, at the junction of segments IV and V. It stores, concentrates (up to 5--10 fold via active sodium and water reabsorption), and releases bile into the duodenum in response to cholecystokinin (CCK) released by duodenal I-cells following a fatty meal. It is anatomically divided into four regions: the fundus (projecting past the inferior liver border at the intersection of the right midclavicular line and the 9th costal cartilage), body, infundibulum (which may form an outpouching known as Hartmann's pouch, a common site of stone impaction), and neck, which tapers into the cystic duct containing the spiral folds of Heister. Arterial supply is provided by the cystic artery, typically arising from the right hepatic artery within the hepatobiliary triangle (triangle of Calot, bounded by the cystic duct, common hepatic duct, and inferior liver surface). Venous drainage passes directly into liver sinusoids or via cystic veins into the portal vein. Clinical conditions include cholelithiasis (gallstones), acute cholecystitis (characterized by right upper quadrant pain, fever, leukocytosis, and an inspiratory arrest on palpation known as Murphy's sign), choledocholithiasis, gallstone pancreatitis, biliary colic, and gallbladder carcinoma.

\medterm{Ganglion}-- A collection of neuronal cell bodies located outside the central nervous system (CNS). Ganglia are classified as: (1) Sensory ganglia (dorsal root ganglia of spinal nerves, cranial nerve ganglia--trigeminal, geniculate, vestibular, spiral, petrosal, nodose), containing pseudounipolar neurons whose peripheral processes receive sensory input and central processes enter the CNS. (2) Autonomic ganglia (sympathetic chain ganglia, prevertebral ganglia--celiac, superior mesenteric, inferior mesenteric, and parasympathetic ganglia--ciliary, pterygopalatine, submandibular, otic, enteric--Auerbach and Meissner plexuses), containing multipolar neurons with efferent output to target organs. (3) Basal ganglia (CNS structures involved in motor control: caudate, putamen, globus pallidus, substantia nigra, subthalamic nucleus--non-pathological usage differs from peripheral ganglia). Clinically: dorsal root ganglia are sites of latency for herpes zoster; basal ganglia dysfunction causes Parkinson disease, Huntington disease, and dystonia. The term "ganglion" is also used for benign cystic swellings arising from tendon sheaths or joint capsules (ganglion cyst), which contain mucinous fluid and are not neural in origin.

\medterm{Gastrocnemius Muscle}-- Most superficial posterior leg muscle forming the calf prominence, with medial/lateral femoral condyle heads merging with soleus to form the Achilles tendon on the calcaneus. Innervated by tibial nerve (S1-S2). Plantarflexes ankle and flexes knee (two-joint muscle). Active in walking/running/jumping. Achilles tendon is the body's strongest but vulnerable to rupture during explosive activities.

\medterm{Gastrointestinal Villi}-- Finger-like projections of the small-intestinal mucosa that greatly increase absorptive surface area. Each villus contains blood capillaries and a central lacteal for transport of absorbed nutrients and lipids.

\medterm{Genitalia}-- The external and internal reproductive organs. The term includes the vulva and vagina, or the penis, scrotum, and associated structures, together with the gonads and reproductive ducts when used broadly; anatomy varies with sex, age, and developmental stage.

\medterm{Gland}-- An organized group of epithelial cells specialized for secretion. Glands are classified by: (1) Location: exocrine (secrete onto a surface via ducts--sweat, salivary, mammary, sebaceous, lacrimal, pancreatic exocrine) or endocrine (secrete hormones directly into the bloodstream--thyroid, pituitary, adrenal, pancreatic islets, parathyroid, gonads). (2) Number of cells: unicellular (goblet cells) or multicellular. (3) Duct structure: simple (unbranched duct) or compound (branched duct). (4) Secretory unit shape: tubular, acinar (alveolar), or tubuloacinar. (5) Secretion mechanism: merocrine (exocytosis, e.g., salivary, pancreatic), apocrine (apical cell membrane pinches off, e.g., mammary, axillary sweat), holocrine (entire cell disintegrates, e.g., sebaceous). (6) Secretion type: serous (watery, enzyme-rich), mucous (viscous, mucin-rich), mixed (seromucous).

\medterm{Glenohumeral Joint}-- Ball-and-socket synovial joint between humeral head and glenoid cavity. Most mobile body joint allowing flexion, extension, abduction, adduction, medial/lateral rotation, circumduction. Mobility trades for stability, making it the most commonly dislocated joint. Stabilized by glenoid labrum, glenohemoral ligaments, and rotator cuff. Loose, redundant capsule.

\medterm{Glenoid Cavity}-- Shallow pear-shaped depression on the lateral scapular angle articulating with the humeral head. The glenoid labrum deepens the socket by approximately 50 percent. Due to its shallowness, the glenohumeral joint is the most commonly dislocated joint. Bankart lesions (inferior labral tears) are associated with anterior dislocations and recurrent instability.

\medterm{Glisson's Capsule}-- Thin fibrous connective tissue envelope surrounding the liver, sending septa inward dividing it into lobules. Richly innervated by phrenic and intercostal nerve pain fibers, sensitive to distension and inflammation. Covered by visceral peritoneum except at the bare area. Hepatic injuries may cause subcapsular hematomas that can rupture causing hemoperitoneum.

\medterm{Globin}-- A family of heme-binding proteins involved in oxygen transport and storage. Structure: globin fold (8 alpha-helices arranged in a globular structure), binds one heme molecule (iron-protoporphyrin IX) per subunit. Types: (1) hemoglobin (alpha2-beta2 tetramer in adults, HbA; alpha2-gamma2 in fetal, HbF): transports oxygen in red blood cells. (2) Myoglobin (single subunit): stores oxygen in skeletal and cardiac muscle, facilitates oxygen diffusion. (3) Neuroglobin: expressed in neurons, protects against hypoxia. (4) Cytoglobin: expressed in fibroblasts and other tissues, involved in nitric oxide metabolism and oxidative stress. Mutations in globin genes cause haemoglobinopathies: sickle cell disease (beta-globin Glu6Val), thalassaemias (reduced alpha or beta globin chain production).

\medterm{Globulin}-- A group of serum proteins (globular proteins) that are less soluble in water but soluble in dilute salt solutions. On serum protein electrophoresis, globulins migrate in the alpha-1, alpha-2, beta, and gamma regions. Alpha-1 globulins: alpha-1-antitrypsin, alpha-fetoprotein. Alpha-2 globulins: haptoglobin, alpha-2-macroglobulin, ceruloplasmin. Beta globulins: transferrin, beta-2-microglobulin, complement factors (C3, C4). Gamma globulins: immunoglobulins (antibodies). Abnormal globulin levels: polyclonal increase (chronic infection, autoimmune disease, cirrhosis), monoclonal gammopathy (multiple myeloma producing a monoclonal peak, MGUS, Waldenström macroglobulinaemia), decreased (immunodeficiency, nephrotic syndrome).

\medterm{Glomerulus}-- The tuft of capillaries in the renal corpuscle (Bowman's capsule) where blood is filtered to form urine. The glomerular filtration barrier consists of: (1) fenestrated endothelium (70-100 nm pores), (2) glomerular basement membrane (GBM, 300-350 nm, type IV collagen network), (3) podocyte foot processes (slit diaphragms formed by nephrin, NEPH1, podocin, and CD2AP). The barrier selectively retains cells and proteins >69 kDa while allowing water and small solutes to pass. Glomerular filtration rate (GFR): 125 mL/min (180 L/day). Mesangial cells support the capillary loops and regulate flow by contraction. Glomerular injury causes proteinuria and haematuria (glomerulonephritis).

\medterm{Glossopharyngeal Nerve (Cranial Nerve IX)}-- The ninth cranial nerve (CN IX), a mixed cranial nerve emerging from the post-olivary sulcus of the rostral medulla oblongata and exiting the skull through the jugular foramen. Functions: (1) Special visceral afferent: taste sensation from the posterior one-third of the tongue; (2) General visceral afferent: visceral sensation from the posterior tongue, palatine tonsils, pharynx, eustachian tube, tympanic cavity, and baroreceptors/chemoreceptors in the carotid sinus and carotid body; (3) Visceral motor (parasympathetic): preganglionic fibers via the lesser petrosal nerve synapsing in the otic ganglion to stimulate parotid gland secretion; (4) Somatic motor: branchiomotor fibers innervating the stylopharyngeus muscle (elevating pharynx during swallowing). It forms the afferent limb of the gag reflex (CN X is efferent). Pathology includes glossopharyngeal neuralgia, characterized by paroxysms of severe lancinating pain in the tonsillar fossa, pharynx, and ear, often triggered by swallowing.

\medterm{Glucagon}-- A peptide hormone produced by the alpha cells of the pancreatic islets (Langerhans). Primary function: counter-regulatory hormone to insulin, maintaining blood glucose by: (1) stimulating hepatic glycogenolysis (breakdown of glycogen to glucose), (2) promoting gluconeogenesis (synthesis of glucose from lactate, amino acids, glycerol), (3) inhibiting glycolysis and glycogen synthesis in the liver. Secreted in response to: hypoglycaemia, catecholamines, amino acids (arginine), and vagal stimulation. Secretion inhibited by: glucose, insulin, somatostatin. Glucagon also relaxes smooth muscle in the gastrointestinal tract (used therapeutically as IV glucagon for severe hypoglycaemia and for relaxation of the GI tract during endoscopy or imaging). Glucagonoma: a rare pancreatic neuroendocrine tumor producing glucagon, causing diabetes, weight loss, necrolytic migratory erythema (characteristic rash).

\medterm{Glucose}-- A monosaccharide (C6H12O6) that is the primary energy source for cells. Types: D-glucose (dextrose, naturally occurring enantiomer) and L-glucose (synthetic). Routes of entry: dietary carbohydrates (digested to glucose), gluconeogenesis (synthesis from non-carbohydrate precursors), glycogenolysis (breakdown of stored glycogen). Cellular metabolism: glycolysis (cytosol, 2 ATP per glucose), TCA cycle (mitochondria, 36 ATP total), pentose phosphate pathway (generates NADPH and ribose-5-phosphate). Glucose transport into cells: facilitated diffusion via GLUT transporters (GLUT1-4, 12 isoforms; GLUT4 is insulin-sensitive in muscle and adipose tissue). Blood glucose regulation: insulin (lowers), glucagon, cortisol, catecholamines, growth hormone (raise). Clinical significance: hyperglycaemia (diabetes), hypoglycaemia, glucose tolerance test (OGTT) for diabetes diagnosis. Intravenous dextrose is used for treating hypoglycaemia and as a component of parenteral nutrition.

\medterm{Gluteus Maximus Muscle}-- Largest, most superficial gluteal muscle forming the buttock bulk, from posterior iliac gluteal line, sacrum, coccyx, and sacrotuberous ligament to the iliotibial tract and gluteal tuberosity. Innervated by inferior gluteal nerve (L5-S2). Primary hip extensor and lateral rotator. Important for stair climbing, rising from seated, and posture. Common intramuscular injection site.

\medterm{Gluteus Medius Muscle}-- Thick fan-shaped muscle between maximus and minimus, from external ilium to the greater trochanter. Innervated by superior gluteal nerve (L4-S1). Primary hip abductor essential for pelvic stability during single-leg stance. Superior gluteal nerve injury causes Trendelenburg sign with contralateral pelvic drop. Commonly involved in hip pain and trochanteric bursitis.

\medterm{Gluteus Minimus Muscle}-- Smallest, deepest gluteal muscle from external ilium to the greater trochanter. Innervated by superior gluteal nerve (L4-S1). Hip abductor and medial rotator stabilizing the hip during walking. Works synergistically with medius for pelvic stability. Tendinopathy commonly causes lateral hip pain mimicking trochanteric bursitis.

\medterm{Glycogen}-- A branched polysaccharide composed of glucose units linked by alpha-1,4 glycosidic bonds (linear) and alpha-1,6 bonds (branching points, every 8-10 residues). Glycogen serves as the main storage form of glucose in animals. Stored primarily in: liver (up to 100 g, 3-5\% wet weight, released as glucose for systemic use) and skeletal muscle (approximately 400 g, 1-2\% wet weight, used locally for exercise). Glycogenesis: synthesis of glycogen from glucose (insulin-stimulated). Glycogenolysis: breakdown of glycogen to glucose-1-phosphate (glucagon-stimulated in liver, epinephrine-stimulated in muscle). Glycogen storage diseases: von Gierke (type I, G6Pase deficiency, severe fasting hypoglycaemia), Pompe (type II, acid alpha-glucosidase deficiency, fatal cardiomyopathy), Cori (type III, debranching enzyme), McArdle (type V, muscle phosphorylase deficiency, exercise intolerance).

\medterm{Gray Matter}-- Regions of the central nervous system (CNS) containing a high concentration of neuronal cell bodies (somata), dendrites, initial unmyelinated axon segments, synaptic terminals, glial cells (astrocytes and microglia), and rich capillary networks. It forms the outer folded cerebral cortex and cerebellar cortex, deep subcortical nuclei (basal ganglia, thalamus, amygdala, cerebellar nuclei), and the internal butterfly-shaped core of the spinal cord (ventral, dorsal, and lateral horns). It functions as the major site of information processing, synaptic integration, decision-making, and reflex computation. Pathologies selectively affecting gray matter include Alzheimer disease (cortical atrophy), poliomyelitis (anterior horn motor neuron destruction), and stroke.

\medterm{Greater Curvature}-- Convex lateral stomach border from gastroesophageal junction to pylorus, approximately four times longer than the lesser curvature. Attachment site for the greater omentum. Blood supply from short gastric and left gastroepicolic vessels. Common site for gastric ulcers and cancer. The gastrocolic ligament connects it to the transverse colon.

\medterm{Greater Trochanter}-- Large quadrilateral lateral proximal femoral prominence. Insertion for gluteus medius/minimus and piriformis. Palpable landmark and site of trochanteric bursa. Intertrochanteric fractures between trochanters are extracapsular with lower avascular necrosis risk because periosteal blood supply is preserved.

\medterm{Habenula}-- A paired epithalamic structure forming part of the dorsal diencephalon. It integrates limbic and basal-forebrain signals with midbrain monoaminergic systems and contributes to reward, aversion, stress, and motivational behavior.

\medterm{Hair Follicle}-- Complex dermal/hypodermal mini-organ producing hair. Has infundibulum (surface opening), isthmus (from sebaceous duct to arrector pili insertion), and lower segment (bulb with matrix and dermal papilla). Hair cycle: anagen (active growth, 2-7 years), catagen (regression, 2-3 weeks), telogen (resting, 2-3 months). Arrector pili contraction causes hair erection (goosebumps). Androgenetic alopecia results from DHT-induced follicle miniaturization. Alopecia areata is autoimmune follicle destruction.

\medterm{Hamate}-- Wedge-shaped distal row carpal bone on the ulnar side with the hook of the hamate. The hook forms Guyon's canal medial border and serves as flexor retinaculum attachment. Articulates with 4th/5th metacarpals. Hook fractures associated with racquet sports and ulnar nerve compression.

\medterm{Hamstring Muscles}-- Three posterior thigh muscles: biceps femoris (long/short heads), semitendinosus, semimembranosus, mostly from the ischial tuberosity (short head of biceps from linea aspera). Long head/semimembranosus/semitendinosus innervated by tibial sciatic nerve; short head by common peroneal division. Flex knee, extend hip. Essential for running/deceleration. Strains are among the most common sports injuries.

\medterm{Heart}-- A four-chambered muscular pump (two atria, two ventricles) located in the mediastinum, responsible for circulating blood throughout the body. The right side receives deoxygenated blood from the systemic veins and pumps it to the lungs (pulmonary circulation); the left side receives oxygenated blood from the pulmonary veins and pumps it to the systemic circulation. The heart wall has three layers: epicardium (visceral pericardium), myocardium (cardiac muscle), and endocardium. Valves ensure unidirectional flow: tricuspid, pulmonary, mitral, aortic. Coronary arteries supply the myocardium. Conduction system: SA node (pacemaker), AV node, bundle of His, Purkinje fibers. Conditions: coronary artery disease, heart failure, arrhythmia, valvular disease, cardiomyopathy, congenital defects.

\medterm{Hepatic Portal System}-- The venous network draining blood from the gastrointestinal tract, spleen, and pancreas to the liver before it reaches the systemic circulation. Formed by the union of the superior mesenteric and splenic veins behind the pancreatic neck; receives inferior mesenteric, left gastric, and other tributaries. The portal vein delivers nutrients and toxins to the liver for processing (first-pass metabolism). Portal hypertension (cirrhosis) causes esophageal varices, ascites, splenomegaly, caput medusae, and portosystemic shunting. Portal venous pressure normally 5-10 mmHg.

\medterm{Hepatic Vein}-- Veins draining the liver into the inferior vena cava. The right, middle, and left hepatic veins emerge from the posterior surface of the liver at the bare area and drain into the IVC just below the diaphragm. They carry deoxygenated blood leaving the liver (after the hepatic portal system delivers blood for processing). Hepatic vein occlusion (Budd-Chiari syndrome) causes hepatic congestion, ascites, and liver failure. Hepatic veins are key reference points for liver resection (segment anatomy) and transplantation.

\medterm{Hip Joint}-- Multiaxial ball-and-socket synovial joint between femoral head and acetabulum. Designed for stability and mobility, bearing entire upper body weight. Acetabular labrum deepens the socket; strong ligaments (iliofemoral, pubofemoral, ischiofemoral) reinforce it. Ligamentum teres contains the femoral head artery. Allows flexion, extension, abduction, adduction, rotation, circumduction. Total hip arthroplasty is a common treatment for osteoarthritis.

\medterm{Hippocampus}-- Curved medial temporal lobe structure, part of the limbic system, critical for new declarative memory formation (episodic and spatial) and short-to-long-term memory conversion. Has dentate gyrus, hippocampus proper (CA1-CA4), subiculum, and entorhinal cortex. One of few adult neurogenesis sites (dentate gyrus). Bilateral damage (herpes simplex encephalitis, bilateral PCA strokes) causes anterograde amnesia. Atrophy is an early Alzheimer's finding.

\medterm{Humeroradial Joint}-- Between capitulum and radial head, allowing flexion-extension and contributing to pronation-supination via radial head rotation. Radial collateral ligament provides lateral stability. Capitulum has articular cartilage only anteriorly. Less stable than humeroulnar joint; more susceptible to dislocation, particularly nursemaid's elbow in children.

\medterm{Humeroulnar Joint}-- Uniaxial hinge between trochlea and trochlear notch of the ulna, allowing only flexion/extension (0 to approximately 145 degrees). Ulnar collateral ligament is the primary valgus restraint. Trochlear notch provides inherent bony stability, making it the most congruent elbow joint. Commonly affected by osteoarthritis and rheumatoid arthritis.

\medterm{Humerus}-- Single upper arm bone, the longest of the upper limb. Proximal end: head, anatomical/surgical necks, tubercles. Shaft: deltoid tuberosity, radial groove, epicondyles. Distal end: trochlea, capitulum, olecranon fossa. The surgical neck is the common fracture site potentially injuring the axillary nerve. Midshaft fractures may injure the radial nerve in the radial groove.

\medterm{Hyaline Cartilage}-- The most common type of cartilage, containing fine type II collagen fibers in a firm, glassy matrix. It supports airways, covers articular surfaces, and forms much of the fetal skeleton and growth plates.

\medterm{Hyoid Bone}-- U-shaped bone in the anterior neck at C3, unique for not directly articulating with any other bone. Suspended by muscles and the stylohyoid ligament. Provides attachment for suprahyoid (digastric, stylohyoid, mylohyoid, geniohyoid) and infrahyoid (sternohyoid, sternothyroid, thyrohyoid, omohyoid) muscles. Critical for swallowing and speech. Fracture is a classic finding in manual strangulation.

\medterm{Hypodermis}-- Subcutaneous tissue/superficial fascia, the deepest skin layer with loose connective tissue and adipose tissue. Attaches skin to underlying muscles/bones via fibrous septa. Provides thermal insulation, energy storage (fat), mechanical cushioning, and contains blood vessels, lymphatics, and nerves. Subcutaneous injections are administered here. Fat distribution is hormonally influenced (estrogen promotes gynoid; testosterone promotes android). Lipoatrophy may occur from repeated insulin injections, causing subcutaneous fat loss and sometimes requiring rotation of injection sites to prevent dermal scarring.

\medterm{Hypogastric Region}-- Central inferior region below the umbilical region, containing the urinary bladder, sigmoid colon, rectum, and in females the uterus and ovaries. Assessed for bladder distension, uterine enlargement, and pelvic masses. Suprapubic pain may indicate cystitis, calculi, or pelvic pathology.

\medterm{Hypoglossal Nerve (Cranial Nerve XII)}-- The twelfth cranial nerve (CN XII), a purely somatic motor nerve supplying all the intrinsic and extrinsic muscles of the tongue (except the palatoglossus, which is supplied by the vagus nerve). Its nucleus lies in the dorsal medulla oblongata (hypoglossal trigone). The nerve exits the medulla between the pyramid and the olive, leaves the cranial cavity through the hypoglossal canal, descends in the neck deep to the posterior belly of the digastric, hooks around the occipital artery, and passes across the carotid sheath into the submandibular region. During its course, fibers from C1 hitchhike with it before forming the superior root of the ansa cervicalis and supplying the thyrohyoid and geniohyoid muscles. An isolated lower motor neuron lesion of CN XII causes flaccid paralysis and fasciculations of the ipsilateral tongue; upon protrusion, the tongue deviates toward the side of the lesion ('points toward the weak side').

\medterm{Hypothalamus}-- Small region below the thalamus forming the third ventricle's floor and lower lateral walls, weighing approximately 4 grams, the master homeostasis regulator. Controls autonomic and endocrine systems (via the pituitary through the hypophyseal portal system), temperature, hunger, thirst, sleep-wake cycles, and emotion. Key nuclei: suprachiasmatic (circadian rhythm), arcuate (appetite), supraoptic/paraventricular (ADH/oxytocin), mammillary bodies (memory). Lesions cause diabetes insipidus, hypothalamic dysfunction (temperature dysregulation, appetite disturbance), and endocrine disorders such as hypopituitarism.

\medterm{Ileum}-- Longest distal segment (3.5 m) with thinner walls, shorter villi, and Peyer's patches on the antimesenteric border. Absorbs vitamin B12, bile salts, remaining nutrients. Ileocecal valve prevents colonic backflow. Crohn's disease most commonly affects the terminal ileum; resection may cause B12 deficiency and bile salt malabsorption. Meckel's diverticulum is a vitelline duct remnant here.

\medterm{Ilium}-- Largest hip bone forming the superior portion. Body contributes to acetabulum; wing has a concave iliac fossa. Iliac crest runs ASIS to PSIS, a palpable landmark for lumbar puncture level. Sacroiliac joint articulates with sacrum. Gluteal muscles originate from its external surface; iliacus fills its fossa.

\medterm{Inferior}-- A directional term indicating a position situated nearer to the feet, toward the lower part of the body, or lower relative to another structure. Synonymous with caudal. Opposed to superior (cranial). Examples include the stomach lying inferior to the diaphragm, and the ankle being inferior to the knee.

\medterm{Inferior Mesenteric Artery}-- Third major abdominal aortic branch at L3, the hindgut artery. Supplies distal transverse colon through rectum. Branches: left colic, sigmoid, superior rectal. The marginal artery of Drummond provides SMA-IMA collateral circulation. Occlusion may cause ischemic colitis at the splenic flexure watershed (Griffiths' point).

\medterm{Inferior Vena Cava}-- Body's largest vein, formed by common iliac vein union at L5, returning lower body deoxygenated blood to the right atrium. Ascends right of the aorta, pierces the diaphragm at T8, enters the pericardium. Receives hepatic, renal, gonadal, and lumbar veins. IVC thrombosis/compression causes bilateral lower extremity edema and DVT.

\medterm{Inguinal Canal}-- An oblique passage (4 cm) through the anterior abdominal wall, from the deep (internal) inguinal ring to the superficial (external) inguinal ring, transmitting the spermatic cord in males and round ligament in females, plus ilioinguinal and genital branch of genitofemoral nerves. Boundaries: floor (inguinal ligament), roof (internal oblique, transversus abdominis), anterior wall (external oblique aponeurosis), posterior wall (transversalis fascia). Site of indirect (through deep ring, congenital) and direct (through posterior wall/Hesselbach triangle, acquired) inguinal hernias.

\medterm{Inguinal Region}-- Area at the thigh-trunk junction traversed by the inguinal ligament (ASIS to pubic tubercle). Clinically significant for inguinal hernias: indirect (through deep inguinal ring, most common) and direct (through Hesselbach's triangle). Contains the spermatic cord in males and round ligament in females, along with the inferior epigastric vessels.

\medterm{Integumentary System}-- The body system comprising the skin and its appendages: hair, nails, sweat glands, and sebaceous glands. It forms a protective barrier, limits fluid loss, contributes to thermoregulation and sensation, supports immune defence, and participates in vitamin D synthesis. The skin has an epidermis, dermis, and underlying subcutaneous tissue; the subcutaneous layer is closely associated with but is not technically part of the skin itself.

\medterm{Intermediate Cuneiform}-- Smallest cuneiform, centrally located between medial and lateral cuneiforms. Articulates with navicular, medial/lateral cuneiforms, and second metatarsal. Forms the transverse arch apex. Relatively fixed, contributing to midfoot stability. Injuries are typically associated with high-energy midfoot trauma.

\medterm{Internal Iliac Artery}-- Smaller common iliac terminal branch (hypogastric), supplying the pelvis, perineum, gluteal region, and medial thigh. Posterior division: iliolumbar, lateral sacral, superior gluteal. Anterior division: inferior gluteal, internal pudendal, obturator, vesical, uterine, vaginal, middle rectal arteries. Primary uterine blood supply, ligated during hysterectomy for hemorrhage. Embolization controls pelvic hemorrhage.

\medterm{Internal Jugular Vein}-- Descends in the carotid sheath with the internal carotid artery and vagus nerve, draining brain, face, and neck. Begins at the jugular foramen as a sigmoid sinus continuation, joining the subclavian vein to form the brachiocephalic vein. Contains backflow-preventing valves. Used for central venous catheterization. JVD assessment indicates right heart failure or elevated central venous pressure.

\medterm{Internal Oblique Muscle}-- Middle flat abdominal layer deep to external oblique, from thoracolumbar fascia, iliac crest, and lateral inguinal ligament to lower ribs, linea alba, and pubic crest. Innervated by T7-T12 and L1. Ipsilateral rotation (opposite of external oblique). Contributes to inguinal canal and abdominal wall integrity.

\medterm{Intervertebral Disc}-- A fibrocartilaginous cushion between adjacent vertebral bodies, composed of a gelatinous nucleus pulposus surrounded by a tough annulus fibrosus. It absorbs shock and permits limited movement of the vertebral column.

\medterm{Inversion}-- A specialized movement of the foot occurring at the subtalar and transverse tarsal joints, in which the sole of the foot is turned medially toward the median plane, raising the medial border of the foot. It is produced by the tibialis anterior (deep fibular nerve) and tibialis posterior (tibial nerve) muscles. Inversion is the most common mechanism of ankle injury, stretching or tearing the lateral ligament complex (particularly the anterior talofibular ligament).

\medterm{Ipsilateral}-- A relational anatomical term denoting a structure, lesion, or symptom located on or affecting the same side of the body as another reference structure. Opposed to contralateral. For example, cerebellar hemispheric lesions produce ipsilateral ataxia, and the right lung is ipsilateral to the right lobe of the liver.

\medterm{Iris}-- Colored circular muscular diaphragm between cornea and lens, controlling pupil size and light entry. Two muscles: sphincter pupillae (parasympathetic CN III, constricts in bright light) and dilator pupillae (sympathetic, dilates in darkness). Iris root attaches to the ciliary body. The anterior chamber angle between iris and cornea is where aqueous humor drains. Narrow angles may cause acute angle-closure glaucoma with sudden IOP elevation.

\medterm{Ischium}-- Inferior-posterior hip bone. Ischial tuberosity bears seated weight and serves as hamstring attachment. Ischial spine is a pudendal nerve block landmark and determines pelvic outlet level. Lesser sciatic notch lies between spine and tuberosity.

\medterm{Jejunum}-- Middle segment (2.5 m) with prominent circular folds and tall villi for a velvety appearance. Richer blood supply and thicker walls than the ileum. Primary nutrient absorption site for carbohydrates, amino acids, fatty acids, vitamins, minerals. Fewer Peyer's patches. More commonly affected by Crohn's disease. Robust blood supply preferred for surgical anastomoses.

\medterm{Joint}-- Articulation between two or more bones permitting movement. Structural classification: fibrous (sutures, syndesmoses, gomphoses -- immovable/slightly movable), cartilaginous (synchondroses, symphyses -- slightly movable), synovial (freely movable -- plane, hinge, pivot, condyloid, saddle, ball-and-socket). Synovial joint features: articular cartilage (hyaline), joint capsule (fibrous outer, synovial inner), synovial fluid (lubrication, nutrition), ligaments (stability), menisci/labrum (congruence). Conditions: osteoarthritis, rheumatoid arthritis, septic arthritis, dislocation, sprain.

\medterm{Jugular Vein}-- One of the major veins draining blood from the head, face, and neck. There are two pairs: internal jugular vein (IJV) and external jugular vein (EJV). Internal jugular vein: the larger vessel, runs within the carotid sheath (medial to the carotid artery), descends from the jugular foramen (at the base of the skull) to the brachiocephalic vein. Receives blood from the brain (via dural venous sinuses, sigmoid sinus), face, and superficial neck. The IJV is the preferred site for central venous access (internal jugular central line) because of its straight course, predictable anatomy, and lower risk of pneumothorax compared to the subclavian approach. It is also used for measuring central venous pressure (CVP) and for venous sampling in endocrine testing (inferior petrosal sinus sampling for ACTH-dependent Cushing syndrome). External jugular vein: superficial vein formed by the retromandibular and posterior auricular veins, descends obliquely across the sternocleidomastoid muscle, pierces the deep fascia, and drains into the subclavian vein. Easily visible when it distends (Valsalva, heart failure, superior vena cava obstruction). The EJV is used for peripheral venous access when other sites are compromised. Jugular venous pressure (JVP): the height of the IJV pulsation in the neck reflects right atrial pressure, assessed with the patient at 45 degrees; the a-wave reflects atrial contraction (absent in atrial fibrillation, cannon a-wave in tricuspid atresia, AV dissociation), and the v-wave reflects atrial filling during ventricular systole (giant v-wave in tricuspid regurgitation).

\medterm{Kidney}-- A paired, bean-shaped retroperitoneal organ responsible for filtering blood, excreting waste products as urine, and regulating fluid and electrolyte balance, acid-base homeostasis, and blood pressure (via the renin-angiotensin-aldosterone system). Each kidney contains approximately 1 million nephrons, the functional filtering units. Structure: outer renal cortex, inner renal medulla (organized into renal pyramids), renal pelvis (collecting urine). The renal artery supplies blood; the renal vein drains it. Kidneys also produce erythropoietin (stimulates red blood cell production) and activate vitamin D. Common diseases: chronic kidney disease, nephrolithiasis, glomerulonephritis, pyelonephritis.

\medterm{Knee Joint}-- Largest, most complex joint with three compartments: medial/lateral tibiofemoral and patellofemoral. Tibiofemoral is a modified hinge allowing flexion-extension with rotation when flexed. Stabilized by cruciate (ACL/PCL) and collateral (MCL/LCL) ligaments; menisci cushion and distribute load. The suprapatellar bursa is the body's largest. Common injuries: ACL tears, meniscal tears, patellar dislocations.

\medterm{Lacrimal Apparatus}-- Produces, distributes, and drains tears for corneal moisture and protection. Components: lacrimal gland (superolateral orbit, parasympathetic CN V1, produces aqueous tear layer), superior/inferior puncta (medial eyelid openings), canaliculi, lacrimal sac (lacrimal groove), and nasolacrimal duct (drains into inferior meatus). Dry eye syndrome (keratoconjunctivitis sicca) from decreased production or excessive evaporation is common. Nasolacrimal duct obstruction causes epiphora and may cause dacryocystitis (lacrimal sac infection) presenting with epiphora and mucopurulent discharge.

\medterm{Lacrimal Gland}-- The paired almond-shaped exocrine glands that secrete the aqueous layer of the tear film, located in the lacrimal fossa of the superior lateral orbit (orbital part) and the upper eyelid (palpebral part). The lacrimal gland is a tubuloacinar serous gland, innervated by parasympathetic fibers from the facial nerve (CN VII) via the pterygopalatine ganglion and the zygomatic and lacrimal nerves. Its secretion is stimulated by corneal irritation, emotional responses, and trigeminal reflexes. Tears drain via the lacrimal puncta (upper and lower) into the lacrimal canaliculi, common canaliculus, lacrimal sac, and nasolacrimal duct, emptying into the inferior meatus of the nasal cavity. The tear film has three layers: lipid (meibomian glands, prevents evaporation), aqueous (lacrimal gland, contains lysozyme, lactoferrin, IgA, electrolytes), and mucin (conjunctival goblet cells, anchors tears to cornea). Diseases: dacryoadenitis (inflammation, usually viral or autoimmune--Sjogren syndrome, sarcoidosis), dacryocystitis (infection of the lacrimal sac, commonly obstructed nasolacrimal duct), and tumors (pleomorphic adenoma most common benign, adenoid cystic carcinoma most common malignant).

\medterm{Large Intestine}-- Final GI segment (approximately 1.5 meters) from ileocecal valve to rectum. Six segments: cecum, ascending, transverse, descending, sigmoid colon, rectum. Absorbs water, electrolytes, and vitamins (K, biotin) from colonic bacteria. Characterized by haustra, taeniae coli, and epiploic appendages. Houses approximately 100 trillion bacteria forming a diverse microbiome essential for metabolism and immunity.

\medterm{Larynx}-- A cartilaginous organ between the pharynx and trachea that maintains an open airway, protects the lower respiratory tract during swallowing, and produces voice. It contains the vocal folds and epiglottis.

\medterm{Lateral}-- A directional term indicating a position situated further away from the median plane (midline) of the body, or toward the outer flank/side. Opposed to medial. Examples include the ears being lateral to the nose, the thumb (pollex) being on the lateral side of the hand in anatomical position, and the fibula being lateral to the tibia in the leg.

\medterm{Lateral Cuneiform}-- Located between intermediate cuneiform and cuboid, articulating with navicular, intermediate cuneiform, cuboid, and third metatarsal. Contributes to the transverse arch and lateral foot column. Injuries typically from high-energy midfoot trauma.

\medterm{Lateral Rotation}-- An external rotation of a limb or body segment around its long axis away from the median plane (midline) of the body. At the glenohumeral joint, it turns the anterior humerus laterally (produced by infraspinatus and teres minor). At the hip, it turns the anterior femur laterally (produced by gluteus maximus, obturator internus, superior and inferior gemelli, quadratus femoris, and piriformis). Opposed to medial rotation.

\medterm{Latissimus Dorsi Muscle}-- Broadest back muscle forming the posterior axillary wall, from T7-L5 spinous processes, thoracolumbar fascia, iliac crest, and lower ribs to the intertubercular groove floor. Innervated by thoracodorsal nerve (C6-C8). Adducts, medially rotates, extends the arm powerfully in swimming/rowing/climbing. Common reconstructive surgery flap. Susceptible to overhead strain.

\medterm{Left Atrium}-- Receives oxygenated blood from 4 pulmonary veins. Most posterior chamber forming the cardiac base on chest radiography. Smooth walls from absorbed pulmonary veins; small auricle with pectinate muscles. Mitral valve separates it from the left ventricle. Enlargement is common in mitral disease and atrial fibrillation, potentially causing esophageal compression or recurrent laryngeal nerve hoarseness.

\medterm{Left Hypochondriac Region}-- Upper left abdomen lateral to the epigastric region, containing the spleen, stomach fundus, splenic flexure, and pancreatic tail. Important for evaluating splenomegaly, gastric pathology, and splenic infarction. The spleen is normally not palpable below the costal margin; enlargement into this region is a significant finding.

\medterm{Left Iliac Region}-- Lower left abdomen containing the descending colon, sigmoid colon, and left ureter. Important for evaluating diverticulosis and diverticulitis, most commonly affecting the sigmoid colon here. Palpation may reveal masses from inflammatory bowel disease or colorectal cancer.

\medterm{Left Ventricle}-- Thickest cardiac chamber (13-15 mm), pumping oxygenated blood into the aorta for systemic distribution. Generates approximately 120 mmHg systolic pressure. Receives blood through the mitral valve; ejects through the aortic valve. Left ventricular hypertrophy from chronic hypertension or aortic stenosis is a risk factor for heart failure, arrhythmias, and sudden death.

\medterm{Lens}-- Transparent biconvex avascular structure behind the iris/pupil, suspended by zonular fibers from the ciliary body. Provides approximately one-third of refractive power and allows accommodation (focusing on near objects by changing shape). Concentric lens fiber layers from surface epithelium, surrounded by a capsule. Absorbs ultraviolet light. Cataracts (opacification) are the most common cause of reversible blindness worldwide. Accommodation decreases with age (presbyopia).

\medterm{Lesser Curvature}-- Concave medial stomach border from gastroesophageal junction to pylorus. Attachment site for the lesser omentum connecting to the liver, containing right/left gastric arteries. Most common gastric ulcer site. The angular incisure marks the body-antrum junction. Cancer here may cause early obstruction from the narrow lumen.

\medterm{Lesser Trochanter}-- Small conical posteromedial proximal femoral prominence. Insertion for ilioposmus (primary hip flexor). Isolated fractures in adults may indicate pathologic fracture from metastatic disease; in adolescents may result from iliopsoas avulsion. Visible on lateral hip radiographs.

\medterm{Ligament}-- A dense, fibrous connective tissue band that connects bone to bone, providing passive joint stability and guiding joint motion. Ligaments are composed predominantly of type I collagen fibers (parallel orientation, crimped pattern) and elastin fibers (10-20\% in some ligaments), arranged in bundles. The collagen fibers give ligaments high tensile strength (up to 30-40 N/mm2 for the ACL), while elastin allows elastic recoil. Ligaments have a limited blood supply (attachments are relatively avascular), leading to slow healing after injury. Mechanoreceptors (Ruffini, Pacinian, Golgi-type) within ligaments provide proprioceptive feedback about joint position and tension. The ligament response to injury: acute inflammatory phase (days), proliferative phase (weeks, fibroblast proliferation, collagen deposition, initially type III then type I), and remodelling phase (months-years, orientation of collagen along lines of stress). Classification of ligament injuries (grades): grade 1 (mild stretching, microscopic tears, no instability), grade 2 (partial tear, some laxity), grade 3 (complete rupture, gross instability). Healing is better in intra-articular ligaments (ACL has poor healing potential, often requires surgical reconstruction) vs extra-articular (MCL heals well non-operatively). Ligaments weaken with age and corticosteroid use, and sprains are common ankle sprain (anterior talofibular ligament most common), ACL tear (knee, pivot sport), MCL tear (knee, valgus stress), and ulnar collateral ligament (thumb, skier's thumb/gamekeeper's thumb).

\medterm{Linea Aspera}-- Prominent longitudinal posterior femoral shaft ridge from converging muscle attachment lines. Insertion for adductor magnus, vastus medialis/lateralis, and short head of biceps femoris. Thickest femoral cortex, resisting bending forces. Diverges distally into supracondylar lines. The profunda femoris artery pierces the adductor magnus near it.

\medterm{Liver}-- The largest internal organ, located in the right upper quadrant of the abdomen beneath the diaphragm. Functions: metabolism of carbohydrates, lipids, and proteins; detoxification of drugs and toxins; synthesis of proteins (albumin, clotting factors, complement); bile production for fat digestion; storage of glycogen, vitamins, and iron; immune function (Kupffer cells). Dual blood supply: hepatic artery (oxygenated) and portal vein (nutrient-rich from intestines). Drains into hepatic veins to IVC. Bile drains via biliary tree to gallbladder and duodenum. Regenerative capacity allows recovery from partial resection. Conditions: hepatitis, cirrhosis, fatty liver, liver cancer, gallstones, drug-induced liver injury.

\medterm{Liver Lobule}-- The functional unit of the liver: a hexagonal structural unit centered on a central (hepatic) venule, with portal triads (portal vein, hepatic artery, bile ductule) at the corners. Hepatocytes arranged in cords radiate from the center; sinusoids (lined by Kupffer cells) receive mixed portal venous and hepatic arterial blood flowing toward the central vein; bile canaliculi flow in the opposite direction. Zones 1-3 differ in oxygen and metabolic function (zone 3 most vulnerable to ischemia/toxins—centrilobular necrosis). Alternative concepts: portal lobule and hepatic acinus (Rappaport).

\medterm{Lumbar Vertebrae (L1-L5)}-- Largest and strongest mobile vertebrae for weight-bearing. Large kidney-shaped bodies, short hatchet-shaped spinous processes, sagittally-oriented facets allowing primarily flexion/extension. Lordotic curvature. Most common disc herniation site, particularly L4-L5 and L5-S1 causing sciatica. Lumbar puncture performed between L3-L4 or L4-L5.

\medterm{Lunate}-- Crescent-shaped proximal row carpal bone between scaphoid and triquetrum. Most commonly dislocated carpal bone from falls on extended wrists. Dislocation may compress the median nerve causing acute carpal tunnel syndrome. Vulnerable to avascular necrosis (Kienbock disease) from its precarious blood supply.

\medterm{Lung}-- The paired, cone-shaped respiratory organs located in the thoracic cavity, responsible for gas exchange (oxygen uptake, carbon dioxide elimination). Each lung is divided into lobes: the right lung has three (superior, middle, inferior), the left lung has two (superior, inferior) with a cardiac notch accommodating the heart. Structure: trachea branches into bronchi, bronchioles, alveolar ducts, and alveoli (300--500 million, providing a surface area of approximately 70--100 m\textsuperscript{2}). The pleura (visceral and parietal) surrounds each lung. Ventilation is driven by the diaphragm and intercostal muscles. Common diseases: pneumonia, COPD, asthma, lung cancer, pulmonary embolism.

\medterm{Lymphatic System}-- A network of lymphatic vessels, lymph nodes, and lymphoid organs (tonsils, thymus, spleen, Peyer patches) that returns interstitial fluid (lymph) to the bloodstream, transports dietary fat (chyle), and serves as the conduit for immune surveillance. Lymphatic capillaries collect excess tissue fluid; lymph passes through lymph nodes (filtration, antigen presentation), then trunks and ducts (thoracic duct on the left, right lymphatic duct) into the venous system. Dysfunction causes lymphedema; lymphangitic spread is a major route of cancer metastasis.

\medterm{Lymph Node Structure}-- Small bean-shaped lymphoid organs (1-25 mm) distributed along lymphatic vessels, filtering lymph and mounting immune responses. Each has a cortex (outer, containing B-cell follicles with germinal centers), paracortex (deep cortex, containing T-cells and dendritic cells), and medulla (inner, containing medullary cords and sinuses). Afferent vessels enter the cortex; efferent vessels exit at the hilum (also containing the artery, vein, and efferent lymphatic). Macrophages and dendritic cells within the node process and present antigens to lymphocytes, initiating adaptive immune responses and germinal center formation.

\medterm{Macula Lutea}-- Small oval yellowish central retinal area approximately 5.5 mm in diameter for high-acuity vision. Yellow color from xanthophyll pigments (lutein, zeaxanthin) filtering blue light and protecting photoreceptors. The fovea centralis is the central depression (1.5 mm) with highest cone density and no rods. Fovea is avascular, supplied by underlying choriocapillaris. Age-related macular degeneration is the leading cause of irreversible vision loss in developed countries.

\medterm{Mammary Gland}-- Modified sweat (apocrine) gland in the breast, overlying pectoralis major, with 15-20 lactiferous ducts converging on the nipple. Composed of glandular lobules (alveoli), fibrous stroma (Cooper's ligaments), and fat. Development regulated by estrogen, progesterone, prolactin, and oxytocin (milk let-down). Lymphatic drainage: primarily to axillary nodes (level I-III), also internal mammary and supraclavicular—critical for breast cancer staging. Blood supply: internal thoracic, lateral thoracic, intercostal arteries.

\medterm{Mandible}-- Largest facial bone forming the lower jaw, with horseshoe-shaped body and two rami ending in condylar (TMJ) and coronoid processes. Mental foramen transmits the mental nerve. Mandibular foramen is the entry point for inferior alveolar nerve blocks. Most commonly fractured facial bone, typically at the condyle, angle, or symphysis.

\medterm{Masseter Muscle}-- Thick rectangular muscle of mastication on the mandibular ramus, from the zygomatic arch to mandibular angle/ramus. Innervated by the masseteric nerve (V3). Most powerful masticatory muscle, elevating the mandible with up to 70 kg force. Involved in bruxism, TM disorders, and myofascial pain. Hypertrophy can widen facial appearance.

\medterm{Maxilla}-- Paired bone forming the upper jaw, central face, and parts of orbit, nasal cavity, and hard palate. Body contains the maxillary sinus (largest paranasal sinus). Four processes: frontal, zygomatic, alveolar (upper teeth), and palatine (anterior hard palate). Fractures classified using the Le Fort system. The infraorbital nerve emerges from its foramen.

\medterm{Medial}-- A directional term indicating a position situated closer to the median plane (midline) of the body, or toward the middle of a structure. Opposed to lateral. Examples include the sternum being medial to the ribs, the hallux (great toe) being on the medial side of the foot, and the ulna being the medial bone of the forearm in anatomical position.

\medterm{Medial Cuneiform}-- Largest cuneiform, on the medial foot between navicular and first metatarsal. Insertion for tibialis anterior and posterior tendons supporting the medial longitudinal arch. The first tarsometatarsal joint (Lisfranc joint) is an important landmark; Lisfranc injuries involve fracture-dislocation at this level.

\medterm{Medial Rotation}-- An internal rotation of a limb or body segment around its long axis toward the median plane (midline) of the body. At the glenohumeral joint (shoulder), it turns the anterior humerus medially (produced by subscapularis, pectoralis major, latissimus dorsi, and teres major). At the hip, it turns the anterior femur medially (produced by gluteus medius, gluteus minimus, and tensor fasciae latae). Opposed to lateral rotation.

\medterm{Median Nerve}-- From lateral and medial cords (C5-T1), passing through the arm without branching, through the cubital fossa, into the forearm innervating most flexor muscles. In the hand innervates thenar muscles (opponens, abductor brevis, flexor brevis pollicis) and first two lumbricals. Sensation to lateral 3.5 digits. Carpal tunnel syndrome is the most common entrapment, causing thenar wasting and lateral palm/finger sensory loss.

\medterm{Median Plane}-- A sagittal plane that passes directly through the midline of the body, dividing it into equal right and left halves.

\medterm{Mediastinum}-- The central interpleural tissue space and compartment of the thoracic cavity bounded anteriorly by the sternum, posteriorly by the thoracic vertebral column (T1--T12), laterally by the mediastinal parietal pleura of each lung, superiorly by the superior thoracic aperture (thoracic inlet), and inferiorly by the thoracic diaphragm. It is divided into superior and inferior mediastina by a transverse anatomical plane passing from the sternal angle of Louis anteriorly to the T4--T5 intervertebral disc posteriorly. The superior mediastinum contains the thymus, brachiocephalic veins, superior vena cava, aortic arch with its three major branches, trachea, esophagus, thoracic duct, phrenic nerves, vagus nerves, and left recurrent laryngeal nerve. The inferior mediastinum is further subdivided into anterior (containing thymus remnants, internal thoracic vessels, lymph nodes), middle (containing the pericardium, heart, ascending aorta, pulmonary trunk, superior and inferior vena cava terminations, phrenic nerves), and posterior (containing the descending thoracic aorta, azygos and hemiazygos veins, thoracic duct, esophagus with esophageal plexus, vagus nerves, and sympathetic trunks). Pathologies include mediastinitis, mediastinal masses (e.g., thymoma, teratoma, thyroid goiter, 'terrible' lymphoma in anterior compartment; neurogenic tumors in posterior compartment), and pneumomediastinum.

\medterm{Medulla Oblongata}-- Most inferior brainstem segment, continuous with the spinal cord at the foramen magnum. Contains vital autonomic centers: cardiovascular center (heart rate, blood pressure), respiratory center (dorsal/ventral respiratory groups), and area postrema (vomiting center). Houses the pyramidal decussation (corticospinal tracts crossing, explaining contralateral motor control), olives (inferior olivary nuclei for cerebellar motor learning), and cranial nerve IX-XII nuclei. Lesions (Chiari malformation, syringobulbia, lateral medullary syndrome) cause ipsilateral cranial nerve deficits and contralateral long tract signs, often with vertigo, dysphagia, and respiratory compromise.

\medterm{Melanocytes}-- Dendritic pigment-producing neural crest-derived cells in the epidermal basal layer and hair follicles, comprising approximately 5-10 percent of basal keratinocytes. Produce melanin (eumelanin, brown-black; pheomelanin, yellow-red) in melanosomes, transferred to keratinocytes via dendritic processes for UV protection. Each melanocyte supplies approximately 30-40 keratinocytes (epidermal melanin unit). UV exposure stimulates melanin production (tanning). Cell of origin for melanoma, a potentially fatal skin cancer that metastasises early and requires prompt surgical excision with clear margins.

\medterm{Meninges}-- The three membranes enveloping the brain and spinal cord: dura mater (outermost, tough), arachnoid mater (middle, avascular), and pia mater (innermost, adherent to the CNS surface). The subarachnoid space (between arachnoid and pia) contains cerebrospinal fluid. The epidural and subdural spaces are sites of hemorrhage. Inflammation of the meninges (meningitis) and meningeal carcinomatosis (tumor spread) are serious clinical conditions. Dural reflections form the falx cerebri and tentorium cerebelli; dural venous sinuses lie between dural layers.

\medterm{Meniscus (Knee)}-- A pair of crescentic, wedge-shaped fibrocartilaginous pads (medial and lateral menisci) interposed between the femoral condyles and the tibial plateau within the knee joint. Composed primarily of type I collagen organized into circumferential and radial fibers. Functions include shock absorption (transmitting ~50--70 percent of joint load), deepening the shallow tibial articular surfaces to improve joint congruency, lubricating the joint, and proprioception. The medial meniscus is larger, C-shaped, and firmly anchored to the tibial collateral ligament (MCL), making it less mobile and approximately three times more prone to injury than the smaller, more circular, mobile lateral meniscus. Vascularity is limited: only the outer 10--30 percent ('red-red zone') has blood supply from genicular arteries; tears in the avascular inner zone ('white-white zone') cannot heal spontaneously. Tears present with joint line tenderness, clicking, effusion, and mechanical locking; diagnosed with the McMurray and Apley grind tests, confirmed by MRI.

\medterm{Mesentery}-- The fan-shaped double layer of peritoneum attaching the small intestine (jejunum and ileum) to the posterior abdominal wall, carrying the superior mesenteric artery, vein, lymphatics, and autonomic nerves between its layers. Root extends from duodenojejunal junction to ileocecal junction. Related mesenteries: transverse mesocolon, sigmoid mesocolon, mesoappendix. Inflammation of the mesentery (mesenteric adenitis, ischemic mesenteric disease) causes abdominal pain; mesenteric ischemia/volvulus is a surgical emergency. The mesentery functions to suspend and provide blood supply to the gut.

\medterm{Metacarpals}-- Five long bones (I-V) forming the palm framework. Each has base, shaft, and head. The first is shortest and most mobile, articulating with the trapezium at a saddle joint for opposition. Second/third are firmly fixed; fourth/fifth more mobile. Bennett's fracture affects the first metacarpal base; boxer's fracture affects the fifth neck.

\medterm{Metacarpophalangeal Joint}-- A condyloid synovial joint between a metacarpal head and the base of a proximal phalanx. It permits flexion, extension, abduction, adduction, and limited circumduction.

\medterm{Metaphysis}-- The region of a long bone between the epiphysis (end) and diaphysis (shaft), containing the epiphyseal growth plate (physis) in children—the zone of active cartilage growth (endochondral ossification). The metaphysis is highly vascularized with abundant trabecular bone; common site of osteomyelitis (sluggish blood flow) and bone tumors (osteosarcoma, Ewing sarcoma) in children. After skeletal maturity, the growth plate closes and the metaphysis fuses with the epiphysis.

\medterm{Metatarsals}-- Five long bones (I-V) forming the forefoot. Each has base, shaft, and head. First is shortest/thickest, bearing approximately one-third of forefoot load. Stress fractures commonly affect the second and third (march fractures). Morton's neuroma typically affects the nerve between 3rd/4th metatarsal heads. MTP joints are condyloid allowing flexion, extension, limited abduction/adduction.

\medterm{Midaxillary Line}-- An imaginary vertical anatomical reference line passing through the apex of the axilla (armpit), midway between the anterior axillary line (formed by the lower border of pectoralis major) and the posterior axillary line (formed by the latissimus dorsi and teres major). It is a key clinical landmark for chest tube (tube thoracostomy) placement (typically in the 4th or 5th intercostal space along the anterior or midaxillary line, within the 'triangle of safety'), thoracocentesis, and percussion of the spleen and lung bases.

\medterm{Midbrain}-- Most superior brainstem segment connecting pons to diencephalon. Contains cerebral peduncles (corticospinal/corticobulbar tracts), substantia nigra (dopaminergic neurons degenerated in Parkinson's), red nucleus (rubrospinal tract for coordination), and periaqueductal gray (pain modulation). Superior/inferior colliculi on the tectum process visual/auditory reflexes. The cerebral aqueduct connects third/fourth ventricles. Syndromes (Weber's, Benedikt's, Claude's, Nothnagel's) result from vascular lesions affecting specific midbrain regions, producing characteristic combinations of ipsilateral cranial nerve palsy and contralateral hemiplegia or ataxia.

\medterm{Midclavicular Line}-- An imaginary vertical anatomical reference line passing through the precise midpoint of the clavicle, running parallel to the midsternal line down the anterior thoracic and abdominal walls. It serves as a standard reference for clinical examination and surface anatomy: the apex beat (cardiac apical impulse) is normally palpated in the 5th left intercostal space along or just medial to the midclavicular line; the liver span is standardly percussed along the right midclavicular line (normal 6--12 cm); and the gallbladder fundus is located where the right midclavicular line crosses the 9th costal cartilage.

\medterm{Mitral Valve}-- Left atrioventricular (bicuspid) valve with anterior and posterior cusps supported by chordae tendineae and papillary muscles. Anterior cusp is larger and proximate to the aortic valve. Mitral valve prolapse is the most common valvular abnormality. Mitral stenosis (most commonly rheumatic) and regurgitation are significant. Most commonly affected valve in rheumatic heart disease.

\medterm{Motor Cortex}-- The frontal-lobe regions that plan and initiate voluntary movement. The primary motor cortex lies in the precentral gyrus and controls the opposite side of the body in an orderly somatotopic map.

\medterm{Mouth}-- The initial entry portal and segment of the alimentary canal, comprising the oral vestibule (the cleft between the lips/cheeks and the teeth/gingiva) and the oral cavity proper (bounded anterolaterally by the teeth and alveolar arches, superiorly by the hard and soft palates, inferiorly by the muscular floor of the mouth and tongue, and posteriorly communicating with the oropharynx via the faucial isthmus). Lined by protective stratified squamous epithelium (partially keratinized over the hard palate and gingiva), it receives secretions from the major and minor salivary glands and initiates both mechanical digestion (mastication by teeth) and chemical digestion (salivary amylase and lingual lipase).

\medterm{Muscle}-- A specialized contractile tissue responsible for movement, posture, and heat production. Three types: skeletal muscle (voluntary, striated, attached to bones via tendons, responsible for locomotion), cardiac muscle (involuntary, striated, found only in the heart, with intercalated discs for synchronized contraction), and smooth muscle (involuntary, non-striated, found in walls of blood vessels, gastrointestinal tract, bladder, uterus). Skeletal muscle structure: epimysium (outer sheath), perimysium (surrounds fascicles), endomysium (surrounds individual fibers). Contraction mechanism: sliding filament theory (actin-myosin cross-bridge cycling, ATP-dependent, calcium-regulated via troponin-tropomyosin).

\medterm{Muscle Tissue}-- One of the four fundamental tissue types, specialized for excitability and contractility to produce movement, maintain posture, propel fluids, and generate body heat through shivering and metabolism. It comprises elongated contractile cells (muscle fibers or myocytes) containing organized filaments of actin and myosin that slide past one another in response to intracellular calcium release. Classified into three morphologically and functionally distinct types: (1) Skeletal muscle: voluntary, striated, multinucleated cylindrical fibers attached to bones via tendons, innervated by somatic motor neurons; (2) Cardiac muscle: involuntary, striated, branching myocytes with central nuclei, joined end-to-end by intercalated discs containing desmosomes and gap junctions for electrical synchrony, found exclusively in the myocardium; and (3) Smooth muscle: involuntary, non-striated, spindle-shaped cells with a single central nucleus, found in the walls of blood vessels and hollow viscera, regulated by the autonomic nervous system, hormones, and local metabolic factors. Pathologies include muscular dystrophies (e.g., Duchenne, dystrophin mutation), rhabdomyolysis, myasthenia gravis, and cardiomyopathy.

\medterm{Musculocutaneous Nerve}-- From lateral cord (C5-C7), piercing coracobrachialis, running between biceps and brachialis to supply these muscles and lateral forearm skin (lateral antebrachial cutaneous nerve). Sole motor supply to the anterior arm compartment. Injury causes weakness of elbow flexion/forearm supination with lateral forearm sensation loss. Vulnerable in axillary injuries and proximal humeral fractures.

\medterm{Myelin}-- A lipid-rich protective sheath that wraps around nerve fibers (axons), insulating them and speeding up electrical signal transmission.

\medterm{Myelin Sheath}-- The lipid-rich insulating layer around axons formed by Schwann cells (peripheral nervous system) and oligodendrocytes (central nervous system), which wrap the axon in concentric membrane layers. Myelin provides electrical insulation and enables saltatory conduction (action potentials jump between nodes of Ranvier), greatly increasing conduction velocity and reducing energy demands. Damage/demyelination (multiple sclerosis, Guillain-Barre syndrome, vitamin B12 deficiency) slows or blocks conduction, causing weakness, sensory loss, and autonomic dysfunction.

\medterm{Myocardium}-- The thick, middle muscular layer of the heart wall that constitutes the bulk of the cardiac mass, responsible for pumping blood through the pulmonary and systemic circulations. Histologically, it consists of involuntary, striated cardiac muscle cells (cardiomyocytes) with central nuclei, branching architecture, and abundant mitochondria (accounting for ~40 percent of cell volume). Cardiomyocytes are joined end-to-end by intercalated discs containing fascia adherens (mechanical anchoring), macula adherens/desmosomes, and low-resistance gap junctions that permit rapid electrical coupling, allowing the myocardium to behave as an electrical and functional syncytium. The left ventricular myocardium is approximately three times thicker (12--15 mm) than the right ventricular myocardium (3--5 mm), reflecting the higher systemic afterload. Pathologies include myocardial infarction (ischemic myocyte necrosis), myocarditis (inflammation, often viral), and cardiomyopathies (dilated, hypertrophic, restrictive).

\medterm{Nasal Cavity}-- The air-filled cavity extending from the external nares to the choanae. It warms, humidifies, and filters inspired air; its superior region contains the olfactory epithelium.

\medterm{Navicular}-- Boat-shaped tarsal bone between talus and three cuneiforms. Navicular tuberosity is a palpable medial prominence serving as posterior tibial tendon insertion, key medial longitudinal arch stabilizer. Central third has a precarious blood supply, susceptible to stress fractures and avascular necrosis (Kohler disease in children). Most commonly fractured midfoot tarsal bone.

\medterm{Nephron}-- The structural and functional unit of the kidney, consisting of a renal corpuscle and renal tubule. It filters blood and modifies the filtrate to form urine.

\medterm{Nerve}-- A bundle of axons (nerve fibers) surrounded by connective tissue sheaths that transmits electrical impulses between the central nervous system and the periphery. Structure: epineurium (outer layer), perineurium (surrounds fascicles), endoneurium (surrounds individual axons). Types: sensory (afferent, carry impulses toward CNS), motor (efferent, carry impulses away from CNS to muscles/glands), and mixed nerves. Myelinated axons conduct impulses faster (saltatory conduction) than unmyelinated axons. Cranial nerves (12 pairs) arise from the brain; spinal nerves (31 pairs) arise from the spinal cord. Nerve regeneration is limited in the CNS but possible in the PNS.

\medterm{Nervous Tissue}-- One of the four basic tissue types, forming the central and peripheral nervous systems, specialized for detecting sensory stimuli, processing and integrating information, and transmitting electrical impulses to coordinate bodily activities. It consists of two main cellular populations: (1) Neurons: excitable conducting units composed of a soma (cell body) containing the nucleus and Nissl bodies, dendrites that receive synaptic inputs, and a single axon that transmits action potentials away from the soma to synaptic terminals; and (2) Neuroglia (glial cells): non-excitable supporting cells that outnumber neurons, including astrocytes (blood-brain barrier maintenance, metabolic support, glutamate recycling), oligodendrocytes (CNS myelination), microglia (resident CNS phagocytes), ependymal cells (ventricular lining, CSF production), Schwann cells (PNS myelination), and satellite cells (sensory/autonomic ganglion support). Pathologies include neurodegenerative diseases (Alzheimer, Parkinson), demyelinating conditions (multiple sclerosis, Guillain-Barré syndrome), gliomas, and stroke.

\medterm{Neuron}-- The fundamental excitable cell of the nervous system specialized for receiving, integrating, and transmitting electrical and chemical signals. Structural components: cell body (soma), dendrites (receptive), axon (conducting), and synaptic terminals (releasing neurotransmitters). Types: sensory, motor, and interneurons; classified by structure (unipolar, bipolar, multipolar) and function. Neurons are highly polarized cells with resting membrane potentials and generate action potentials; they rely on axonal transport and are largely post-mitotic (limited regeneration). Neurotransmitters include glutamate, GABA, dopamine, serotonin, and acetylcholine.

\medterm{Nose}-- The external nasal structure and nasal cavities, the primary organ of olfaction and the entry to the respiratory tract. External nose: bone (nasal bones, frontal process of maxilla) and cartilage (septal, upper/lower lateral cartilages). Nasal cavity: divided by septum, lined by respiratory epithelium (ciliated pseudostratified columnar) and olfactory epithelium (superior turbinate/septum). Turbinates (conchae): superior, middle, inferior (increase surface area, warm/humidify/filter air). Paranasal sinuses drain into nasal cavity (maxillary, frontal, ethmoid, sphenoid). Blood supply: sphenopalatine, greater palatine, facial arteries. Conditions: rhinitis, sinusitis, epistaxis, deviated septum, nasal polyps, olfactory loss.

\medterm{Oblique Plane}-- A plane that passes through the body or an organ at an angle other than a right angle (diagonal plane).

\medterm{Obturator}-- A structure that occludes an opening. In anatomy, the obturator foramen is a large opening in the pelvic bone, closed by the obturator membrane. The obturator nerve (L2--L4) innervates the adductor muscles of the thigh and provides sensory innervation to the medial thigh. The obturator artery supplies the adductor compartment. The term also applies to surgical obturators—prosthetic devices used to close palatal defects (cleft palate, post-surgical).

\medterm{Occipital Bone}-- Forms the posterior-inferior cranium surrounding the foramen magnum (spinal cord passage). Has basilar, lateral, and squamous parts. Occipital condyles articulate with C1 for nodding movement. The external occipital protuberance provides muscle attachment. Fractures through the foramen magnum may cause brainstem injury and are life-threatening.

\medterm{Occipital Lobe}-- The most posterior of the four cerebral lobes, located behind the parieto-occipital sulcus and above the tentorium cerebelli. The occipital lobe is the primary visual processing center of the brain. The primary visual cortex (V1, Brodmann area 17) receives input from the lateral geniculate nucleus via the optic radiations; damage causes contralateral homonymous hemianopia or quadrantanopia. Higher-order visual areas (V2--V5) process colour, motion, and object recognition. The lingual gyrus and cuneus are key occipital landmarks. Lesions beyond V1 can cause visual agnosia, alexia without agraphia, and visual hallucinations (peduncular hallucinosis).

\medterm{Oculomotor Nerve (Cranial Nerve III)}-- The third cranial nerve (CN III), predominantly a motor nerve containing somatic motor fibers to extraocular muscles and parasympathetic fibers to intraocular smooth muscles. Its nuclei lie in the midbrain at the level of the superior colliculus: the oculomotor nucleus (somatic motor) and the accessory oculomotor (Edinger-Westphal) nucleus (parasympathetic). The nerve emerges from the interpeduncular fossa, courses between the posterior cerebral and superior cerebellar arteries, travels through the lateral wall of the cavernous sinus, and enters the orbit via the superior orbital fissure within the common tendinous ring. Somatic motor fibers innervate four extraocular muscles: superior rectus (elevation), medial rectus (adduction), inferior rectus (depression), and inferior oblique (elevation/extorsion), plus the levator palpebrae superioris (eyelid elevation). Preganglionic parasympathetic fibers synapse in the ciliary ganglion, giving rise to short ciliary nerves that supply the sphincter pupillae (miosis) and ciliary muscle (accommodation for near vision). Complete CN III palsy produces ptosis, a 'down-and-out' abducted and depressed eye, and a fixed, dilated pupil.

\medterm{Odontoid Process}-- A tooth-like projection superiorly from the body of the axis (C2 vertebra). The odontoid process projects into the atlas (C1) and is held in place by the transverse ligament of the atlas, allowing rotation of the head (the atlantoaxial joint accounts for ~50\% of cervical rotation). Fractures of the odontoid process (Anderson--D'Alonzo classification types I--III) can cause atlantoaxial instability and spinal cord compression. Type II fractures (at the base of the dens) have the highest non-union rate. Os odontoideum is a congenital variant where the dens is not fused to the axis body. Also called the dens.

\medterm{Olfactory Nerve}-- Cranial nerve I (CN I), responsible for the sense of smell. Olfactory receptor neurons in the olfactory epithelium of the nasal roof project axons through the cribriform plate of the ethmoid bone to the olfactory bulb, then via the olfactory tract to the primary olfactory cortex (piriform cortex, amygdala, entorhinal cortex). Loss of smell (anosmia) can result from head trauma (shearing of fibers at the cribriform plate), viral infections (COVID-19, common cold), nasal polyps, neurodegenerative diseases (Parkinson, Alzheimer, Lewy body dementia), and meningiomas of the olfactory groove. Testing uses the University of Pennsylvania Smell Identification Test (UPSIT) or sniffin' sticks.

\medterm{Omentum}-- A double layer of peritoneum connecting the stomach to other abdominal organs. The greater omentum (gastrocolic omentum) hangs from the greater curvature of the stomach, covers the transverse colon and small intestine, and is rich in fat and immune cells (milky spots—aggregates of macrophages and lymphocytes). It is mobile and can migrate to sites of inflammation to wall off infection (the 'policeman of the abdomen'). The lesser omentum (gastrohepatic omentum) connects the lesser curvature of the stomach and duodenum to the liver, containing the hepatic artery, portal vein, and common bile duct (within the hepatoduodenal ligament). Omentectomy is performed in gastric and ovarian cancer surgery.

\medterm{Opposition}-- A complex, unique movement of the human hand in which the pulp of the thumb is brought into contact with the pad of any of the other four fingertips. It combines abduction, flexion, and medial rotation occurring primarily at the saddle-type first carpometacarpal (trapeziometacarpal) joint. It is produced primarily by the opponens pollicis muscle (assisted by flexor pollicis brevis and abductor pollicis brevis), innervated by the recurrent branch of the median nerve (C8--T1). Opposition is essential for fine precision grip and tool manipulation; loss of opposition occurs in carpal tunnel syndrome or median nerve laceration ('ape hand' deformity).

\medterm{Optic Nerve}-- CN II exiting through the optic disc (blind spot), passing through the optic canal, and partially decussating at the optic chiasm. Technically a CNS tract with oligodendrocyte myelination and meningeal coverings, making it susceptible to demyelinating disease (optic neuritis in MS). Papilledema (disc swelling from raised ICP) and optic atrophy are important findings. Carries approximately 1.2 million ganglion cell axons.

\medterm{Oral Cavity}-- The space extending from the lips to the oropharynx, bounded by the cheeks, palate, and tongue. It receives food, begins mechanical and chemical digestion, and contributes to speech.

\medterm{Organ}-- A structure composed of two or more different types of tissues that has a recognizable shape and performs specific functions.

\medterm{Organ System}-- A coordinated, higher-order anatomical and physiological functional grouping of distinct organs working synergistically to perform complex, life-sustaining bodily processes. The human body comprises 11 major organ systems: integumentary, skeletal, muscular, nervous, endocrine, cardiovascular, lymphatic/immune, respiratory, digestive, urinary, and reproductive. These systems are tightly interdependent, communicating via endocrine hormones, autonomic neurotransmission, and humoral signaling to maintain dynamic homeostatic equilibrium.

\medterm{Ossicles}-- Three smallest body bones forming the ossicular chain in the middle ear: malleus (hammer, attached to tympanic membrane), incus (anvil), stapes (stirrup, footplate in oval window). Amplify and transmit sound vibrations from tympanic membrane to inner ear, providing approximately 22-fold amplification (impedance matching between air and fluid). Tensor tympani (CN V3) attaches to the malleus; stapedius (CN VII) to the stapes, forming the acoustic reflex dampening loud sounds. Ossicular erosion or disruption (from chronic otitis media, cholesteatoma, or trauma) causes conductive hearing loss, often correctable with tympanoplasty or ossicular reconstruction.

\medterm{Osteoblast}-- A bone-forming cell derived from mesenchymal stem cells in the bone marrow. Osteoblasts synthesise and deposit the bone matrix (osteoid), composed primarily of type I collagen and non-collagenous proteins (osteocalcin, osteonectin, alkaline phosphatase). They then regulate matrix mineralisation by releasing matrix vesicles containing calcium and phosphate. Osteoblasts also produce RANKL (receptor activator of nuclear factor kappa-B ligand), which stimulates osteoclast differentiation, and osteoprotegerin (OPG), which inhibits RANKL. Once encased in mineralised matrix, they become osteocytes. Impaired osteoblast function contributes to osteoporosis (uncoupling of bone formation and resorption).

\medterm{Osteoclast}-- A large, multinucleated bone-resorbing cell derived from the monocyte--macrophage lineage (hematopoietic stem cells). Osteoclasts attach to bone surfaces via integrins, form a sealed acidic compartment (ruffled border), and secrete hydrogen ions (dissolve hydroxyapatite) and proteolytic enzymes (cathepsin K, matrix metalloproteinases) to degrade the bone matrix. They are regulated by RANK--RANKL signalling (stimulates differentiation/activity) and calcitonin (inhibits). Overactive osteoclasts cause bone loss in osteoporosis, Paget disease of bone, rheumatoid arthritis, and multiple myeloma. Bisphosphonates, denosumab (RANKL inhibitor), and calcitonin are therapeutic inhibitors of osteoclast activity.

\medterm{Osteocyte}-- The mature, terminally differentiated bone cell embedded within the bone matrix in lacunae, derived from osteoblasts. Osteocytes are connected to each other and to bone surface cells via canaliculi, forming a mechanosensory network that senses mechanical loading and fluid flow. They regulate bone remodeling by secreting signaling molecules (RANKL, osteoprotegerin, sclerostin) that coordinate osteoblast and osteoclast activity. Osteocyte apoptosis with aging reduces bone quality. Sclerostin (osteocyte product) inhibits bone formation; anti-sclerostin antibodies (romosozumab) treat osteoporosis.

\medterm{Ovary}-- The female gonad, a paired almond-shaped organ located in the pelvic cavity on each side of the uterus. Functions: production of oocytes (ova) and secretion of female sex hormones (oestrogen, progesterone, and small amounts of androgens). Each ovary contains follicles at various stages of development. Monthly ovulation releases a mature oocyte. After ovulation, the follicle becomes the corpus luteum, which secretes progesterone. Ovarian structure: outer cortex (containing follicles) and inner medulla (containing blood vessels and connective tissue). Common conditions: polycystic ovary syndrome, ovarian cysts, ovarian tumors, premature ovarian insufficiency.

\medterm{Pancreas}-- Retroperitoneal gland (15-20 cm) behind the stomach, spanning from duodenum to spleen. Exocrine function: acinar cells secrete digestive enzymes (amylase, lipase, proteases) and bicarbonate into pancreatic duct to duodenum. Endocrine function: islets of Langerhans secrete insulin (beta cells), glucagon (alpha cells), somatostatin (delta cells), pancreatic polypeptide. Blood supply: splenic artery, superior/inferior pancreaticoduodenal arteries. Conditions: acute/chronic pancreatitis, pancreatic cancer, diabetes mellitus, cystic fibrosis, pancreatic cysts.

\medterm{Pancreatic Duct}-- The main duct of the pancreas, conveying exocrine pancreatic secretions to the duodenum. It commonly joins the common bile duct before opening at the major duodenal papilla.

\medterm{Papillary Muscle}-- Muscular projections of the ventricular myocardium attached by chordae tendineae to the atrioventricular valve leaflets (anterolateral and posteromedial in the left ventricle; anterior, septal, and posterior in the right ventricle). They contract during ventricular systole to tense the chordae, preventing AV valve prolapse and regurgitation. Papillary muscle rupture (acute MI, especially posteromedial) causes acute severe mitral regurgitation, pulmonary edema, and cardiogenic shock—a surgical emergency.

\medterm{Parasympathetic Nervous System}-- The craniosacral division of the autonomic nervous system mediating 'rest-and-digest' functions. Preganglionic fibers from cranial nerves III, VII, IX, X, and sacral nerves S2-S4 synapse on ganglia near or within target organs (short postganglionic fibers). Neurotransmitter: acetylcholine at both ganglia (nicotinic) and end organs (muscarinic). Effects: bradycardia, bronchoconstriction, increased GI motility and secretion, miosis, bladder contraction, salivation, and penile erection. Opposes the sympathetic system; tonic balance maintains homeostasis.

\medterm{Parathyroid Glands}-- Four small, yellowish-brown, ovoid endocrine glands (each ~6 mm long, weighing 30--50 mg) typically situated on the posterior surface of the lateral lobes of the thyroid gland, divided into two superior parathyroids (derived from the fourth pharyngeal pouch) and two inferior parathyroids (derived from the third pharyngeal pouch, alongside the thymus, resulting in more variable ectopic locations). Composed predominantly of chief (principal) cells that synthesize and secrete parathyroid hormone (PTH) in response to low serum ionized calcium, and oxyphil cells that increase in number with age. PTH maintains systemic calcium homeostasis by stimulating osteoclast-mediated bone resorption, increasing renal distal tubular calcium reabsorption and inhibiting phosphate reabsorption, and activating renal 1-alpha-hydroxylase to convert vitamin D to active calcitriol (which enhances intestinal calcium absorption). Pathologies include primary hyperparathyroidism (most commonly due to a benign solitary parathyroid adenoma, presenting with 'stones, bones, groans, and psychiatric overtones'), secondary hyperparathyroidism in chronic kidney disease, and hypoparathyroidism (often iatrogenic following thyroidectomy), causing hypocalcemic tetany, Chvostek's sign, and Trousseau's sign.

\medterm{Parietal Bones}-- Paired curved bones forming the cranial vault roof and sides. Articulate at the sagittal suture, with frontal at coronal, temporal at squamous, and occipital at lambdoid sutures. The pterion (where four bones meet) overlies the middle meningeal artery; fractures here may cause epidural hematomas.

\medterm{Parietal Lobe}-- The cerebral lobe posterior to the central sulcus and above the lateral sulcus, containing the primary somatosensory cortex (postcentral gyrus, Brodmann 3,1,2) processing touch, pain, temperature, and proprioception, plus the somatosensory association cortex. Functions: spatial awareness, body schema, visuospatial processing, and (dominant hemisphere) skilled movements, apraxia; (nondominant) neglect. Lesions cause contralateral sensory loss, agraphesthesia, astereognosis, hemineglect (nondominant), Gerstmann syndrome (dominant: agraphia, acalculia, finger agnosia, left-right confusion), and homonymous hemianopia (optic radiation).

\medterm{Parotid Gland}-- The largest of the three major salivary glands, located in the preauricular region, extending from the zygomatic arch superiorly to the angle of the mandible inferiorly. The parotid gland is encased in the parotid fascia (investing layer of deep cervical fascia). It produces serous saliva (thin, watery, rich in amylase), secreted via Stensen duct (parotid duct), which pierces the buccinator muscle and opens opposite the second upper molar tooth. Structures passing through the parotid gland (from superficial to deep): facial nerve (CN VII—its five branches: temporal, zygomatic, buccal, marginal mandibular, and cervical), retromandibular vein, and external carotid artery (dividing into maxillary and superficial temporal arteries). Clinical: parotitis (mumps—most common cause), Sjögren syndrome, sialolithiasis (parotid duct stones), and parotid tumors (80\% benign—pleomorphic adenoma most common; 20\% malignant—mucoepidermoid carcinoma, adenoid cystic carcinoma). Parotidectomy is the surgical treatment; the facial nerve must be identified and preserved.

\medterm{Patella}-- Largest sesamoid bone, embedded in the quadriceps tendon. Articulates with the trochlear groove forming the patellofemoral joint. Increases quadriceps mechanical advantage by approximately 30 percent. Triangular with inferior apex. Covered with the thickest articular cartilage (up to 7 mm). Common conditions: fractures, dislocations, chondromalacia.

\medterm{Pectoralis Major Muscle}-- Thick fan-shaped chest muscle with clavicular and sternocostal heads, inserting on the lateral intertubercular groove lip. Innervated by medial/lateral pectoral nerves (C5-T1). Clavicular head flexes; sternocostal head extends from flexion. Adducts and medially rotates powerfully. Rupture uncommon, may occur during weightlifting at the musculotendinous junction.

\medterm{Pectoralis Minor Muscle}-- Thin triangular muscle deep to pectoralis major, from ribs 3-5 to the coracoid process. Innervated by medial pectoral nerve (C8-T1). Draws scapula anteriorly/inferiorly, stabilizing it against the thorax. Assists forced inspiration. Tightness contributes to rounded shoulder posture and thoracic outlet syndrome.

\medterm{Pelvic Cavity}-- The lower funnel-shaped portion of the abdominopelvic cavity enclosed by the bony pelvis (sacrum, coccyx, and paired hip bones), bounded superiorly by the pelvic brim (pelvic inlet) and inferiorly by the muscular pelvic diaphragm (levator ani and coccygeus muscles). It is distinguished from the greater (false) pelvis above the pelvic brim and is designated the lesser (true) pelvis. It houses the urinary bladder, terminal ureters, pelvic loop of the sigmoid colon, rectum, and the pelvic internal reproductive organs (uterus, fallopian tubes, and ovaries in females; prostate, seminal vesicles, and pelvic vas deferens in males). The peritoneum reflects over the pelvic viscera, creating important clinical dependency spaces: the rectovesical pouch in males, and the vesicouterine and rectouterine (pouch of Douglas) pouches in females. Pathologies include pelvic inflammatory disease, endometriosis, cul-de-sac fluid collections, urinary retention, pelvic organ prolapse, and pelvic fractures.

\medterm{Pelvis}-- Bony ring from 2 hip bones, sacrum, and coccyx. Each hip bone fuses from ilium, ischium, and pubis meeting at the acetabulum. Divided into false (greater, above brim) and true (lesser, below brim) pelvis. Males: narrower, deeper, heart-shaped inlet. Females: wider, shallower, oval inlet for childbirth.

\medterm{Penis}-- The male external genital organ, composed of three erectile columns: two dorsal corpora cavernosa and one ventral corpus spongiosum (containing the urethra, expanding into the glans). Root is attached to the pubic arch (crura and bulb). Erection: parasympathetic (pelvic splanchnic, NO-cGMP) vasodilation filling the corpora; detumescence by sympathetic. Blood supply: deep and dorsal arteries of the penis (from internal pudendal); drainage to superficial/deep dorsal veins. Conditions: phimosis, paraphimosis, Peyronie disease, priapism, penile fracture, carcinoma.

\medterm{Pericardial Cavity}-- The narrow, fluid-filled space between the layers of the pericardium surrounding the heart.

\medterm{Pericardium}-- The fibroserous sac surrounding the heart and proximal great vessels. Outer fibrous pericardium (anchors to diaphragm, sternum) and inner serous pericardium (parietal layer lining the sac; visceral layer/epicardium on the heart) with the pericardial cavity (15-50 mL serous fluid) between them. Functions: limits cardiac distension, reduces friction, barriers to infection. Pericarditis causes chest pain (relieved by leaning forward), friction rub, and ECG changes; pericardial effusion may cause cardiac tamponade (Beck triad). Pulsus paradoxus in tamponade.

\medterm{Periosteum}-- The dense connective tissue membrane covering the outer surface of bones (except at articular surfaces, which are covered by cartilage). Two layers: outer fibrous (collagen, vessels, nerves) and inner cambium (osteoprogenitor cells, osteoblasts). Supplies blood to bone, provides attachment for tendons/ligaments, and mediates appositional bone growth, callus formation, and fracture healing. Stripping of the periosteum (fracture, surgery) can devascularize the underlying bone. Pain-sensitive, mediating bone pain.

\medterm{Peripheral}-- A directional term describing structures situated away from the center or midline, toward the outer boundaries, extremities, or surface of the body or an organ. Opposed to central. Examples include the Peripheral Nervous System (cranial and spinal nerves), peripheral arterial disease in the limbs, and peripheral resistance vessels (arterioles).

\medterm{Peripheral Nervous System}-- The portion of the nervous system consisting of nerves and ganglia outside the brain and spinal cord. The PNS connects the central nervous system (CNS) to peripheral tissues and organs. Components: (1) Cranial nerves (12 pairs)—emerging from the brain/brainstem, innervating head, neck, and viscera. (2) Spinal nerves (31 pairs)—emerging from the spinal cord via dorsal and ventral roots, each dividing into dorsal and ventral rami. (3) Autonomic nervous system—sympathetic (thoracolumbar, 'fight or flight') and parasympathetic (craniosacral, 'rest and digest'), regulating involuntary functions. Peripheral nerves are composed of axons bundled into fascicles by connective tissue layers: endoneurium (surrounds individual axons), perineurium (surrounds fascicles—blood–nerve barrier), and epineurium (outer layer). Myelination (by Schwann cells in the PNS) enables saltatory conduction. Peripheral nerve injury is classified by Seddon (neuropraxia—conduction block, axonotmesis—axon disrupted but sheath intact, neurotmesis—complete transection) or Sunderland (grades I--V). Peripheral neuropathy (polyneuropathy, mononeuritis multiplex) has many causes.

\medterm{Peritoneum}-- The serous membrane lining the abdominal cavity and covering the abdominal organs. Parietal peritoneum lines the abdominal wall (innervated by somatic nerves—pain is well-localized); visceral peritoneum covers organs (innervated by visceral afferents—pain is dull, poorly localized). The peritoneal cavity (potential space) contains serous fluid. Mesenteries, omenta, and ligaments are peritoneal folds. Peritonitis (infection/inflammation) causes abdominal rigidity, rebound tenderness, and ileus. Peritoneal dialysis uses this membrane. Ascites is fluid accumulation in the peritoneal cavity.

\medterm{Phalanges}-- 14 finger bones. Thumb: 2 (proximal, distal). Each other finger: 3 (proximal, middle, distal). MCP joints are condyloid allowing flexion, extension, abduction, adduction, circumduction. IP joints are hinge joints allowing only flexion/extension. Mallet finger and jersey finger are common tendon avulsion injuries.

\medterm{Pharyngeal Tonsil}-- A collection of lymphoid tissue situated in the submucosa of the roof and posterior wall of the nasopharynx, forming the superior aspect of Waldeyer's ring. See \hyperlink{Adenoids}{Adenoids}.

\medterm{Pharynx}-- A muscular tube extending from the base of the skull (clivus) to the level of the cricoid cartilage (C6), where it continues as the esophagus. The pharynx is divided into three parts: nasopharynx (posterior to the nasal cavity, contains the pharyngeal tonsil/adenoids, Eustachian tube openings), oropharynx (posterior to the oral cavity, contains the palatine tonsils, lingual tonsil, base of tongue, soft palate), and laryngopharynx/hypopharynx (posterior to the larynx, extends from the hyoid bone to the cricoid cartilage, leads to esophagus and larynx). The pharynx functions in swallowing (deglutition—voluntary oral phase followed by involuntary pharyngeal phase, coordinated by the swallowing center in the medulla), breathing (conducts air from nasal/oral cavities to the larynx), and speech (resonance). The pharyngeal constrictor muscles (superior, middle, inferior) contract sequentially in peristalsis to propel the bolus. The pharynx is innervated by the pharyngeal plexus (CN IX and X).

\medterm{Phrenic Nerve}-- Primarily from C4 (with C3 and C5: "C3, 4, 5 keeps the diaphragm alive"), traveling between anterior and middle scalene muscles, descending anterior to the anterior scalene across the first rib. Right phrenic passes anterior to the lung root; left passes anterior and lateral. Innervates the ipsilateral diaphragm and provides sensory innervation to central diaphragmatic pleura, pericardium, and peritoneum. Injury causes ipsilateral hemidiaphragm paralysis reducing vital capacity. Phrenic nerve palsy causes ipsilateral hemidiaphragmatic paralysis, detected as an elevated hemidiaphragm on chest radiograph with paradoxical movement on fluoroscopic sniff test.

\medterm{Pia Mater}-- The innermost meningeal layer, a delicate vascularized membrane firmly adherent to the surface of the brain and spinal cord, following all gyri and sulci. It carries blood vessels that penetrate the CNS and forms the tela choroidea, contributing to the choroid plexus (CSF production). The pia and the spinal cord are anchored by the denticulate ligaments. Inflammation (meningitis) involves all meninges; subpial hemorrhages occur in some conditions.

\medterm{Pinna}-- Visible external ear portion, elastic cartilage covered by skin, shaped to funnel sound into the external auditory canal. Key features: helix (outer rim), antihelix (inner ridge), tragus (canal prominence), antitragus, concha (canal-leading depression), lobule (earlobe, only non-cartilaginous part). Blood supply: superficial temporal and posterior auricular arteries. Innervated by great auricular nerve (C2-C3) and auriculotemporal nerve (V3). Perichondritis can cause permanent deformity.

\medterm{Pisiform}-- Small pea-shaped sesamoid bone in the flexor carpi ulnaris tendon, in the proximal row. Articulates only with the triquetrum. Forms the medial border of Guyon's canal (ulnar nerve/artery passage). Only sesamoid carpal bone, subcutaneous and easily palpable on the ulnar palm.

\medterm{Pituitary Gland}-- Pea-sized "master gland" in the sella turcica, connected to hypothalamus by infundibulum. Anterior pituitary (adenohypophysis): GH, TSH, ACTH, FSH, LH, prolactin. Posterior pituitary (neurohypophysis): ADH (vasopressin), oxytocin (from hypothalamus). Regulates other endocrine glands and bodily functions. Conditions: pituitary adenoma, hypopituitarism, acromegaly, Cushing disease, diabetes insipidus, prolactinoma.

\medterm{Pleura}-- The serous membrane covering the lungs (visceral pleura) and lining the thoracic cavity (parietal pleura: costal, mediastinal, diaphragmatic, cervical/cupola). The pleural cavity is a potential space containing a thin film of serous fluid that couples the lung to the chest wall during ventilation. Parietal pleura is pain-sensitive (innervated by intercostal and phrenic nerves); visceral pleura is not. Pleural effusion, pneumothorax, hemothorax, empyema, and pleurisy (pleuritic chest pain) are pleural disorders. Pleural fluid analysis distinguishes transudates from exudates (Light criteria).

\medterm{Pleural Cavity}-- The narrow, fluid-filled space between the two layers of the pleura surrounding each lung.

\medterm{Plexus}-- A complex, interconnected branching network of intersecting nerves, blood vessels, or lymphatic channels. In the somatic nervous system, ventral rami of spinal nerves interweave to form major somatic plexuses (cervical C1--C4, brachial C5--T1, lumbar L1--L4, sacral L4--S4), ensuring that individual peripheral nerves receive fibers from multiple spinal levels so that single-root injury does not paralyze an entire limb muscle. Autonomic plexuses (cardiac, pulmonary, celiac, mesenteric, hypogastric) convey sympathetic and parasympathetic fibers to viscera. Vascular plexuses include the internal vertebral venous plexus and Kiesselbach's plexus in the anterior nasal septum (common site of epistaxis).

\medterm{Pons}-- Middle brainstem segment with prominent anterior bulge containing pontine nuclei (relaying cortical input to cerebellum), corticospinal fibers, and cranial nerve V, VI, VII, and VIII nuclei. Fourth ventricle floor is partly formed by the pons. Pneumotaxic center (parabrachial nucleus) regulates breathing pattern. Bilateral ventral pontine infarction causes locked-in syndrome with preserved consciousness and vertical eye movements but complete paralysis.

\medterm{Popliteal Artery}-- Continues from femoral artery at the adductor hiatus, dividing into anterior/posterior tibial arteries at the lower popliteus border. Deepest popliteal fossa structure, related anteriorly to vein and tibial nerve. Palpable at the fossa's inferior margin. Popliteal aneurysms are the most common peripheral arterial aneurysms, frequently bilateral and associated with AAA. Vulnerable to injury in knee dislocations.

\medterm{Popliteal Fossa}-- Diamond-shaped space on the posterior knee, bounded by hamstrings superiorly and gastrocnemius heads inferiorly. Contains popliteal artery/vein, tibial nerve, and common peroneal nerve. The popliteal pulse is palpated at the inferior margin. Significant for evaluating aneurysms, Baker's cysts, and vascular status in peripheral artery disease.

\medterm{Portal System}-- Venous network draining the GI tract, spleen, and pancreas to the liver via the portal vein (formed by superior mesenteric and splenic vein union behind the pancreatic neck). Carries nutrient-rich deoxygenated blood for hepatic processing before systemic entry. Receives left gastric, right gastric, cystic, and paraumbilical veins. Portal hypertension (cirrhosis) causes esophageal varices, caput medusae, and ascites.

\medterm{Portal Vein}-- The vein carrying nutrient-rich blood from the gastrointestinal tract, spleen, pancreas, and gallbladder to the liver, formed behind the pancreatic neck by the union of the superior mesenteric and splenic veins. It enters the porta hepatis with the hepatic artery and bile duct. Within the liver it ramifies, delivering blood to sinusoids for processing (first-pass metabolism). Portal hypertension (cirrhosis) leads to esophageal varices, ascites, splenomegaly, and portosystemic shunting; portal vein thrombosis causes portal hypertension without cirrhosis.

\medterm{Posterior}-- A directional term indicating a position situated at, toward, or closer to the back (dorsal) surface of the body or an organ. In human anatomy, posterior is generally synonymous with dorsal. Opposed to anterior (ventral). Examples include the esophagus lying posterior to the trachea, and the scapulae located on the posterior aspect of the thoracic wall.

\medterm{Posterior Cruciate Ligament (PCL)}-- A strong, thick, intra-articular, extrasynovial ligament of the knee joint that acts as the primary stabilizer against posterior tibial displacement. It originates from the anterolateral aspect of the medial femoral condyle in the intercondylar notch and passes posteriorly, laterally, and distally to insert onto the posterior intercondylar fossa of the tibia, approximately 1 cm below the articular line. Composed of a larger anterolateral bundle (tight in flexion) and a smaller posteromedial bundle (tight in extension). It provides approximately 95 percent of total resistance to posterior translation of the tibia relative to the femur and restrains hyperextension and rotation. PCL injuries typically result from high-energy direct anterior tibial trauma in a flexed knee ('dashboard injury' in motor vehicle collisions) or severe hyperextension. Evaluated clinically by the posterior drawer test, posterior sag sign (Godfrey's test), and quadriceps active test; confirmed by MRI.

\medterm{Posterior Tibial Artery}-- Arises from the popliteal artery, courses deep in the posterior leg, palpable behind the medial malleolus. Divides into medial/lateral plantar arteries for the sole. Supplies posterior leg muscles. Important for peripheral vascular disease assessment. The pedal pulse behind the medial malleolus is a key clinical examination component.

\medterm{Pronation}-- A rotational movement of the forearm occurring at the proximal and distal radioulnar joints, wherein the radius rotates medially across the anterior surface of the ulna, turning the palm of the hand to face posteriorly in the anatomical position (or downward when the elbow is flexed 90 degrees). It is produced primarily by the pronator teres and pronator quadratus muscles, innervated by the median nerve (C6--T1). In the foot, pronation describes a complex triplanar movement combining eversion, dorsiflexion, and abduction.

\medterm{Prostate}-- Walnut-sized gland (approximately 20 g) surrounding the urethra just below the bladder, producing approximately 30 percent of seminal fluid. It has peripheral, central, and transition zones. The peripheral zone is the most common site of prostate cancer; the transition zone enlarges in benign prostatic hyperplasia (BPH), compressing the urethra and causing urinary obstruction. The prostate secretes prostate-specific antigen (PSA), citric acid, and zinc. The ejaculatory ducts pass through the prostate and open at the verumontanum in the prostatic urethra.

\medterm{Proximal}-- A directional term used primarily in the extremities to describe a position nearer to the point of attachment of a limb to the trunk (shoulder or hip), or nearer to the source/origin of a vessel, nerve, or tubular organ. Opposed to distal. Examples include the elbow being proximal to the wrist, the femur being proximal to the tibia, and the proximal convoluted tubule being closest to the glomerulus.

\medterm{Proximal Radioulnar Joint}-- Pivot synovial joint between radial head circumference and ulnar radial notch, held by the annular ligament. Allows pronation/supination with approximately 150 degrees range. Parallel bones in supination; radius crosses ulna in pronation. Significant in radial head fractures, Monteggia fracture-dislocations, and nursemaid's elbow.

\medterm{Pubis}-- Anterior-inferior hip bone with body, superior ramus, and inferior ramus. Two pubic bones meet at the pubic symphysis (cartilaginous joint widening in childbirth). Pubic tubercle is the lateral inguinal ligament attachment. Direct hernias emerge medial to inferior epigastric vessels near it.

\medterm{Pulmonary Artery}-- The blood vessel that carries deoxygenated blood from the right ventricle of the heart to the lungs.

\medterm{Pulmonary Lobes}-- The anatomical divisions of the lungs separated by interlobar fissures. The right lung has superior, middle, and inferior lobes; the left lung has superior and inferior lobes and a lingula.

\medterm{Pulmonary Valve}-- Semilunar valve between right ventricle and pulmonary trunk with three cusps. Most superiorly located cardiac valve with thinner cusps than the aortic valve. Pulmonic stenosis is a congenital defect causing right ventricular hypertrophy. Less commonly affected by endocarditis than mitral/aortic valves. A systolic ejection murmur at the left upper sternal border suggests pulmonic stenosis.

\medterm{Pulmonary Vein}-- The four veins (two from each lung) returning oxygenated blood from the lungs to the left atrium, a unique exception in the venous system (carrying oxygenated blood). Superior and inferior pulmonary veins enter the left atrium posteriorly. Anomalous drainage (partial or total anomalous pulmonary venous return) causes left-to-right shunt/cyanosis. Pulmonary vein isolation (ablation) is a treatment for atrial fibrillation because ectopic foci in the pulmonary vein sleeves trigger AF. Pulmonary vein stenosis complicates ablation.

\medterm{Pupil}-- The central circular aperture in the iris that regulates the amount of light entering the posterior segment of the eye. Pupillary diameter is controlled by two smooth muscle groups: the circumferentially arranged sphincter pupillae (innervated by parasympathetic postganglionic fibers from the ciliary ganglion via CN III, producing constriction or miosis down to 2--3 mm) and the radially arranged dilator pupillae (innervated by sympathetic postganglionic fibers from the superior cervical ganglion, producing dilation or mydriasis up to 8 mm). The pupillary light reflex tests the integrity of CN II (optic nerve afferent limb to the pretectal nucleus) and CN III (oculomotor efferent limb via the Edinger-Westphal nuclei); illumination of one eye elicits bilateral pupillary constriction (direct and consensual responses). Clinical abnormalities include the Argyll Robertson pupil (bilateral, small, irregular pupils that accommodate to near vision but do not react to light, classic for neurosyphilis), Horner syndrome (miosis, ptosis, and anhidrosis from sympathetic pathway disruption), and a unilaterally fixed, dilated pupil indicating uncal herniation compressing the ipsilateral oculomotor nerve.

\medterm{Quadriceps Femoris Muscle}-- Four-muscle anterior thigh group: rectus femoris, vastus lateralis, vastus medialis, vastus intermedius, all inserting on the tibial tuberosity via the patellar ligament. Innervated by femoral nerve (L2-L4). Primary knee extensor for walking, running, rising from seated. VMO is important for patellar tracking. Weakness common after knee injury, surgery, and aging.

\medterm{Radial Artery}-- The smaller terminal branch of the brachial artery, arising at the cubital fossa, running laterally along the radius to the wrist (palpable at the anatomical snuffbox) and hand. Supplies the lateral forearm and the thumb side of the hand; forms the deep palmar arch. Clinical importance: radial artery pulse, Allen test (collateral circulation before radial artery cannulation or CABG conduit harvest), arterial blood gas sampling, and coronary angiography access. The radial artery is a preferred conduit for coronary bypass grafting.

\medterm{Radial Nerve}-- The largest terminal branch of the posterior cord of the brachial plexus, carrying fibers from spinal roots C5--T1. In the axilla and arm, it travels posterior to the brachial artery, passes through the lower triangular space, and descends obliquely along the spiral (radial) groove of the humerus accompanied by the profunda brachii artery, giving motor branches to the triceps brachii, anconeus, brachioradialis, and extensor carpi radialis longus. Anterior to the lateral epicondyle of the humerus, it bifurcates into a superficial sensory branch (supplying skin of the radial two-thirds of the dorsal hand and dorsal proximal phalanges of the lateral 3.5 digits) and a deep motor branch (posterior interosseous nerve, supplying the extensor-supinator compartment of the forearm). A midshaft humeral shaft fracture or prolonged compression ('Saturday night palsy') damages the nerve in the radial groove, causing wrist drop, finger drop, and sensory loss over the first dorsal web space; although the flexor muscles are intact, functional grip strength is severely diminished due to active insufficiency of the wrist and finger flexors unless the wrist is passively extended.

\medterm{Radius}-- Lateral forearm bone in the anatomical position. Proximal end: head, neck, radial tuberosity. Shaft. Distal end: styloid process, articular surfaces. The radial head articulates with the capitulum and ulnar notch, allowing pronation/supination. Bears approximately 80 percent of wrist axial load. Colles' fracture (distal radius with dorsal displacement) is among the most common fractures.

\medterm{Rectum}-- The terminal approximately 12--15 cm of the large intestine between the sigmoid colon and anal canal. It stores feces and participates in defecation.

\medterm{Rectus Abdominis Muscle}-- Long flat vertical anterior abdominal wall muscle from pubic symphysis to xiphoid process/costal cartilages 5-7, divided by tendinous intersections ("six-pack"). Innervated by T7-T12. Flexes trunk, compresses abdomen, increases intra-abdominal pressure. Diastasis recti common in pregnancy/obesity. Common TRAM flap donor site for breast reconstruction.

\medterm{Rectus Abdominis Sheath}-- A fibrous aponeurotic envelope formed by the abdominal wall muscles and surrounding the rectus abdominis muscle. It contains the superior and inferior epigastric vessels and related nerves.

\medterm{Reflex Arc}-- The neural pathway that mediates a reflex action, bypassing conscious brain processing for rapid, automatic responses. A reflex arc consists of five components: (1) Receptor—sensory nerve ending detecting a stimulus (e.g., pain, stretch). (2) Sensory (afferent) neuron—carries the signal from the receptor to the spinal cord (dorsal root ganglion). (3) Integration center—a single synapse (monosynaptic reflex, e.g., patellar reflex) or multiple interneurons (polysynaptic reflex, e.g., withdrawal reflex) within the spinal cord gray matter. (4) Motor (efferent) neuron—carries the signal from the spinal cord to the effector organ. (5) Effector—muscle or gland that executes the response. Examples: monosynaptic—stretch reflex (patellar reflex, testing L3--L4 nerve root). Polysynaptic—flexor withdrawal reflex (pain causes limb withdrawal), crossed extensor reflex (opposite limb extends to maintain posture), and superficial reflexes (abdominal, cremasteric, plantar—Babinski sign). The reflex arc is the functional unit of the nervous system and is assessed in the neurological examination to localise lesions.

\medterm{Renal Artery}-- A paired branch of the abdominal aorta supplying a kidney. It divides into segmental, interlobar, arcuate, and interlobular arteries within the renal parenchyma.

\medterm{Renal Cortex}-- Outer kidney portion between capsule and medulla. Granular tissue with glomeruli, convoluted tubules, and cortical collecting ducts. Extends between pyramids as renal columns. Receives approximately 90 percent of renal blood flow; primary filtration site. Cortical nephrons (85 percent) have cortical glomeruli; juxtamedullary (15 percent) are near the corticomedullary junction. Essential for GFR regulation and erythropoietin production.

\medterm{Renal Medulla}-- Inner kidney with 8-18 renal pyramids containing loops of Henle, vasa recta, and medullary collecting ducts in countercurrent arrangement for urine concentration. Pyramid bases face cortex; apices (papillae) project into minor calyces. Lower blood flow (approximately 10 percent) and higher osmolality (up to 1200 mOsm/kg at papillary tip). Corticomedullary gradient is critical for ADH-mediated water reabsorption. More susceptible to ischemic injury than cortex.

\medterm{Renal Pelvis}-- The funnel-shaped expansion of the ureter within the kidney, formed by the union of major calyces. It collects urine from the renal papillae and directs it into the ureter.

\medterm{Renal Vein}-- A paired vein draining each kidney into the inferior vena cava. The left renal vein is longer and crosses anterior to the aorta, receiving the left gonadal and suprarenal veins.

\medterm{Reticular Formation}-- A network of interconnected nuclei in the brainstem (midbrain, pons, medulla) that regulates arousal, consciousness, sleep-wake transitions, and alertness (ascending reticular activating system), as well as muscle tone, posture, cardiovascular and respiratory control (vasomotor and respiratory centers), pain modulation, and attention. Damage to the ARAS causes coma or reduced consciousness; certain sedatives and anesthetics act here. Collateral projections from sensory pathways activate the ARAS.

\medterm{Retina}-- Innermost eye layer, a thin neural tissue layer lining the posterior two-thirds of the globe for phototransduction (light to neural signals). Ten layers, with photoreceptors (rods for low-light vision, cones for color vision) in the outer nuclear layer. Macula lutea is the central high-acuity area; fovea centralis has the highest cone density and sharpest vision. Supplied by the central retinal artery; drained by the central retinal vein. Conditions: retinal detachment, diabetic retinopathy, age-related macular degeneration (ARMD), retinal vein occlusion, and retinitis pigmentosa.

\medterm{Retroperitoneum}-- Space posterior to the peritoneum. Primary retroperitoneal (never intraperitoneal): kidneys, adrenals, ureters, aorta, IVC. Secondary (originally intraperitoneal): duodenum 2nd-4th parts, pancreas head/body, ascending/descending colon. Divided into anterior pararenal, perirenal, and posterior pararenal spaces. Contains retroperitoneal fibrosis, hematomas, and tumors.

\medterm{Rib}-- One of 12 paired, curved bones forming the thoracic cage, protecting the heart, lungs, and great vessels. True ribs (1--7) attach directly to the sternum via costal cartilage. False ribs (8--10) attach indirectly to the sternum via the costal cartilage of rib 7. Floating ribs (11--12) have no anterior attachment. Each rib has a head, neck, tubercle, shaft, and costal groove (housing the intercostal neurovascular bundle). Intercostal muscles occupy the spaces between ribs and are essential for respiration. Rib fractures are common and can cause pneumothorax or haemothorax.

\medterm{Right Atrium}-- Receives deoxygenated blood from superior vena cava, inferior vena cava, and coronary sinus. Smooth posterior part (sinus venosus) and rough anterior auricle with pectinate muscles. The fossa ovalis is the fetal foramen ovale remnant. The SA node pacemaker is near the superior vena cava opening. Contracts to fill the right ventricle through the tricuspid valve.

\medterm{Right Hypochondriac Region}-- One of nine abdominal regions in the upper right abdomen, bounded medially by the right midclavicular line and superiorly by the right costal margin. Contains the cecum, appendix, ascending colon, and inferior right liver lobe. Pain here may indicate appendicitis (McBurney's point), cecal pathology, or hepatic disease. Important for assessing hepatomegaly and gallbladder referred pain.

\medterm{Right Iliac Region}-- Lower right abdomen containing the cecum, appendix, and terminal ileum. McBurney's point (junction of lateral and middle thirds from umbilicus to ASIS) is critical for appendicitis diagnosis. Also significant for evaluating inguinal and femoral hernias, and the right ureter's course.

\medterm{Right Lymphatic Duct}-- Short lymphatic vessel (approximately 1.5 cm) draining lymph from the right upper quadrant: right arm, right side of the head and neck, and right thorax. Formed by the junction of right jugular, subclavian, and bronchomediastinal trunks, emptying into the right venous angle (junction of right internal jugular and subclavian veins). Blockage may cause right-sided chylothorax, though this is less common than left-sided from thoracic duct injury.

\medterm{Right Ventricle}-- Receives deoxygenated blood from the right atrium through the tricuspid valve, pumping into the pulmonary trunk via the pulmonary valve. Thinner walls (3-5 mm) than the left ventricle due to low pulmonary resistance. Features: trabeculae carneae, 3 papillary muscles with tricuspid chordae tendineae, moderator band with right bundle branch. Pressure overload from pulmonary hypertension causes hypertrophy and failure.

\medterm{Rostral}-- A directional term relating to or situated toward the snout, rostrum, beak, or anterior end of the head. In neuroanatomy, because the human neuraxis bends approximately 80 degrees at the cephalic flexure, rostral means toward the frontal pole/forehead within the cerebrum, and toward the upper brainstem/cranial pole within the spinal cord and brainstem. Opposed to caudal.

\medterm{Rotator Cuff}-- A group of four muscles and their tendons that stabilize the glenohumeral (shoulder) joint and allow a wide range of motion. The muscles are the supraspinatus (abduction), infraspinatus (external rotation), teres minor (external rotation), and subscapularis (internal rotation)—collectively remembered by the acronym SITS. Rotator cuff tears and impingement syndrome are common causes of shoulder pain and weakness.

\medterm{Sacral Vertebrae (S1-S5)}-- Five vertebrae fused to form the sacrum. They transmit body weight from the vertebral column to the pelvic girdle and contain anterior and posterior sacral foramina for spinal nerves.

\medterm{Sacroiliac Joint}-- The synovial joint between the sacrum and the ilium (auricular surfaces), forming the connection between the axial skeleton and the pelvis. The sacroiliac joint is reinforced by strong ligaments (anterior, interosseous, and posterior sacroiliac ligaments, sacrotuberous and sacrospinous ligaments). It has limited mobility (2--4 degrees of rotation, 1--2 mm of translation) and functions primarily to transmit weight from the upper body to the lower limbs and provide stability. Sacroiliac joint dysfunction (sacroiliitis) is a common cause of low back pain, often referred to the buttock, groin, and posterior thigh. Causes: inflammatory (ankylosing spondylitis, psoriatic arthritis, reactive arthritis, inflammatory bowel disease-associated arthropathy), degenerative (osteoarthritis), pregnancy and postpartum (hormonally-mediated ligamentous laxity), trauma, and infection (septic arthritis, rare). Diagnosis: clinical examination (FABER test, Gaenslen sign, compression/distraction), imaging (X-ray, CT, MRI with STIR sequences for inflammation), and diagnostic intra-articular injection. Treatment: physiotherapy (core strengthening, pelvic stabilisation), NSAIDs, local corticosteroid injection, and radiofrequency ablation for refractory pain.

\medterm{Sacrum}-- Large triangular bone from 5 fused sacral vertebrae between the hip bones. Articulates with L5 superiorly and coccyx inferiorly. The sacral promontory is an obstetric landmark. Contains 4 pairs of sacral foramina. Forms the posterior pelvic wall. Sacroiliac joints are important for force transmission between spine and lower limbs.

\medterm{Sagittal Plane}-- A vertical anatomical plane oriented anteroposteriorly, parallel to the sagittal suture of the skull, that divides the body or an organ into right and left portions. The median (midsagittal) plane passes precisely through the midline of the body, dividing it into symmetrical right and left halves. Parasagittal planes run parallel to the median plane on either side. Movements occurring in the sagittal plane include flexion and extension. Sagittal MRI views are essential for neuroimaging the brainstem, corpus callosum, and spinal cord.

\medterm{Salivary Glands}-- Exocrine glands that produce saliva. Three pairs of major salivary glands: parotid glands (largest, located preauricularly, secrete serous saliva via Stensen duct), submandibular glands (located in the submandibular triangle, secrete mixed serous and mucous saliva via Wharton duct), and sublingual glands (smallest, located in the floor of the mouth, secrete predominantly mucous saliva via multiple ducts of Rivinus). Numerous minor salivary glands (600--1000) are scattered throughout the oral mucosa, palate, pharynx, and lips. Saliva (1--1.5 L/day) contains water (99\%), electrolytes, amylase (starch digestion), mucus (lubrication), lysozyme and IgA (antimicrobial), and lingual lipase (fat digestion). Salivary secretion is controlled by the autonomic nervous system: parasympathetic (CN VII and IX—profuse, watery secretion) and sympathetic (scanty, viscous secretion).

\medterm{Scaphoid}-- Largest proximal carpal row bone on the radial wrist side. Most commonly fractured carpal bone from falls on an outstretched hand. Unique distal blood supply makes waist fractures prone to proximal fragment avascular necrosis. Anatomical snuffbox tenderness is characteristic. Initial radiographs may be negative, requiring repeat imaging or MRI.

\medterm{Scapula}-- Flat triangular bone on the posterior thorax between ribs 2-7. Key features: acromion, coracoid process, glenoid cavity, spine, supraspinous/infraspinous fossae. Stabilized by 17 muscles and serves as rotator cuff attachment. Has 3 borders, 3 angles, and 2 surfaces. Fractures are uncommon, usually from high-energy trauma.

\medterm{Sciatic Nerve}-- Body's largest nerve (L4-S3), formed from the lumbosacral plexus, exiting via the greater sciatic foramen below the piriformis. Descends in the posterior thigh, dividing into tibial and common peroneal nerves at the popliteal fossa. Innervates hamstrings and provides lower leg/foot motor/sensory innervation via its branches. Sciatica is pain radiating along its distribution, commonly from L4-L5 or L5-S1 disc herniation.

\medterm{Sclera}-- Dense fibrous opaque outer eye coat providing structural integrity and extraocular muscle attachment. White, covering approximately five-sixths of the globe. Continuous with cornea anteriorly (limbus) and dura mater posteriorly through the optic nerve sheath. May appear yellowish in elderly or with jaundice (icterus). Scleritis (inflammation) is associated with autoimmune diseases like rheumatoid arthritis and may be sight-threatening.

\medterm{Scrotum}-- Pouch of skin and superficial fascia containing the testes, epididymides, and lower spermatic cords. The scrotal wall has multiple layers reflecting the abdominal wall layers. The dartos muscle (smooth muscle in the superficial fascia) and cremaster muscle (skeletal muscle from the internal oblique) regulate testicular temperature by contracting (bringing testes closer to the body in cold) or relaxing (lowering them in warmth). The scrotum is divided by a median septum into two compartments, each containing one testis with its epididymis and the lower portion of the spermatic cord.

\medterm{Sebaceous Glands}-- Holocrine glands associated with hair follicles (pilosebaceous units) secreting sebum (triglycerides, wax esters, squalene, cholesterol). Lubricate and waterproof skin/hair, have antimicrobial properties (fatty acids), and maintain skin barrier function. Most numerous on the face and scalp. Hyperactivity contributes to acne vulgaris; hypoactivity causes dry skin. Stimulated by androgens (explaining puberty acne). Fordyce spots are ectopic glands on oral mucosa and genitalia.

\medterm{Semicircular Canals}-- Three fluid-filled (endolymph) loops in each inner ear, oriented approximately orthogonally to detect angular (rotational) acceleration in three planes: horizontal (lateral), anterior (superior), and posterior. Each has an ampulla with a crista containing hair cells embedded in a gelatinous cupula. Head rotation deflects the cupula, bending stereocilia and stimulating hair cells. Horizontal canal excitatory direction is ampullopetal flow. BPPV results from displaced otoconia entering the canals, causing benign paroxysmal positional vertigo with brief episodes of spinning provoked by head movement.

\medterm{Seminal Vesicles}-- Paired lobulated glands behind the bladder and anterior to the rectum, each approximately 5 cm long, producing approximately 60-70 percent of seminal volume. The secretion is alkaline (neutralizing vaginal acidity), rich in fructose (energy source for sperm), prostaglandins (stimulating female reproductive tract contractions), and clotting proteins (causing semen coagulation after ejaculation). Each seminal vesicle duct joins the vas deferens to form the ejaculatory duct. Seminal vesicle pathology (infection, cysts, or invasion by prostate cancer) may affect ejaculate volume, cause haemospermia, and impair male fertility.

\medterm{Serratus Anterior Muscle}-- Fan-shaped muscle from ribs 1-8/9 to the medial scapular border. Innervated by long thoracic nerve (C5-C7). Protracts scapula, upwardly rotates glenoid for overhead elevation, stabilizes against thorax. Long thoracic nerve injury causes medial scapular winging, especially visible during wall push-ups. Commonly injured in breast surgery.

\medterm{Sigmoid Colon}-- S-shaped segment from descending colon to rectum in the left iliac fossa, approximately 40-50 cm, with complete mesentery (intraperitoneal, mobile). Most common diverticular disease site from narrow lumen and high pressure. Common volvulus site, particularly in elderly/institutionalized patients. Second most common colorectal cancer site. Transition between water-absorbing colon and rectal storage.

\medterm{Sinuses}-- Air-filled cavities within the skull bones lined by respiratory mucosa, connected to the nasal cavity via ostia. Four paired paranasal sinuses: maxillary (largest, in maxilla, drains to middle meatus), frontal (in frontal bone, drains to middle meatus), ethmoid (3-18 cells in ethmoid bone, anterior/middle/posterior groups), sphenoid (in sphenoid bone, drains to sphenoethmoidal recess). Functions: lighten skull, resonate voice, humidify air, crumple zones for trauma. Conditions: sinusitis (acute/chronic), mucocele, polyps, fungal sinusitis, osteomeatal complex obstruction.

\medterm{Skeletal Muscle}-- Voluntary, striated muscle attached to bones by tendons or directly by connective tissue. Its contraction produces movement, maintains posture, and generates heat.

\medterm{Skin}-- The largest organ of the body (approximately 1.5--2 m\textsuperscript{2}), forming a protective barrier against physical injury, pathogens, UV radiation, and water loss. Layers: epidermis (stratified squamous epithelium, keratinocytes, melanocytes, Langerhans cells, Merkel cells), dermis (connective tissue with collagen and elastin fibers, blood vessels, nerves, hair follicles, sweat and sebaceous glands), and hypodermis (subcutaneous fat). Functions: thermoregulation (sweating, vasodilation/vasoconstriction), sensation (touch, pain, temperature), vitamin D synthesis, immune surveillance. Common diseases: dermatitis, psoriasis, skin cancer (basal cell, squamous cell, melanoma), acne, infections.

\medterm{Skull}-- Bony framework of the head with 22 bones: cranium (8: frontal, 2 parietal, 2 temporal, occipital, sphenoid, ethmoid) and facial (14). Joined by sutures: coronal, sagittal, lambdoid, squamous. The skull base has foramina for cranial nerves and vessels. Protects the brain and provides attachment for muscles of mastication and facial expression.

\medterm{Small Intestine}-- Longest GI segment (approximately 6 meters) from pylorus to ileocecal valve. Primary nutrient digestion and absorption site. Three segments: duodenum (25 cm, chyme/bile/pancreatic enzymes), jejunum (2.5 m, prominent folds/villi for absorption), ileum (3.5 m, Peyer's patches for immunity). Specialized for absorption via plicae circulares, villi, and microvilli. Ileocecal valve prevents colonic backflow.

\medterm{Smooth Muscle}-- Involuntary, non-striated muscle found in the walls of hollow organs and blood vessels. It produces slow, sustained contractions under autonomic, hormonal, and local control.

\medterm{Soleus Muscle}-- Broad flat deep posterior leg muscle from posterior tibia, fibula, and tendinous arch to the calcaneus via the Achilles tendon. Innervated by tibial nerve (S1-S2). Powerful ankle plantarflexor, the primary anti-gravity muscle for upright posture. One-joint muscle. High slow-twitch fiber proportion for fatigue-resistant sustained postural activities and walking.

\medterm{Sphenoid Bone}-- Butterfly-shaped bone at the skull base, with body (sphenoid sinuses), greater/lesser wings, and pterygoid processes. The sella turcica houses the pituitary gland. Important foramina include the optic canal (CN II, ophthalmic artery), superior orbital fissure (CN III, IV, V1, VI), foramen rotundum (V2), ovale (V3), and spinosum (middle meningeal artery). Articulates with most cranial bones.

\medterm{Spinal Cord}-- The cylindrical bundle of nerve fibers extending from the medulla oblongata (at the foramen magnum) to the conus medullaris (L1--L2 in adults), surrounded by the vertebral column and meninges (dura mater, arachnoid mater, pia mater). The spinal cord is divided into 31 segments: 8 cervical, 12 thoracic, 5 lumbar, 5 sacral, and 1 coccygeal. Each segment gives rise to a pair of spinal nerves (dorsal root—sensory, ventral root—motor). Internal structure: central gray matter (butterfly-shaped—dorsal horn: sensory; ventral horn: motor; lateral horn: autonomic) surrounded by white matter (ascending tracts—spinothalamic, dorsal columns, spinocerebellar; descending tracts—corticospinal, reticulospinal, vestibulospinal). The spinal cord conducts sensory information from the periphery to the brain, carries motor commands from the brain to muscles, and mediates spinal reflexes (stretch, withdrawal). Spinal cord injury causes paralysis, sensory loss, and autonomic dysfunction below the level of the lesion.

\medterm{Spinal Nerve}-- A mixed nerve formed by the union of dorsal and ventral roots as it exits the vertebral canal. There are 31 pairs, carrying sensory, motor, and autonomic fibers to and from the spinal cord.

\medterm{Spine}-- A central anatomical reference term used clinically for the vertebral column (the bony column of 33 vertebrae protecting the neural canal) or the spinal cord (the primary cylindrical neural conduit transmitting motor and sensory pathways between the brain and periphery). See \hyperlink{Spinal Cord}{Spinal Cord} and \hyperlink{Vertebral Column}{Vertebral Column}.

\medterm{Spleen}-- Largest lymphoid organ in the left upper quadrant under ribs 9-11, approximately 12 cm long and 150-200 g. Red pulp filters blood (removing old/damaged RBCs and platelets); white pulp mounts immune responses (lymphoid follicles and periarteriolar lymphoid sheaths). Enclosed by a fibroelastic capsule with trabeculae. Connected to the stomach by the gastrosplenic ligament. The splenic hilum is where vessels enter and exit. Asplenia (congenital or surgical) increases infection risk from encapsulated bacteria, particularly Streptococcus pneumoniae, Neisseria meningitidis, and Haemophilus influenzae.

\medterm{Spleen Hilum}-- Concave medial surface where the tortuous splenic artery (celiac trunk branch) enters and splenic vein exits, with lymphatics and nerves. Artery divides into 5 segmental branches at the hilum (basis for partial splenectomy). Splenic vein joins the superior mesenteric vein forming the portal vein. Attaches the splenorenal ligament containing vessels and pancreatic tail. Pancreatic pathology may involve splenic vessels here.

\medterm{Sternoclavicular Joint}-- A saddle-type synovial joint between the clavicle and manubrium of the sternum. It is the only bony connection between the upper limb and axial skeleton and permits movement of the shoulder girdle.

\medterm{Sternocleidomastoid Muscle}-- Strap-like anterior neck muscle from manubrium and medial clavicle to the mastoid process. Innervated by CN XI and C2-C3. Divides neck into anterior/posterior triangles. Unilaterally: laterally flexes to opposite side, rotates face to same side. Bilaterally: flexes neck, assists forced inspiration. Key examination landmark; commonly affected by congenital torticollis.

\medterm{Sternum}-- Flat bone in the anterior midline thorax, with manubrium, body, and xiphoid process. The sternal angle (of Louis) at the manubriosternal junction is a key landmark at the 2nd costal cartilage, T4-T5, and tracheal bifurcation. Common site for bone marrow biopsy. Median sternotomy is the standard cardiac surgical approach.

\medterm{Stomach}-- J-shaped muscular organ in the left upper quadrant, connecting esophagus to duodenum. Regions: cardia, fundus, body, antrum, pylorus. Functions: food storage, mechanical churning, chemical digestion (pepsin, HCl), intrinsic factor secretion (B12 absorption). Mucosa has rugae (folds) allowing expansion. Blood supply: celiac trunk branches (left/right gastric, left/right gastroepiploic, short gastric). Venous drainage to portal system. Nerve supply: vagus (parasympathetic), celiac plexus (sympathetic). Conditions: gastritis, peptic ulcer, gastric cancer, gastroparesis, obstruction.

\medterm{Subclavian Artery}-- The artery supplying the upper limb, arising on the right from the brachiocephalic trunk and on the left directly from the aortic arch, passing over the first rib and becoming the axillary artery. Branches: vertebral, internal thoracic, thyrocervical trunk, costocervical trunk. The subclavian is a site of central venous/arterial access (subclavian line) and is compressed in thoracic outlet syndrome and by cervical ribs. Occlusion causes upper limb ischemia; subclavian steal syndrome (vertebral reversal) causes vertebrobasilar symptoms with arm exercise.

\medterm{Subclavian Vein}-- The continuation of the axillary vein beyond the lateral border of the first rib. It joins the internal jugular vein behind the sternoclavicular joint to form the brachiocephalic vein.

\medterm{Submandibular Gland}-- One of the three paired major salivary glands, located in the submandibular triangle beneath the mandible. Mixed gland (predominantly serous) secreting saliva via Wharton's duct opening at the sublingual caruncle (floor of mouth). Saliva contains amylase, mucins, lysozyme, IgA. Submandibular stones (sialolithiasis) are common (Wharton's duct is long and gravity-dependent); causes painful swelling with eating. Enlargement on palpation bimanually. Blood supply: facial and lingual arteries; innervation: parasympathetic via chorda tympani/submandibular ganglion.

\medterm{Superficial}-- A directional term indicating a location situated closer to the external surface of the body or closer to the exterior of an organ. Opposed to deep (profundus). Examples include the epidermis being superficial to the dermis, and the external oblique muscle being superficial to the internal oblique.

\medterm{Superficial Fascia}-- A layer of loose connective tissue beneath the dermis, also called subcutaneous tissue or the hypodermis. It contains variable amounts of adipose tissue, superficial blood vessels, lymphatics, and cutaneous nerves; it cushions and insulates the body, stores energy, permits skin movement over deeper structures, and provides a route for subcutaneous injections.

\medterm{Superior}-- A directional term indicating a position situated nearer to the head (vertex), toward the upper part of the body, or higher relative to another structure. Synonymous with cranial or cephalic. Opposed to inferior (caudal). Examples include the heart lying superior to the diaphragm, and the eyebrows being superior to the eyes.

\medterm{Superior Mesenteric Artery}-- Second major abdominal aortic branch at L1, the midgut artery. Supplies distal duodenum through proximal two-thirds of transverse colon. Branches: inferior pancreaticoduodenal, jejunal, ileal, ileocolic, right colic, middle colic. Arises anterior to the pancreatic uncinate process, crosses the third duodenal part. Occlusion causes acute mesenteric ischemia with severe pain disproportionate to examination. SMA syndrome compresses the duodenum.

\medterm{Superior Vena Cava}-- Formed by brachiocephalic vein junction at the first right costal cartilage, returning upper body deoxygenated blood to the right atrium. Approximately 7 cm, valveless. Related to right lung/pleura, right phrenic nerve, and azygos vein. SVC syndrome from compression (commonly lung cancer) causes facial swelling, plethora, and upper extremity edema.

\medterm{Supination}-- A rotational movement of the forearm occurring at the proximal and distal radioulnar joints, wherein the radius rotates laterally to become parallel with the ulna, turning the palm to face anteriorly in the anatomical position (or upward when the elbow is flexed 90 degrees). It is produced by the biceps brachii (most powerful supinator, especially when the elbow is flexed; musculocutaneous nerve) and the supinator muscle (radial nerve/PIN). In the foot, supination combines inversion, plantarflexion, and adduction.

\medterm{Sweat Glands}-- Exocrine glands that secrete sweat onto the skin surface or into hair follicles. Eccrine glands are widely distributed and produce a watery secretion for thermoregulation; their ducts open directly onto the surface and are activated mainly by sympathetic cholinergic fibers. Apocrine glands are concentrated in the axillae and anogenital region, empty into hair follicles, and produce a protein- and lipid-rich secretion that develops odor after bacterial breakdown, especially after puberty.

\medterm{Sympathetic Nervous System}-- The division of the autonomic nervous system responsible for the "fight-or-flight" response, accelerating heart rate, dilating bronchioles, inhibiting digestion, and mobilizing energy stores during stress or emergency. Anatomically, it arises from the thoracolumbar spinal cord (T1-L2) and utilizes the sympathetic trunk (paravertebral ganglia) and prevertebral ganglia to distribute postganglionic fibers. It acts in opposition to the parasympathetic nervous system.

\medterm{Synovium}-- The inner layer of the joint capsule in synovial joints, composed of synovial lining cells (type A macrophage-like, type B fibroblast-like) on a vascularized connective tissue base. Produces synovial fluid (hyaluronic acid, lubricin, proteinases) for lubrication and chondrocyte nutrition. Two layers: intimal (1-3 cells thick) and subintimal (vascularized connective tissue). Inflammation (synovitis) causes joint swelling, pain, and destruction in rheumatoid arthritis, gout, infection.

\medterm{Talus}-- Second largest tarsal bone articulating with tibia/fibula proximally at the ankle. Approximately 60 percent of its surface is articular cartilage with no muscular attachments. Transmits entire body weight to the calcaneus. Talar neck fractures may disrupt blood supply causing avascular necrosis in up to 50 percent. Also participates in the subtalar joint for inversion/eversion.

\medterm{Tarsus}-- The collection of seven bones forming the posterior part of the foot, between the tibia/fibula and the metatarsals. The tarsal bones: talus (articulates with tibia and fibula at the ankle joint), calcaneus (heel bone—largest tarsal bone, insertion of Achilles tendon), navicular (medial midfoot), cuboid (lateral midfoot), and three cuneiforms (medial, intermediate, lateral—articulate with the metatarsals). The tarsus transmits body weight from the leg to the foot, provides a flexible arch, and permits inversion/eversion movements at the subtalar and midtarsal joints. Tarsal coalition (congenital fusion of two or more tarsal bones—calcaneonavicular and talocalcaneal most common) causes peroneal spastic flatfoot and limited subtalar motion. Tarsal tunnel syndrome: compression of the posterior tibial nerve beneath the flexor retinaculum (medial ankle), causing burning pain, paraesthesia, and numbness on the plantar surface of the foot.

\medterm{Teeth}-- Hard, calcified, mineralized structures embedded within the alveolar processes of the maxilla and mandible, organized in superior and inferior dental arches for mastication, speech articulation, and facial profile support. Humans exhibit diphyodont dentition: primary (deciduous) dentition comprises 20 teeth (2 incisors, 1 canine, 2 molars per quadrant; formula 2-1-0-2), whereas permanent dentition consists of 32 teeth (2 incisors, 1 canine, 2 premolars, 3 molars per quadrant; formula 2-1-2-3). Each tooth comprises a clinical crown (projecting above the gingiva), a neck (cervix), and one or more roots anchored within bony alveolar sockets (gomphosis joint) by the periodontal ligament. Histologically, the crown is covered by enamel (the hardest substance in the human body, 96 percent inorganic hydroxyapatite, non-regenerative after ameloblast loss), overlying dentin (living mineralized tissue with dentinal tubules formed by odontoblasts). The root is covered by cementum. The central pulp cavity contains vascularized, innervated, loose connective tissue (dental pulp) that enters through the apical foramen. Innervation is mediated by the maxillary nerve (CN V2 via superior alveolar nerves) for upper teeth, and the mandibular nerve (CN V3 via the inferior alveolar nerve) for lower teeth. Pathologies include dental caries (bacterial demineralization of enamel and dentin), pulpitis, periapical abscess, periodontal disease (loss of supporting alveolar bone), and third molar (wisdom tooth) impaction.

\medterm{Temporal Bones}-- Paired bones forming the cranial floor and lateral wall, with squamous, petrous, mastoid, tympanic, and styloid parts. Contain the external auditory canal, middle and inner ear, and foramina: carotid canal (internal carotid), jugular foramen (CN IX-XI, IJV), and foramen lacerum. Fractures may cause hearing loss, facial nerve palsy, and CSF otorrhea.

\medterm{Temporomandibular Joint}-- Synovial ginglymoarthrodial joint between mandibular condyle and temporal bone's mandibular fossa, separated by an articular disc. Only bilateral body joint, allowing hinge and sliding movements. Disc divides into superior/inferior compartments. Temporomandibular disorders cause pain, clicking, limited movement, common with bruxism.

\medterm{Tendon}-- Dense regular connective tissue connecting muscle to bone, composed primarily of type I collagen fibers arranged in parallel bundles. Transmits muscle force to bone to produce movement. Poor vascularity (slower healing). Enclosed in tendon sheath (synovial) at friction points. Common injuries: tendinitis (inflammation), tendinosis (degeneration), rupture (Achilles, rotator cuff, biceps). Examples: Achilles tendon (gastrocnemius/soleus to calcaneus), patellar tendon (quadriceps to tibial tuberosity), rotator cuff tendons.

\medterm{Testis (Testicle)}-- The male gonad, a paired ovoid organ (4--5 cm \(\times\) 2--3 cm, weight 15--20 g) located in the scrotum, responsible for sperm production (spermatogenesis) and testosterone secretion. Each testis is covered by the tunica vaginalis (peritoneal remnant) and tunica albuginea (dense fibrous capsule). Internally: 200--300 lobules, each containing 1--4 seminiferous tubules (site of spermatogenesis—60--70 days), lined by Sertoli cells (support, nourishment, and regulation of spermatogenesis; produce inhibin B, anti-Müllerian hormone) and containing germ cells (spermatogonia $\rightarrow$ primary spermatocytes $\rightarrow$ secondary spermatocytes $\rightarrow$ spermatids $\rightarrow$ spermatozoa). Leydig cells (interstitial cells) in the connective tissue between tubules produce testosterone (regulated by LH). The testis develops in the abdomen and descends into the scrotum at 28--32 weeks gestation under the influence of androgens, guided by the gubernaculum. Undescended testis (cryptorchidism) increases the risk of infertility and testicular cancer. The testis is innervated by the genitofemoral nerve (L1--L2) and supplied by the testicular artery.

\medterm{Thalamus}-- Bilateral ovoid gray matter structure forming the third ventricle's lateral walls, the brain's relay station for nearly all sensory (except olfaction), motor, and limbic information to the cortex. Numerous nuclei: VPL (spinothalamic, dorsal column), VPM (trigeminal), LGN (visual), MGN (auditory), VL/VA (motor). Lesions cause contralateral sensory loss, Dejerine-Roussy thalamic pain syndrome, and cognitive impairment.

\medterm{Thoracic Cage}-- The bony and cartilaginous framework formed by the thoracic vertebrae, ribs, costal cartilages, and sternum. It protects thoracic organs and assists respiratory movements.

\medterm{Thoracic Cavity}-- The upper ventral cavity of the trunk, protected and enclosed by the thoracic cage (sternum, 12 pairs of ribs, costal cartilages, and 12 thoracic vertebrae) and intercostal musculature, bounded inferiorly by the diaphragm. It is compartmentalized into three primary spaces: two lateral pleural cavities, each enclosing a lung and lined by visceral and parietal pleurae, and the central mediastinum housing the heart, pericardium, great vessels, trachea, and esophagus. The thoracic cavity volume dynamically changes during respiration through diaphragm excursion and costal movement ('bucket-handle' and 'pump-handle' mechanics). Clinical conditions include pneumothorax, hemothorax, pleural effusion, empyema, flail chest, and cardiac tamponade.

\medterm{Thoracic Duct}-- Body's largest lymphatic vessel, approximately 40 cm long, draining lymph from the entire body below the diaphragm and left side above it (approximately 75 percent of total lymph). Begins at the cisterna chyli (abdominal lymphatic sac) at the L1-L2 level, ascends through the aortic hiatus, through the posterior mediastinum between the aorta and azygos vein, and empties into the junction of the left internal jugular and left subclavian veins. Blocked or damaged thoracic duct causes chylothorax (leakage of chyle into the pleural space), typically on the left side, manifesting as a milky pleural effusion rich in triglycerides.

\medterm{Thoracic Vertebrae (T1-T12)}-- Have costal facets on bodies and transverse processes for rib articulation. Heart-shaped bodies, long downward spinous processes, coronally-oriented facets allowing rotation and lateral flexion but limiting flexion/extension. Small circular vertebral foramen. Most rigid vertebral region due to the rib cage, least affected by disc herniations.

\medterm{Thymus}-- Primary lymphoid organ in the anterior mediastinum, largest in childhood, involutes with age. Site of T-cell maturation (thymopoiesis): thymocytes undergo positive and negative selection to produce self-tolerant, MHC-restricted T cells. Epithelial cells produce thymic hormones (thymosin, thymopoietin). Blood supply: internal thoracic arteries. Conditions: thymoma, myasthenia gravis, DiGeorge syndrome, thymic hyperplasia, thymic carcinoma.

\medterm{Thyroid Gland}-- An endocrine gland located in the anterior neck, consisting of two lobes (right and left) connected by an isthmus, lying anterior to the trachea between the cricoid cartilage and the suprasternal notch. The thyroid produces tri-iodothyronine ($T_3$) and thyroxine ($T_4$)—iodinated tyrosine derivatives that regulate metabolism, growth, and development—and calcitonin (produced by parafollicular/C-cells, regulates calcium metabolism). Thyroid hormone synthesis: iodide is actively transported into follicular cells (NIS/sodium-iodide symporter), oxidised by thyroid peroxidase (TPO), and iodinated to tyrosine residues on thyroglobulin (Tg) to form MIT and DIT, which couple to form $T_3$ and $T_4$. Secretion is regulated by the hypothalamic–pituitary–thyroid axis (TRH $\rightarrow$ TSH $\rightarrow$ thyroid hormone $\rightarrow$ negative feedback). Disorders: hyperthyroidism (Graves disease, toxic nodular goitre, thyroiditis), hypothyroidism (Hashimoto thyroiditis, iodine deficiency, post-surgical/radioiodine), thyroid nodules and thyroid cancer (papillary—most common, follicular, medullary, anaplastic), goitre, and thyroiditis. The thyroid gland develops from the foramen cecum at the base of the tongue (week 4 of gestation) and descends to its final position; thyroglossal duct remnants can form midline neck cysts.

\medterm{Tibia}-- Larger leg bone, the main weight-bearer transmitting body weight from knee to ankle. Proximal: tibial plateau (medial/lateral condyles) and tibial tuberosity (patellar ligament insertion). Shaft: sharp subcutaneous anterior border. Distal: medial malleolus articulating with the talus at the ankle. Tibial plateau fractures are common in high-energy trauma.

\medterm{Tibialis Anterior Muscle}-- Superficial anterior leg muscle from lateral tibial condyle and shaft to medial cuneiform and first metatarsal base. Innervated by deep peroneal nerve (L4-L5). Primary dorsiflexor for foot clearance during gait swing phase; inverts foot. Weakness causes foot drop with steppage gait. Commonly affected by anterior compartment syndrome and shin splints.

\medterm{Tibialis Posterior Muscle}-- Deepest posterior leg muscle from posterior tibia, fibula, and interosseous membrane to navicular tuberosity, medial cuneiform, and sustentaculum tali. Innervated by tibial nerve (L4-L5). Plantarflexes and inverts foot; primary dynamic medial longitudinal arch stabilizer. Dysfunction causes adult-acquired flatfoot with arch flattening, hindfoot valgus, and medial ankle pain.

\medterm{Tibial Nerve}-- Sciatic division (L4-S3) continuing from posterior thigh through the popliteal fossa into the deep posterior leg, innervating gastrocnemius, soleus, tibialis posterior, flexor digitorum longus, and flexor hallucis longus. Passes behind the medial malleolus as the tarsal tunnel, dividing into medial/lateral plantar nerves for the foot. Sensation to the sole. Tarsal tunnel syndrome causes burning pain and paresthesias in the sole.

\medterm{Tissue}-- An organized cellular and extracellular ensemble of similarly specialized cells and their intercellular matrix that cooperate to perform a specific physiological function. All human organs and body structures are composed of varying combinations of four fundamental primary tissue types: epithelial tissue (covering and lining surfaces, gland formation), connective tissue (support, protection, binding, transport), muscle tissue (contraction and movement), and nervous tissue (sensory reception, impulse transmission, integration). The study of tissues at the microscopic level is histology; disruption of normal tissue architecture is the hallmark of histopathology.

\medterm{Tongue}-- A mobile, highly muscular organ situated in the oral cavity and oropharynx, essential for mastication, deglutition, taste sensation, and articulate speech. It is divided into an anterior two-thirds (oral part, in the oral cavity) and a posterior one-third (pharyngeal part, in the oropharynx) by the V-shaped sulcus terminalis, at the apex of which lies the foramen cecum (embryonic origin of the thyroglossal duct). The dorsal mucosa of the anterior tongue is covered with four varieties of lingual papillae: filiform (keratinized, mechanical, lack taste buds), fungiform (mushroom-shaped, vascular, scattered taste buds), circumvallate (large, 8--12 in a V-row anterior to sulcus terminalis, rich in taste buds and serous von Ebner glands), and foliate (lateral folds). Its musculature comprises four intrinsic muscles (superior longitudinal, inferior longitudinal, transverse, vertical) that alter tongue shape, and four extrinsic muscles (genioglossus, hyoglossus, styloglossus, palatoglossus) that alter position. Motor innervation to all intrinsic and extrinsic muscles is mediated by the hypoglossal nerve (CN XII), except palatoglossus which is innervated by the vagus nerve (CN X via the pharyngeal plexus). Sensory innervation is anatomically tripartite: the anterior two-thirds receives general somatic sensation via the lingual nerve (CN V3) and special taste sensation via the chorda tympani (CN VII traveling with the lingual nerve); the posterior one-third receives both general sensation and taste via the glossopharyngeal nerve (CN IX); and the base near the epiglottis receives sensation via the internal laryngeal nerve (CN X). Lymphatic drainage from the tip passes to submental nodes, lateral anterior tongue to submandibular nodes, and posterior tongue directly to deep cervical (jugulodigastric) nodes. Hypoglossal nerve injury causes the protruded tongue to deviate toward the side of the lesion.

\medterm{Tonsils}-- Aggregated lymphoid tissue forming Waldeyer's ring at the pharyngeal entrance. Pharyngeal tonsils (adenoids) in the nasopharynx, palatine tonsils at the oropharyngeal pillars, lingual tonsils at the tongue base, and tubal tonsils near the Eustachian tube orifices. They are the first line of defense against inhaled or ingested pathogens. Palatine tonsils are most commonly infected (tonsillitis) and surgically removed (tonsillectomy) for recurrent infections or obstructive sleep apnoea. The tonsils are part of Waldeyer's ring, a protective lymphoid barrier at the pharyngeal inlet that samples inhaled and ingested antigens, with maximal activity in early childhood.

\medterm{Trachea}-- The windpipe, a 10-12 cm fibrocartilaginous tube from the larynx (C6) to the carina (T4/T5) where it bifurcates into main bronchi. Sixteen to twenty C-shaped hyaline cartilage rings (open posteriorly, completed by trachealis muscle) maintain patency. Lined by ciliated pseudostratified columnar epithelium with goblet cells (mucociliary escalator). Relations: thyroid isthmus (anterior), esophagus (posterior), recurrent laryngeal nerves (tracheoesophageal groove). Blood supply: inferior thyroid arteries. Conditions: tracheitis, tracheal stenosis, tracheomalacia, tracheoesophageal fistula, tracheostomy.

\medterm{Transverse Colon}-- Longest, most mobile colonic segment from hepatic to splenic flexure, approximately 45 cm, with complete mesentery (transverse mesocolon). Crosses anterior to duodenum, pancreas, and major vessels. Splenic flexure is the highest colon point and common gas trapping site. Most common diverticular disease site along with the sigmoid.

\medterm{Transverse Plane}-- A horizontal anatomical plane (also termed the axial plane) oriented perpendicular to both the sagittal and coronal planes, dividing the body or an organ into superior (upper) and inferior (lower) portions. Movements taking place in or parallel to the transverse plane include rotational movements (pronation, supination, medial and lateral rotation). In cross-sectional diagnostic imaging (axial CT and MRI scans), images are viewed standardly from inferior to superior (as if looking up from the patient's feet).

\medterm{Transversus Abdominis Muscle}-- Deepest flat abdominal muscle with horizontal fibers, from inguinal ligament, iliac crest, thoracolumbar fascia, and inner lower six costal cartilages. Innervated by T7-T12 and L1. Compresses abdomen without trunk movement, providing core stability and increased intra-abdominal pressure during coughing, lifting, Valsalva. Key for core stability exercises and low back pain rehabilitation.

\medterm{Trapezium}-- Saddle-shaped distal row carpal bone on the radial side, articulating with the first metacarpal at the CMC joint of the thumb. This saddle joint allows thumb opposition essential for grip and fine motor function. Basal joint arthritis is common in postmenopausal women, causing pinch grip pain and weakness.

\medterm{Trapezius Muscle}-- Large triangular muscle from occipital bone and C7-T12 spinous processes to clavicle, acromion, and scapular spine. Innervated by CN XI and C3-C4. Upper fibers elevate scapula, middle fibers retract, lower fibers depress. Essential for shoulder stability and scapulothoracic motion. CN XI injury causes scapular winging and difficulty with abduction above 90 degrees.

\medterm{Trapezoid}-- Smallest distal row carpal bone between trapezium and capitate, articulating with the second metacarpal. Wedge shape provides stability for the firmly fixed second metacarpal. Fractures rare (less than 1 percent of carpal fractures). Contributes to the radial column wrist stability.

\medterm{Triceps Brachii Muscle}-- Three-headed posterior arm muscle. Long head from infraglenoid tubercle, lateral head from posterior humerus above radial groove, medial head below. All insert on the olecranon. Innervated by radial nerve (C7-C8). Primary elbow extensor; assists shoulder extension/adduction. Triceps reflex (C7) tests radial nerve/spinal cord integrity.

\medterm{Tricuspid Valve}-- Right atrioventricular valve with anterior, posterior, and septal cusps supported by chordae tendineae and three papillary muscles. Largest cardiac valve. Most commonly affected in IV drug use-associated endocarditis. Tricuspid regurgitation from right ventricular dilation, endocarditis, or carcinoid. Stenosis is rare, mostly rheumatic. Thinner and more delicate than the mitral valve.

\medterm{Trigeminal Nerve}-- Cranial nerve V (CN V), the largest cranial nerve, providing sensory innervation to the face and anterior scalp (ophthalmic $V_1$, maxillary $V_2$, mandibular $V_3$ divisions) and motor innervation to the muscles of mastication (masseter, temporalis, medial and lateral pterygoids, tensor veli palatini, tensor tympani, mylohyoid, anterior belly of digastric). The trigeminal nerve emerges from the lateral pons as a large sensory root (trigeminal ganglion/ semilunar/Gasserian ganglion in Meckel cave) and a smaller motor root. Sensory distribution: $V_1$ (ophthalmic)—forehead, scalp, upper eyelid, cornea, nasal tip; $V_2$ (maxillary)—cheek, lower eyelid, side of nose, upper lip, upper teeth and gums, palate; $V_3$ (mandibular)—lower lip, chin, lower teeth and gums, anterior two-thirds of tongue (general sensation), floor of mouth, temporal region. Motor: $V_3$ supplies the muscles of mastication, tensor veli palatini, tensor tympani, mylohyoid, and anterior belly of digastric. Reflexes: corneal reflex (afferent $V_1$, efferent CN VII), jaw jerk (afferent and efferent $V_3$—monosynaptic). Lesions: trigeminal neuralgia ($V_2$/$V_3$), trigeminal schwannoma, cavernous sinus syndrome, and brainstem stroke.

\medterm{Triquetrum}-- Pyramidal proximal row carpal bone on the ulnar side, articulating with lunate, hamate, pisiform, and TFCC. Second most commonly fractured carpal bone. Contributes to proximal carpal row and ulnar wrist stability. Pisiform articulates with its volar surface.

\medterm{Trochlear Nerve (Cranial Nerve IV)}-- The fourth cranial nerve (CN IV), a purely somatic motor nerve that innervates the superior oblique muscle of the eye. It holds three unique anatomical distinctions among all cranial nerves: it is the smallest cranial nerve by axon count, the only cranial nerve to emerge from the dorsal (posterior) surface of the brainstem (below the inferior colliculi in the midbrain), and the only cranial nerve whose fibers undergo complete decussation within the brainstem prior to emergence, meaning the nucleus controls the contralateral superior oblique. It curves anteriorly around the midbrain, traverses the lateral wall of the cavernous sinus below CN III, and enters the orbit through the superior orbital fissure outside the common tendinous ring. The superior oblique muscle intorts, depresses, and abducts the eye; its depressive action is maximal when the eye is adducted. Trochlear nerve palsy is the most common cause of acquired vertical diplopia; patients present with vertical/torsional diplopia (worse on downward gaze, such as reading or walking downstairs) and characteristically tilt their head toward the opposite side to compensate (Bielschowsky head tilt test).

\medterm{True Ribs}-- Ribs 1-7 articulating directly with the sternum via their own costal cartilages. Each has a head, neck, tubercle, and shaft. The first two are short and strongly curved for neck/shoulder muscle attachment. Costal cartilages provide thoracic cage elasticity for respiratory expansion. Less commonly fractured than false ribs due to shorter length and more protected position.

\medterm{Tunica Adventitia}-- The outer connective-tissue layer of a blood-vessel wall. It contains collagen and elastic fibers that provide structural support, along with nerves and, in larger vessels, small vessels called vasa vasorum that supply the outer wall. It is relatively prominent in veins and helps anchor vessels to surrounding tissues.

\medterm{Tunica Media}-- The middle layer of a blood-vessel wall, located between the tunica intima and tunica adventitia. It consists primarily of smooth muscle and elastic fibers; its contraction or relaxation changes vascular diameter, regulating blood pressure and tissue perfusion. It is thickest in arteries and thinner in veins.

\medterm{Tympanic Membrane}-- Thin translucent cone-shaped membrane separating external and middle ear, approximately 10 mm diameter. Three layers: outer cutaneous, middle fibrous, inner mucous. Malleus handle is visible through it. Light reflex (cone of light) in the anteroinferior quadrant indicates normal position. Perforation from trauma, infection (otitis media), or barotrauma. Myringotomy and tympanostomy tube placement treat chronic otitis media with effusion. Vibrates in response to sound, transmitting energy to the ossicular chain and ultimately to the cochlear fluid for auditory transduction.

\medterm{Ulna}-- Medial forearm bone, longer than the radius. Proximal end: trochlear notch, olecranon, coronoid process. Shaft. Distal end: head, styloid process. The trochlear notch articulates with the trochlea forming the primary elbow hinge. Bears approximately 20 percent of wrist axial load. The olecranon is the bony elbow prominence and triceps insertion. Serves as the forearm rotation axis.

\medterm{Ulnar Artery}-- The larger terminal branch of the brachial artery at the cubital fossa, running medially down the forearm to the wrist (palpable medial to the pisiform) and hand. Supplies the medial forearm and ulnar side of the hand; forms the superficial palmar arch. Provides the common interosseous artery. Clinical: ulnar artery pulse, Allen test for collateral circulation, and ulnar artery cannulation. The ulnar nerve travels with the artery at the wrist (Guyon canal); ulnar artery occlusion causes hypothenar pain and digital ischemia.

\medterm{Ulnar Nerve}-- From medial cord (C8-T1), passing behind the medial epicondyle (cubital tunnel), into the forearm innervating flexor carpi ulnaris and medial FDP half. In the hand innervates hypothenar muscles, interossei, medial two lumbricals, and adductor pollicis. Sensation to medial 1.5 digits. Injury (elbow trauma, Guyon's canal compression) causes claw hand (MCP hyperextension with IP flexion of digits 4-5) and loss of grip/pinch strength.

\medterm{Umbilical Region}-- Central region around the umbilicus containing small intestine, transverse colon, and great vessels. The umbilicus is a landmark for assessing ascites, caput medusae in portal hypertension, and umbilical hernias. Percussion and auscultation here are essential components of the abdominal examination.

\medterm{Umbilicus (Navel)}-- The scar on the anterior abdominal wall marking the site of attachment of the umbilical cord in fetal life, located at the level of the L3-L4 intervertebral disc, approximately midway between the xiphoid process and the pubic symphysis. The umbilicus is a remnant of fetal structures: the umbilical vein (ligamentum teres hepatis, running to the liver), the umbilical arteries (medial umbilical ligaments), and the urachus (median umbilical ligament, running to the bladder apex). It is a common site of hernia (umbilical hernia) and is used as a laparoscopic port entry site (Hasson technique, Veress needle insertion for pneumoperitoneum). Dermatomal supply is T10. The aorta is palpable at the level of the umbilicus, and the iliac bifurcation lies at the L4 level, approximately 1 cm below it.

\medterm{Ureter}-- The paired muscular tubes (25-30 cm) conveying urine from the renal pelvis to the urinary bladder, crossing the pelvic brim at the pelvic inlet. Three anatomical narrowings (pelviureteric junction, pelvic brim, vesicoureteric junction) are common sites of ureteric stone impaction. Peristalsis propels urine; valves (oblique intramural passage) prevent reflux. Blood supply: segmental (renal, gonadal, iliac, vesical arteries). Ureteric obstruction (stones, tumor, compression) causes hydronephrosis and renal colic. Ureteroscopy/ureteric stents are used in management.

\medterm{Urethra}-- The tube conveying urine from the bladder to the exterior. Male urethra (18-20 cm): prostatic, membranous, and spongy (penile) parts; carries urine and semen; longer and more tortuous than female. Female urethra (3-4 cm): short, straight, predisposing to ascending urinary tract infections. The external urethral sphincter (striated, voluntary) provides continence; internal sphincter (smooth). Male urethral trauma occurs with pelvic fracture (membranous) and straddle injuries (bulbar). Urethral strictures cause obstructive voiding symptoms.

\medterm{Urinary Bladder}-- A hollow, distensible muscular organ located in the pelvic cavity posterior to the pubic symphysis, serving as a reservoir for urine. It consists of the detrusor muscle (smooth muscle layer), the trigone (a smooth triangular area at the base bounded by the two ureteral openings and the internal urethral orifice), and transitional epithelium (urothelium) that accommodates stretching. The bladder is controlled by both autonomic and somatic nerves to regulate micturition.

\medterm{Uterus}-- Hollow muscular organ in the pelvis between the bladder and rectum, approximately 7.5 cm long in adult females, responsible for menstruation, implantation, and fetal development during pregnancy. Three layers: endometrium (inner mucosa, cyclically shed during menstruation), myometrium (middle smooth muscle layer, contracts during labor), and perimetrium (outer serous coat). The uterus has a fundus (superior portion above the tube openings), body, and cervix (lower portion protruding into the vagina). The cervix has an internal and external os, and its endocervical canal is a common site for cervical cancer screening by smear testing.

\medterm{Vagina}-- Muscular fibromembranous canal extending from the cervix to the vulva, approximately 7-9 cm long in adult females. Functions include receiving the penis during intercourse, serving as the birth canal, and channeling menstrual flow. The vaginal walls have rugae (transverse folds allowing distension) and are lined with stratified squamous epithelium (non-keratinized, responding to estrogen by producing glycogen). Lactobacilli in the vaginal flora metabolize glycogen to lactic acid, maintaining an acidic environment (pH 4-5) that protects against pathogenic colonisation by opportunistic organisms.

\medterm{Vagus Nerve (Cranial Nerve X)}-- The tenth cranial nerve (CN X), the longest and most widely distributed cranial nerve in the body. It contains sensory, motor, and parasympathetic fibers, exiting the skull via the jugular foramen. It provides parasympathetic innervation to the thoracic and abdominal viscera (heart, lungs, gastrointestinal tract up to the splenic flexure), regulates heart rate, gastrointestinal motility, and speech (via the recurrent laryngeal nerve), and is a key mediator of the vasovagal reflex.

\medterm{Vas Deferens}-- Muscular duct (approximately 45 cm long) transporting sperm from the epididymal tail to the ejaculatory duct. It ascends through the spermatic cord, passes through the inguinal canal, enters the pelvis, crosses over the ureter, and joins the seminal vesicle duct to form the ejaculatory duct (in the prostate). The vas deferens has thick smooth muscle walls providing peristaltic propulsion of sperm. It is the structure cut during vasectomy (male sterilization), typically at the scrotal level where it is easily palpable through the scrotal skin, allowing reliable identification and division with minimal morbidity.

\medterm{Vein}-- A blood vessel that carries deoxygenated blood toward the heart (except pulmonary veins, which carry oxygenated blood). Veins have thinner walls and larger lumens than arteries, with less smooth muscle and elastic tissue. Many veins contain one-way valves (especially in the legs) that prevent backflow and assist venous return against gravity. Venous return is facilitated by the skeletal muscle pump, respiratory pump, and cardiac suction. Venous system: superficial veins (subcutaneous), deep veins (accompany arteries), and perforating veins (connect superficial and deep systems). Conditions: varicose veins, deep vein thrombosis, chronic venous insufficiency.

\medterm{Vena Cava}-- The two largest systemic veins returning deoxygenated blood to the right atrium. Superior Vena Cava (SVC): formed by brachiocephalic veins, drains head, neck, upper limbs, thorax (no valves). Inferior Vena Cava (IVC): formed by common iliac veins at L5, drains abdomen, pelvis, lower limbs (receives hepatic, renal, gonadal, lumbar veins). SVC syndrome: obstruction causes facial/upper body swelling, dyspnea, distended neck veins. IVC filter for PE prevention.

\medterm{Ventral}-- A directional term relating to or situated toward the belly (anterior) surface of the body. In human bipedal anatomy, ventral is largely synonymous with anterior. In neuroanatomy and spinal cord anatomy, the ventral roots carry motor (efferent) fibers, and the ventral horns contain lower motor neuron somas. Opposed to dorsal.

\medterm{Ventricle}-- One of the two lower, thicker-walled chambers of the heart that pump blood out of the heart. The right ventricle pumps deoxygenated blood into the pulmonary artery (low-pressure system); the left ventricle pumps oxygenated blood into the aorta (high-pressure system, generating systemic blood pressure). The left ventricular wall is 3--4 times thicker than the right. The interventricular septum separates the ventricles. Ventricular contraction (systole) ejects blood; ventricular relaxation (diastole) allows filling. Conditions: ventricular hypertrophy, ventricular failure, ventricular septal defect, ventricular tachycardia/fibrillation. Also refers to the four fluid-filled cavities within the brain (two lateral ventricles, third, and fourth ventricles) that produce and circulate cerebrospinal fluid (CSF).

\medterm{Ventricle (Brain)}-- The interconnected system of fluid-filled cavities within the brain producing and containing cerebrospinal fluid: two lateral ventricles (cerebrum), third ventricle (diencephalon), and fourth ventricle (pons/medulla, continuous with the central canal). CSF is produced by the choroid plexus, flows through the interventricular foramina (Monro), cerebral aqueduct (Sylvius), and median/lateral apertures (Magendie/Luschka) to the subarachnoid space. Obstruction (e.g., aqueductal stenosis, tumor) causes hydrocephalus with raised intracranial pressure. Ependymal lining.

\medterm{Venule}-- A small vessel that drains blood from capillary beds into veins. Postcapillary venules are important sites for leukocyte migration and inflammatory fluid exchange.

\medterm{Vertebra}-- The individual bone of the vertebral column (33 vertebrae: 7 cervical, 12 thoracic, 5 lumbar, 5 sacral fused, 4 coccygeal). Parts: vertebral body, vertebral arch (pedicles, laminae), spinous and transverse processes, articular processes/facets, and vertebral foramen (spinal canal). Regional variations: cervical (transverse foramina, bifid spine), thoracic (costal facets for ribs, long downward spinous processes), lumbar (large bodies). Vertebral bodies articulate via intervertebral discs. Fractures, spondylolisthesis, and disc herniation are common clinical conditions.

\medterm{Vertebral Canal}-- The continuous, protective longitudinal bony canal formed by the sequential alignment of the vertebral foramina of all 33 articulated vertebrae, extending from the foramen magnum of the occipital bone superiorly to the sacral hiatus inferiorly. It encloses and protects the spinal cord, spinal nerve roots, cauda equina, meninges (dura, arachnoid, and pia mater), the internal vertebral venous plexus (Batson's plexus), and epidural fat. The spinal cord terminates within the canal at the L1--L2 intervertebral disc level as the conus medullaris, below which the canal houses the lumbosacral roots of the cauda equina bathed in CSF within the lumbar cistern (the target site for lumbar puncture at L3--L4 or L4--L5). Spinal stenosis is narrowing of the canal that compresses neural elements, causing neurogenic claudication.

\medterm{Vertebral Column}-- 33 vertebrae: 7 cervical, 12 thoracic, 5 lumbar, 5 fused sacral, 4 fused coccygeal. Protects the spinal cord, supports the head/trunk, provides muscle/rib attachment. Each typical vertebra has a body, vertebral arch, and 7 processes. Intervertebral discs (nucleus pulposus/annulus fibrosus) provide cushioning. Four curvatures: cervical lordosis, thoracic kyphosis, lumbar lordosis, sacral kyphosis.

\medterm{Vestibule}-- Central bony labyrinth between cochlea and semicircular canals, containing utricle (horizontal linear acceleration and head tilt detection) and saccule (vertical linear acceleration detection). Each has a macula with hair cells covered by gelatinous membrane embedded with otoconia (calcium carbonate crystals). Forces deflect otoconia, bending stereocilia and generating nerve impulses. The vestibular system works with visual and proprioceptive systems for balance. Vestibular neuritis and labyrinthitis cause acute vertigo, nystagmus, and balance disturbance, while Menière's disease involves episodic vertigo with tinnitus and hearing loss.

\medterm{Vestibulocochlear Nerve (Cranial Nerve VIII)}-- The eighth cranial nerve (CN VIII), a purely special sensory nerve responsible for the senses of hearing and equilibrium (balance). It consists of two anatomically distinct divisions: the cochlear nerve and the vestibular nerve. The cochlear nerve originates from bipolar cell bodies in the spiral ganglion of the cochlea, which receive inputs from auditory hair cells in the organ of Corti. The vestibular nerve arises from bipolar cell bodies in the vestibular (Scarpa's) ganglion, receiving inputs from hair cells in the semicircular duct ampullae (detecting angular acceleration) and the utricle and saccule maculae (detecting linear acceleration and gravity). Both divisions unite and emerge from the internal acoustic meatus alongside the facial nerve and labyrinthine artery, traversing the subarachnoid space of the cerebellopontine angle to enter the brainstem at the pontomedullary junction. Pathologies include vestibular schwannoma (acoustic neuroma, a benign tumor of Schwann cells typically arising in the internal acoustic meatus, causing progressive unilateral sensorineural hearing loss, tinnitus, and balance impairment), vestibular neuritis (acute prolonged vertigo without hearing loss), and labyrinthitis.

\medterm{Viscera}-- The internal organs of the body, particularly those within the thoracic and abdominal cavities. The thoracic viscera include the heart, lungs, esophagus, trachea, and thymus. The abdominal viscera are divided into the digestive organs (stomach, small intestine, large intestine, liver, gallbladder, pancreas), the urinary organs (kidneys, ureters, bladder), and the reproductive organs (uterus, ovaries in females; prostate, seminal vesicles in males). Visceral pain is poorly localised, vague, and referred to dermatomes that share the same spinal cord level (e.g., cardiac ischemic pain referred to the left arm and jaw, gallbladder pain referred to the right shoulder). Visceral peritoneum covers the abdominal organs; the mesentery suspends the intestines and carries blood vessels, lymphatics, and nerves. The viscera receive dual autonomic innervation: sympathetic (fight or flight) and parasympathetic (rest and digest, primarily via the vagus nerve).

\medterm{Vocal Folds}-- Paired mucosal folds within the larynx that vibrate as air passes between them to produce sound. They also close to protect the airway and assist in raising intrathoracic pressure.

\medterm{Vulva}-- External female genitalia, including the mons pubis (fatty mound over the pubic symphysis), labia majora (outer fatty folds covered with hair), labia minora (inner hairless folds), clitoris (erectile tissue homologous to the penis, with the glans being the most sensitive erogenous zone), vestibular glands (Bartholin's glands lubricating the vaginal opening), and the urethral and vaginal openings. The vulva is innervated by the pudendal nerve. Vulvar skin conditions include lichen sclerosus, vulvar intraepithelial neoplasia (VIN), and vulvovaginitis, each presenting with pruritus, irritation, and visible changes to the vulvar skin and mucosa.

\medterm{White Matter}-- Regions of the central nervous system (CNS) composed primarily of bundles of myelinated axons, oligodendrocytes (the myelin-producing glial cells), and fibrous astrocytes, exhibiting a characteristic pale white appearance due to the high lipid content of myelin sheaths. It underlies the cerebral and cerebellar cortices and forms the outer ascending and descending funiculi (columns) of the spinal cord, organized into tracts (e.g., corticospinal, spinothalamic, dorsal columns). In the cerebrum, it is classified into commissural fibers (connecting hemispheres, e.g., corpus callosum), association fibers (connecting ipsilateral cortical areas), and projection fibers (connecting cortex with subcortical and spinal structures, e.g., internal capsule). Pathologies include demyelinating diseases (multiple sclerosis, progressive multifocal leukoencephalopathy) and diffuse axonal injury in severe traumatic brain injury.

\medterm{Wrist Joint (Radiocarpal Joint)}-- A biaxial condyloid (ellipsoid) synovial joint connecting the distal upper extremity to the carpus. It is formed proximally by the distal articular surface of the radius and the fibrocartilaginous triangular fibrocartilage complex (TFCC, articular disc), and distally by the proximal row of carpal bones (scaphoid, lunate, and triquetrum; the pisiform is a sesamoid in the flexor carpi ulnaris tendon and does not articulate). The ulna does not articulate directly with the carpal bones. The joint permits flexion, extension, radial deviation (abduction), ulnar deviation (adduction, greater range than abduction), and circumduction; pronation and supination occur at the proximal and distal radioulnar joints, not the radiocarpal joint. Stabilized by dorsal and palmar radiocarpal ligaments, ulnocarpal ligaments, and radial and ulnar collateral ligaments. Clinical conditions include distal radius fractures (Colles' and Smith's fractures), scaphoid fractures (prone to avascular necrosis), TFCC tears causing ulnar-sided wrist pain, carpal instability, and rheumatoid arthritis.

\medterm{Zygoma (Zygomatic Bone)}-- The cheekbone, a paired, quadrangular bone of the facial skeleton (viscerocranium) that forms the prominence of the cheek and the lateral wall and floor of the orbit. The zygomatic bone articulates with four bones: the maxilla (at the zygomaticomaxillary suture, forming the infraorbital rim), the temporal bone (at the zygomaticotemporal suture, forming the zygomatic arch, which is completed posteriorly by the temporal process of the zygomatic bone and anteriorly by the zygomatic process of the temporal bone), the frontal bone (at the zygomaticofrontal suture, forming the lateral orbital margin), and the sphenoid bone (at the zygomaticosphenoid suture). The zygomatic bone has three surfaces: lateral (facial, for attachment of the zygomaticus major and minor muscles), orbital (forming the lateral and inferior orbital walls, containing the zygomaticofacial and zygomaticotemporal foramina for the corresponding nerves), and temporal (forming part of the temporal fossa). The zygomatic nerve (branch of V2, maxillary division of the trigeminal nerve) runs through the zygomatic bone, providing sensory innervation. Zygomatic fractures (tripod/trimalar fracture, zygomaticomaxillary complex fracture) are the second most common facial fracture (after nasal bone fractures) and cause flattening of the malar prominence, periorbital ecchymosis, subconjunctival hemorrhage, trismus (from impingement of the coronoid process of the mandible), and infraorbital nerve paraesthesia (numbness of the cheek, upper lip, and lateral nose).

